% Frame Beamer riordinati dal sorgente fornito dall'utente.
% Frammento da includere nel documento Beamer con \input{...}.
% Richiede amsmath, amssymb (normalmente già disponibili in Beamer).
% Convenzioni: metrica (+---), hbar=c=1 salvo diversa indicazione.
% Nelle sezioni sulle particelle libere u_s(p) è uno spinore numerico;
% nei doppietti di QCD, u(x) indica invece il campo del quark up.

\section{Dalla teoria di Dirac ai campi quantistici}


\subsection{Origine dell'equazione di Dirac}

%----------------------------------------------------------
\begin{frame}{Dalla relazione relativistica all'equazione di Dirac}
\small

\begin{columns}[T,onlytextwidth]
\begin{column}{0.49\textwidth}

Per una particella libera:
\[
E^2=\mathbf p^2+m^2.
\]

Con la corrispondenza quantistica
\[
E\to i\partial_t,\qquad
\mathbf p\to-i\boldsymbol\nabla,
\]
si ottiene l'equazione di Klein--Gordon:
\[
(\partial_t^2-\nabla^2+m^2)\phi=0.
\]

È del secondo ordine nel tempo e descrive
naturalmente un campo relativistico scalare.

\end{column}
\begin{column}{0.49\textwidth}

Dirac cerca un'equazione del primo ordine
nel tempo e nelle derivate spaziali:
\[
i\partial_t\psi=H_D\psi,
\]
\[
H_D=\boldsymbol\alpha\cdot\mathbf p+\beta m.
\]

Qui
$\boldsymbol\alpha=(\alpha_1,\alpha_2,\alpha_3)$
e $\beta$ sono ancora da determinare.

Affinché $H_D^2=\mathbf p^2+m^2$, non
possono essere semplici numeri:
devono essere matrici.

\end{column}
\end{columns}

\end{frame}

%----------------------------------------------------------
\begin{frame}{Perché occorrono quattro componenti?}
\small

L'hamiltoniana di Dirac deve soddisfare,
per ogni $\mathbf p$,
\[
\left(
\boldsymbol\alpha\cdot\mathbf p+\beta m
\right)^2
=(\mathbf p^2+m^2)\mathbf1.
\]

\begin{columns}[T,onlytextwidth]
\begin{column}{0.49\textwidth}

Sviluppando il quadrato, si ottengono
le condizioni algebriche:
\[
\begin{aligned}
\{\alpha_i,\alpha_j\}
&=2\delta_{ij}\mathbf1,\\
\{\alpha_i,\beta\}&=0,\\
\beta^2&=\mathbf1.
\end{aligned}
\]

Qui
\[
\{A,B\}=AB+BA
\]
è l'anticommutatore e $\delta_{ij}$
è il delta di Kronecker;
$i,j=1,2,3$.

\end{column}
\begin{column}{0.49\textwidth}

Servono quattro matrici
$\alpha_1,\alpha_2,\alpha_3,\beta$
mutuamente anticommutanti.

Una rappresentazione complessa
$2\times2$ non è sufficiente:
se $\alpha_i=\sigma_i$
(matrici di Pauli), nessuna matrice
$2\times2$ non nulla anticommuta
con tutte e tre.

In $3+1$ dimensioni, per $m\ne0$,
la rappresentazione minima è $4\times4$:
\[
\psi(x)\in\mathbb C^4.
\]

\end{column}
\end{columns}

\end{frame}

%----------------------------------------------------------
\begin{frame}{Equazione di Dirac e significato dello spinore}
\small

Definiamo
\[
\gamma^0=\beta,\qquad
\gamma^i=\beta\alpha_i
\]
e adottiamo
$g^{\mu\nu}=\operatorname{diag}(1,-1,-1,-1)$.

Allora
\[
\{\gamma^\mu,\gamma^\nu\}
=2g^{\mu\nu}\mathbf1,
\qquad
(i\gamma^\mu\partial_\mu-m)\psi=0.
\]

\begin{columns}[T,onlytextwidth]
\begin{column}{0.49\textwidth}

Nella rappresentazione di Dirac:
\[
\gamma^0=
\begin{pmatrix}
\mathbf1_2&0\\
0&-\mathbf1_2
\end{pmatrix},
\]
\[
\gamma^i=
\begin{pmatrix}
0&\sigma_i\\
-\sigma_i&0
\end{pmatrix}.
\]

Le $\sigma_i$ sono le matrici di Pauli.

\end{column}
\begin{column}{0.49\textwidth}

La funzione incognita è uno
spinore di Dirac:
\[
\psi(x)=
\begin{pmatrix}
\psi_1(x)\\
\psi_2(x)\\
\psi_3(x)\\
\psi_4(x)
\end{pmatrix}.
\]

Le quattro componenti sono ampiezze
complesse accoppiate dalle matrici
$\gamma^\mu$.

Sotto Lorentz, $\psi$ si trasforma
mediante una rappresentazione spinoriale
$S(\Lambda)$, non come un quadrivettore.

\end{column}
\end{columns}

\end{frame}

%----------------------------------------------------------
\begin{frame}{Spin $1/2$, soluzioni libere e indice $s$}
\small

\begin{columns}[T,onlytextwidth]
\begin{column}{0.49\textwidth}

L'operatore di spin è
\[
\mathbf S=\frac12\boldsymbol\Sigma,
\qquad
\Sigma_i=
\begin{pmatrix}
\sigma_i&0\\
0&\sigma_i
\end{pmatrix}.
\]

Per $S_z$, gli autovalori sono
$s=+\frac12,-\frac12$.

Usiamo gli autospinori di Pauli:
\[
\xi_\uparrow=
\begin{pmatrix}1\\0\end{pmatrix},
\qquad
\xi_\downarrow=
\begin{pmatrix}0\\1\end{pmatrix}.
\]

L'indice $s$ etichetta la scelta
dello stato di spin.

\end{column}
\begin{column}{0.49\textwidth}

A riposo $\mathbf p=0$, quindi
\[
H_D=m\beta=
\begin{pmatrix}
m\mathbf1_2&0\\
0&-m\mathbf1_2
\end{pmatrix}.
\]

Esistono due modi con energia $+m$
e due con energia $-m$:
\[
u_s(0)\propto
\begin{pmatrix}\xi_s\\0\end{pmatrix},
\qquad
v_s(0)\propto
\begin{pmatrix}0\\\xi_s\end{pmatrix}.
\]

Il settore a frequenza negativa
è reinterpretato in teoria dei campi
in termini di antiparticelle.

\end{column}
\end{columns}

\medskip

A momento fissato, nel settore $E>0$
restano due spinori indipendenti
$u_s(\mathbf p)$: uno per ciascun valore di $s$.

\end{frame}

\subsection{Spinori relativistici e bilineari}

\begin{frame}{Spinore di Dirac e indice di spin $s$}
\small
\begin{columns}[T,onlytextwidth]
\begin{column}{0.48\textwidth}
Uno spinore di Dirac è un vettore complesso
$u\in\mathbb C^4$. Non è un quadrivettore:
il campo spinoriale trasforma come
\[
 \psi'(x')=S(\Lambda)\psi(x).
\]
Una soluzione libera a energia positiva è
\[
 \psi(x)=u_s(\mathbf p)e^{-ip\cdot x},
\]
con $p\cdot x=Et-\mathbf p\cdot\mathbf x$.
\end{column}
\begin{column}{0.48\textwidth}
Per una particella di massa $m$:
\[
 p^\mu=(E,\mathbf p),\qquad
 E=\sqrt{\mathbf p^2+m^2}.
\]
Fissati $\mathbf p$ ed $E>0$, l'equazione
\[
 (\gamma^\mu p_\mu-m)u_s(\mathbf p)=0
\]
ha due soluzioni indipendenti.
L'indice $s=\pm\frac12$ le distingue:
è la proiezione dello spin lungo un asse
scelto nel sistema di riposo.
\end{column}
\end{columns}
\medskip
Le quattro componenti di $u_s$ sono ampiezze spinoriali;
$s$ non è un indice che numera queste componenti.
\end{frame}

\begin{frame}{Come si determinano le quattro componenti?}
\small
\begin{columns}[T,onlytextwidth]
\begin{column}{0.48\textwidth}
Nella rappresentazione di Dirac:
\[
 \gamma^0=\begin{pmatrix}I_2&0\\0&-I_2\end{pmatrix},
 \quad
 \gamma^i=\begin{pmatrix}0&\sigma^i\\-\sigma^i&0\end{pmatrix}.
\]
Scriviamo $u_s=(\varphi_s,\chi_s)^T$, con
$\varphi_s,\chi_s\in\mathbb C^2$.
L'equazione di Dirac fornisce
\[
 \begin{aligned}
 (E-m)\varphi_s-\boldsymbol\sigma\!\cdot\!\mathbf p\,\chi_s&=0,\\
 \boldsymbol\sigma\!\cdot\!\mathbf p\,\varphi_s-(E+m)\chi_s&=0.
 \end{aligned}
\]
\end{column}
\begin{column}{0.48\textwidth}
$\sigma^i$ sono le matrici di Pauli. Dalla seconda equazione:
\[
 \chi_s=\frac{\boldsymbol\sigma\cdot\mathbf p}{E+m}\varphi_s.
\]
Scegliamo un autospinore di spin
$\xi_s\in\mathbb C^2$ con
$\xi_s^\dagger\xi_s=1$ e poniamo
\[
 \varphi_s=\sqrt{E+m}\,\xi_s.
\]
Così $u_s^\dagger u_s=2E$.
\end{column}
\end{columns}
\medskip
La soluzione esplicita è
\[
 u_s(\mathbf p)=
 \begin{pmatrix}
  \sqrt{E+m}\,\xi_s\\[2pt]
  \dfrac{\boldsymbol\sigma\cdot\mathbf p}{\sqrt{E+m}}\,\xi_s
 \end{pmatrix}.
\]
\end{frame}

\begin{frame}{Esempio: spin su/giù e momento lungo $z$}
\small
Scegliamo $\mathbf p=(0,0,p)$, $p>0$, e la quantizzazione dello spin lungo $z$.
\begin{columns}[T,onlytextwidth]
\begin{column}{0.48\textwidth}
Per $s=+\frac12$:
\[
 \xi_+=\begin{pmatrix}1\\0\end{pmatrix},\qquad
 \frac{\sigma_z}{2}\xi_+=\frac12\xi_+.
\]
\[
 u_+(p)=
 \begin{pmatrix}
 \sqrt{E+m}\\0\\[1pt]
 p/\sqrt{E+m}\\0
 \end{pmatrix}.
\]
\end{column}
\begin{column}{0.48\textwidth}
Per $s=-\frac12$:
\[
 \xi_-=\begin{pmatrix}0\\1\end{pmatrix},\qquad
 \frac{\sigma_z}{2}\xi_-=-\frac12\xi_-.
\]
\[
 u_-(p)=
 \begin{pmatrix}
 0\\\sqrt{E+m}\\[1pt]
 0\\-p/\sqrt{E+m}
 \end{pmatrix}.
\]
\end{column}
\end{columns}
\medskip
A riposo $u_\pm(0)=\sqrt{2m}\,(\xi_\pm,0)^T$.
Le coppie superiore/inferiore in questa rappresentazione
non coincidono con le componenti chirali $L/R$.
\end{frame}

%\begin{frame}{Equazione di Dirac: campi e soluzioni libere}
%\small
%\begin{columns}[T,onlytextwidth]
%\begin{column}{0.48\textwidth}
%Un campo spinoriale classico $\psi(x)$ obbedisce all'equazione libera
%\[
% (i\gamma^\mu\partial_\mu-m)\psi(x)=0.
%\]
%Le matrici $\gamma^\mu$ soddisfano
%\[
% \{\gamma^\mu,\gamma^\nu\}=2g^{\mu\nu}\mathbf1_4,
%\]
%con $g^{\mu\nu}=\mathrm{diag}(1,-1,-1,-1)$.
%\end{column}
%\begin{column}{0.48\textwidth}
%Una soluzione a energia positiva ha la forma
%\[
% \psi(x)=u_s(\mathbf p)e^{-ip\cdot x},
% \qquad p^\mu=(E,\mathbf p),
%\]
%\[
% (\gamma^\mu p_\mu-m)u_s(\mathbf p)=0,
% \qquad E=\sqrt{\mathbf p^2+m^2}.
%\]
%$u_s$ è uno spinore numerico a quattro componenti; $s$ specifica lo spin.
%\end{column}
%\end{columns}
%\medskip
%Sotto una trasformazione di Lorentz $x'=\Lambda x$, il campo trasforma come
%$\psi'(x')=S(\Lambda)\psi(x)$: non è un quadrivettore.
%\end{frame}

\begin{frame}{Coniugato di Dirac e bilineari covarianti}
\small
\begin{columns}[T,onlytextwidth]
\begin{column}{0.48\textwidth}
Il coniugato hermitiano $\psi^\dagger$ non produce, da solo, uno scalare di Lorentz.
Si introduce
\[
 \bar\psi\equiv\psi^\dagger\gamma^0.
\]
Dalla proprietà
\[
 S^\dagger\gamma^0=\gamma^0S^{-1}
\]
segue $\bar\psi'(x')=\bar\psi(x)S^{-1}$.
\end{column}
\begin{column}{0.48\textwidth}
Esempi di bilineari:
\[
 \bar\psi\psi \quad\text{(scalare)},
\]
\[
 j^\mu=\bar\psi\gamma^\mu\psi
 \quad\text{(quadrivettore)}.
\]
In particolare, $j^0=\psi^\dagger\psi$ è una densità nel sistema di riferimento scelto, non uno scalare di Lorentz.
\end{column}
\end{columns}
\medskip
La notazione $\bar\psi\psi$ ricomparirà nel termine di massa, nel condensato chirale e nell'accoppiamento scalare della RMF.
\end{frame}

\subsection{Proiezione chirale}

\begin{frame}{Chiralità: proiettori e autospazi di $\gamma^5$}
\small
\begin{columns}[T,onlytextwidth]
\begin{column}{0.48\textwidth}
Definiamo
\[
 \gamma^5=i\gamma^0\gamma^1\gamma^2\gamma^3,
 \quad (\gamma^5)^2=\mathbf1_4.
\]
I proiettori sono
\[
 P_L=\frac{1-\gamma^5}{2},\qquad
 P_R=\frac{1+\gamma^5}{2}.
\]
Valgono $P_L+P_R=1$ e $P_LP_R=0$.
\end{column}
\begin{column}{0.48\textwidth}
Per un campo o uno spinore $\psi$:
\[
 \psi_L=P_L\psi,\quad \psi_R=P_R\psi,
 \quad \psi=\psi_L+\psi_R.
\]
Inoltre
\[
 \gamma^5\psi_L=-\psi_L,\qquad
 \gamma^5\psi_R=+\psi_R.
\]
La decomposizione è univoca; il fattore $1/2$ nei proiettori non impone norme uguali.
\end{column}
\end{columns}
\medskip
Chiralità ed elicità sono concetti distinti: per stati liberi di particella massless, a energia positiva, coincidono.
\end{frame}

\begin{frame}{Esempio: proiettori nella rappresentazione di Weyl}
\small
\begin{columns}[T,onlytextwidth]
\begin{column}{0.48\textwidth}
Nella rappresentazione chirale:
\[
 \gamma^5=\begin{pmatrix}-\mathbf1_2&0\\0&\mathbf1_2\end{pmatrix},
 \quad
 \gamma^0=\begin{pmatrix}0&\mathbf1_2\\\mathbf1_2&0\end{pmatrix}.
\]
Per $\psi=(\chi_L,\chi_R)^T$, con $\chi_{L,R}$ spinori a due componenti,
\[
 P_L\psi=\binom{\chi_L}{0},\quad
 P_R\psi=\binom{0}{\chi_R}.
\]
\end{column}
\begin{column}{0.48\textwidth}
I proiettori selezionano sottospazi, non quote fisse della norma:
\[
 \|\psi_L\|^2=\psi^\dagger P_L\psi,
 \quad \|\psi_R\|^2=\psi^\dagger P_R\psi.
\]
Vale
\[
 \|\psi_L\|^2+\|\psi_R\|^2=\|\psi\|^2.
\]
Esempio: $\chi_R=0$ implica $\|\psi_L\|^2=\|\psi\|^2$.
\end{column}
\end{columns}
\end{frame}

\begin{frame}{Coniugazione chirale e termine scalare}
\small
\begin{columns}[T,onlytextwidth]
\begin{column}{0.48\textwidth}
Poiché $\{\gamma^0,\gamma^5\}=0$:
\[
 P_L\gamma^0=\gamma^0P_R,\quad
 P_R\gamma^0=\gamma^0P_L.
\]
Pertanto
\[
 \overline{\psi_L}=\bar\psi P_R,\qquad
 \overline{\psi_R}=\bar\psi P_L.
\]
Il pedice $L$ del coniugato si riferisce al campo originario $\psi_L$.
\end{column}
\begin{column}{0.48\textwidth}
Ne segue
\[
 \bar\psi_L\psi_L=\bar\psi P_RP_L\psi=0,
\]
\[
 \bar\psi_R\psi_R=0.
\]
Quindi
\[
 \bar\psi\psi=\bar\psi_L\psi_R+
 \bar\psi_R\psi_L.
\]
In base di Weyl:
$\bar\psi\psi=\chi_R^\dagger\chi_L+\chi_L^\dagger\chi_R$.
\end{column}
\end{columns}
\end{frame}

\begin{frame}{Termini scalari e vettoriali: comportamento chirale}
\small
\begin{columns}[T,onlytextwidth]
\begin{column}{0.48\textwidth}
Il termine scalare accoppia $L$ e $R$:
\[
 -m\bar\psi\psi
 =-m(\bar\psi_L\psi_R+\bar\psi_R\psi_L).
\]
Le equazioni chirali libere sono
\[
 i\gamma^\mu\partial_\mu\psi_L=m\psi_R,
\]
\[
 i\gamma^\mu\partial_\mu\psi_R=m\psi_L.
\]
Per $m=0$ risultano disaccoppiate.
\end{column}
\begin{column}{0.48\textwidth}
La corrente vettoriale conserva i due settori:
\[
 \bar\psi\gamma^\mu\psi=
 \bar\psi_L\gamma^\mu\psi_L+
 \bar\psi_R\gamma^\mu\psi_R.
\]
Questa struttura compare nelle interazioni QED e QCD e nel campo $\omega_\mu$ della RMF.
\end{column}
\end{columns}
\medskip
Un campo scalare efficace $\sigma$ compare invece nel termine $g_\sigma\sigma\bar\psi\psi$, con la stessa struttura chirale del termine di massa.
\end{frame}

\subsection{Scomposizione in modi del campo di Dirac libero - Dal modo libero al campo quantizzato}

\begin{frame}{Equazione di Dirac e modi a momento definito}
\small
Per un campo libero, nella rappresentazione di Dirac,
\[
 i\partial_t\psi=H_D\psi,\qquad
 H_D=\boldsymbol\alpha\cdot(-i\boldsymbol\nabla)+\beta m.
\]
\begin{columns}[T,onlytextwidth]
\begin{column}{0.49\textwidth}
Scomponiamo il campo in onde piane spaziali:
\[
 \psi(t,\mathbf x)=\int\!\frac{d^3p}{(2\pi)^3}
 e^{i\mathbf p\cdot\mathbf x}\,\psi_{\mathbf p}(t).
\]
Per ogni $\mathbf p$:
\[
 i\partial_t\psi_{\mathbf p}(t)
 =H_D(\mathbf p)\psi_{\mathbf p}(t).
\]
\end{column}
\begin{column}{0.49\textwidth}
Poiche' $H_D(\mathbf p)^2=E_{\mathbf p}^2\mathbf1_4$,
\[
 E_{\mathbf p}=\sqrt{|\mathbf p|^2+m^2},
\]
ci sono quattro autovettori:
\[
 \begin{array}{c|c}
 \text{autovalore}&\text{molteplicita'}\\ \hline
 +E_{\mathbf p}&2\\
 -E_{\mathbf p}&2
 \end{array}
\]
L'indice $s=1,2$ distingue i due stati di spin.
\end{column}
\end{columns}
\medskip
La soluzione generale si ottiene sovrapponendo tutti i modi,
positivi e negativi, e tutti i momenti $\mathbf p$.
\end{frame}

%----------------------------------------------------------
\begin{frame}{I quattro modi indipendenti: spinori $u_s$ e $v_s$}
\small
Per $p^\mu=(E_{\mathbf p},\mathbf p)$, con $p^2=m^2$,
\[
 (\gamma^\mu p_\mu-m)u_s(p)=0,\qquad
 (\gamma^\mu p_\mu+m)v_s(p)=0.
\]
\begin{columns}[T,onlytextwidth]
\begin{column}{0.49\textwidth}
\textbf{Soluzioni a frequenza positiva}
\[
 u_s(p)=\begin{pmatrix}
 \sqrt{E+m}\,\xi_s\\[2mm]
 \dfrac{\boldsymbol\sigma\cdot\mathbf p}{\sqrt{E+m}}\,\xi_s
 \end{pmatrix}.
\]
\end{column}
\begin{column}{0.49\textwidth}
\textbf{Soluzioni a frequenza negativa}
\[
 v_s(p)=\begin{pmatrix}
 \dfrac{\boldsymbol\sigma\cdot\mathbf p}{\sqrt{E+m}}\,\eta_s\\[2mm]
 \sqrt{E+m}\,\eta_s
 \end{pmatrix}.
\]
\end{column}
\end{columns}
\medskip
Qui $E=E_{\mathbf p}$; $\xi_s,\eta_s\in\mathbb C^2$ sono basi
ortonormali di spinori di Pauli, per esempio
$\xi_1=(1,0)^T$ e $\xi_2=(0,1)^T$.
\[
 u_s^\dagger(p)u_{s'}(p)=v_s^\dagger(p)v_{s'}(p)
 =2E\,\delta_{ss'}.
\]
\end{frame}

%----------------------------------------------------------
\begin{frame}{Dalla soluzione generale alla scomposizione in modi}
\small
\begin{columns}[T,onlytextwidth]
\begin{column}{0.49\textwidth}
Per ogni momento, la soluzione dell'equazione
$ i\partial_t\psi_{\mathbf p}=H_D(\mathbf p)\psi_{\mathbf p}$
contiene due frequenze:
\[
 e^{-iEt}\quad(+E),\qquad
 e^{+iEt}\quad(-E).
\]
$u_s(\mathbf p)$ sono autovettori di $H_D(\mathbf p)$
con autovalore $+E$.
\end{column}
\begin{column}{0.49\textwidth}
Gli spinori $v_s(-\mathbf p)$ sono autovettori
di $H_D(\mathbf p)$ con autovalore $-E$.

Rietichettando $\mathbf p\to-\mathbf p$ nella parte
a frequenza negativa, questa assume la forma
$ v_s(p)e^{+ip\cdot x}$.

Definiamo $p\cdot x=Et-\mathbf p\cdot\mathbf x$.
\end{column}
\end{columns}
\medskip
Per un campo spinoriale libero (ancora non quantizzato):
\[
 \psi(x)=\sum_{s=1}^{2}\int\frac{d^3p}{(2\pi)^3\sqrt{2E_{\mathbf p}}}
 \left[c_s(\mathbf p)u_s(p)e^{-ip\cdot x}
 +d_s(\mathbf p)v_s(p)e^{+ip\cdot x}\right].
\]
I coefficienti $c_s,d_s$ specificano la sovrapposizione dei modi;
la scelta $1/\sqrt{2E_{\mathbf p}}$ fissa una normalizzazione comoda.
\end{frame}

%----------------------------------------------------------
\begin{frame}{Quantizzazione: operatori $a_s$ e $b_s^\dagger$}
\small
\begin{columns}[T,onlytextwidth]
\begin{column}{0.49\textwidth}
Nel passaggio alla teoria quantistica dei campi,
le componenti di $\psi(x)$ diventano operatori.

Il coefficiente del modo di particella diventa
$a_s(\mathbf p)$, che annichila una particella.

Quello del modo a frequenza negativa diventa
$b_s^\dagger(\mathbf p)$, che \emph{crea} un'antiparticella.
\end{column}
\begin{column}{0.49\textwidth}
Le anticommutazioni non nulle sono
\[
 \begin{aligned}
 \{a_s(\mathbf p),a_{s'}^\dagger(\mathbf p')\}
 &=(2\pi)^3\delta_{ss'}\delta^3(\mathbf p-\mathbf p'),\\
 \{b_s(\mathbf p),b_{s'}^\dagger(\mathbf p')\}
 &=(2\pi)^3\delta_{ss'}\delta^3(\mathbf p-\mathbf p').
 \end{aligned}
\]
Gli altri anticommutatori si annullano.
\end{column}
\end{columns}
\medskip
La scomposizione del \textbf{campo quantizzato} e'
\[
 \psi(x)=\sum_{s=1}^{2}\int\frac{d^3p}{(2\pi)^3\sqrt{2E_{\mathbf p}}}
 \bigl[a_s(\mathbf p)u_s(p)e^{-ip\cdot x}
 +b_s^\dagger(\mathbf p)v_s(p)e^{+ip\cdot x}\bigr].
\]
\end{frame}

%----------------------------------------------------------
\begin{frame}{Perche' $b_s^\dagger$ anziche' un annichilatore?}
\small
\begin{columns}[T,onlytextwidth]
\begin{column}{0.49\textwidth}
Nell'equazione classica sono presenti modi a
frequenza $\pm E$. Dopo la quantizzazione, non
li interpretiamo come particelle di energia fisica negativa.

Per un campo fermionico libero, l'hamiltoniana
normal-ordinata diventa
\[
 :H:=\sum_s\int\!\frac{d^3p}{(2\pi)^3}
 E_{\mathbf p}(a_s^\dagger a_s+b_s^\dagger b_s).
\]
Particelle e antiparticelle hanno entrambe energia positiva.
\end{column}
\begin{column}{0.49\textwidth}
L'operatore di carica conservata, in unita' di
carica del campo, ha invece la struttura
\[
 :Q:=\sum_s\int\!\frac{d^3p}{(2\pi)^3}
 (a_s^\dagger a_s-b_s^\dagger b_s).
\]
I due tipi di eccitazione hanno carica opposta.

Per il campo del quark, il numero barionico
si ottiene moltiplicando questo operatore per $1/3$.
\end{column}
\end{columns}
\medskip
La scelta $b_s^\dagger$ segue dalla quantizzazione fermionica:
il settore a frequenza negativa viene reinterpretato come
creazione di antiparticelle di energia positiva.
\end{frame}

%----------------------------------------------------------
\begin{frame}{Completezza dei modi e significato del campo}
\small
\begin{columns}[T,onlytextwidth]
\begin{column}{0.49\textwidth}
A momento $\mathbf p$ fissato, gli spinori
$u_s(\mathbf p)$ e $v_s(-\mathbf p)$ completano
lo spazio spinoriale:
\[
 \sum_{s=1}^{2}\left[
 u_s(\mathbf p)u_s^\dagger(\mathbf p)
 +v_s(-\mathbf p)v_s^\dagger(-\mathbf p)
 \right]=2E_{\mathbf p}\,\mathbf1_4.
\]
Per questo nessuna delle quattro soluzioni indipendenti
puo' essere omessa dallo sviluppo generale.
\end{column}
\begin{column}{0.49\textwidth}
Gli anticommutatori dei modi e la completezza
garantiscono, a tempi uguali,
\[
 \{\psi_\alpha(t,\mathbf x),
 \psi_\beta^\dagger(t,\mathbf y)\}
 =\delta_{\alpha\beta}\delta^3(\mathbf x-\mathbf y).
\]
$\alpha,\beta=1,\ldots,4$ sono gli indici spinoriali.
\end{column}
\end{columns}
\medskip
Il campo $\psi(x)$ e' un operatore spinoriale; $u_s,v_s$
sono invece spinori numerici che descrivono i modi liberi.
In QCD, per $q_{fa\alpha}(x)$ si aggiungono gli indici
di sapore $f$ e colore $a$.
\end{frame}

\begin{frame}{Spinore numerico e operatore di campo}
\small
\begin{columns}[T,onlytextwidth]
\begin{column}{0.48\textwidth}
Nel problema a una particella $u_s(\mathbf p)$ e $v_s(\mathbf p)$ sono spinori numerici, soluzioni dell'equazione di Dirac.
\end{column}
\begin{column}{0.48\textwidth}
Il campo $q_\alpha(x)$, dopo la quantizzazione, è invece un operatore per ogni indice $\alpha=1,\ldots,4$.
\end{column}
\end{columns}
\medskip
Un'espansione in modi per un campo fermionico libero è
\[
 q(x)=\sum_s\int\frac{d^3p}{(2\pi)^3\sqrt{2E_p}}
 \left[a_s(\mathbf p)u_s(\mathbf p)e^{-ip\cdot x}
 +b_s^\dagger(\mathbf p)v_s(\mathbf p)e^{ip\cdot x}\right].
\]
$a_s$ annichila una particella; $b_s^\dagger$ crea un'antiparticella. La normalizzazione è $u_s^\dagger u_s=v_s^\dagger v_s=2E_p$.

Lo stesso proiettore agisce su spinori e campi: $u_L=P_Lu$ e $q_L=P_Lq$.
\end{frame}

\begin{frame}{Come un campo produce uno spinore di particella}
\small
\begin{columns}[T,onlytextwidth]
\begin{column}{0.48\textwidth}
Per uno stato libero $|\mathbf p,s\rangle$ con una particella, un elemento di matrice del campo è proporzionale al modo spinoriale:
\[
 \langle0|q(x)|\mathbf p,s\rangle
 \propto u_s(\mathbf p)e^{-ip\cdot x}.
\]
La proporzionalità dipende dalla normalizzazione degli stati.
\end{column}
\begin{column}{0.48\textwidth}
In QCD $u(x)$ e $d(x)$ indicano i campi dei sapori up e down, non lo spinore libero $u_s(\mathbf p)$.

Un campo di quark ha indici
\[
 q_{f a\alpha}(x),
\]
con sapore $f$, colore $a=1,2,3$, indice di Dirac $\alpha=1,\ldots,4$.
\end{column}
\end{columns}
\medskip
La quantità $u^\dagger P_Lu$ misura un peso chirale del modo libero; $\langle0|\bar q_Lq_R|0\rangle$ è invece un condensato di operatori di campo nel vuoto interagente.
\end{frame}

\section{QCD e simmetria chirale}
\subsection{Interazioni fondamentali e sapori}

\begin{frame}{Da Dirac a QED e QCD: derivata covariante}
\small
\begin{columns}[T,onlytextwidth]
\begin{column}{0.48\textwidth}
Per un fermione libero:
\[
 \mathcal L_D=\bar\psi(i\gamma^\mu\partial_\mu-m)\psi.
\]
In QED ($U(1)$):
\[
 D_\mu=\partial_\mu+ieQA_\mu,
\]
\[
 \mathcal L_{\rm int}^{\rm QED}
 =-eQ\bar\psi\gamma^\mu\psi A_\mu.
\]
$e>0$; $Q$ è la carica in unità di $e$ (es.: per l'elettrone $Q=-1$).
\end{column}
\begin{column}{0.48\textwidth}
In QCD ($SU(3)_c$):
\[
 D_\mu=\partial_\mu+ig_sT^aG_\mu^a,
\]
\[
 \mathcal L_{\rm int}^{\rm QCD}
 =-g_s\bar q\gamma^\mu T^a q\,G_\mu^a.
\]
$g_s$: accoppiamento forte; $T^a$: generatori del colore; $G_\mu^a$: otto gluoni ($a=1,\ldots,8$).
\end{column}
\end{columns}
\medskip
In QED il fotone è elettricamente neutro; nella QCD non abeliana i gluoni portano colore e si auto-interagiscono. Si aggiungono le lagrangiane dei campi di gauge.
\end{frame}

\begin{frame}{Sapore, colore e chiralità dei campi di quark}
\small
\begin{columns}[T,onlytextwidth]
\begin{column}{0.48\textwidth}
Per i sapori leggeri definiamo il doppietto
\[
 q(x)=\binom{u(x)}{d(x)},\qquad
 q_{L,R}=P_{L,R}q.
\]
Le sue componenti sono
\[
 q_L=\binom{u_L}{d_L},\qquad
 q_R=\binom{u_R}{d_R}.
\]
Ogni sapore ha entrambe le chiralità.
\end{column}
\begin{column}{0.48\textwidth}
Tre strutture distinte:
\begin{itemize}
\item $P_L,P_R$ agiscono sugli indici di Dirac;
\item le matrici di sapore $L,R$ su $(u,d)$;
\item $SU(3)_c$ sui tre colori.
\end{itemize}
Il gruppo $SU(2)$ che segue riguarda il \emph{sapore}, non lo spin.
\end{column}
\end{columns}
\end{frame}

\subsection{Simmetria chirale e termini di massa}

\begin{frame}{Simmetria chirale della lagrangiana QCD}
\small
\begin{columns}[T,onlytextwidth]
\begin{column}{0.48\textwidth}
Per i due sapori $u,d$:
\[
 \mathcal L_q=\bar q(i\gamma^\mu D_\mu-M)q,
 \qquad M=\begin{pmatrix}m_u&0\\0&m_d\end{pmatrix}.
\]
Nel limite $m_u=m_d=0$:
\[
 \mathcal L_q=\bar q_Li\gamma^\mu D_\mu q_L
 +\bar q_Ri\gamma^\mu D_\mu q_R.
\]
\end{column}
\begin{column}{0.48\textwidth}
Sono ammesse trasformazioni globali indipendenti:
\[
 q_L\mapsto Lq_L,\qquad q_R\mapsto Rq_R,
\]
\[
 L,R\in SU(2),\qquad
 L^\dagger L=R^\dagger R=1.
\]
La simmetria risultante è $SU(2)_L\times SU(2)_R$.
\end{column}
\end{columns}
\medskip
La derivata covariante agisce sul colore ma non mescola i sapori: per $L$ globale vale $D_\mu(Lq_L)=L D_\mu q_L$ e il termine sinistro è invariante per $L^\dagger L=1$ (analogamente il destro).
\end{frame}

\begin{frame}{Esempio concreto di $L\neq R$}
\small
\begin{columns}[T,onlytextwidth]
\begin{column}{0.48\textwidth}
Una possibile trasformazione assiale è
\[
 L=e^{i\theta\tau^3/2},\qquad
 R=e^{-i\theta\tau^3/2},
\]
con $\tau^3=\mathrm{diag}(1,-1)$.

Per esempio:
\[
 u_L\mapsto e^{i\theta/2}u_L,
 \qquad u_R\mapsto e^{-i\theta/2}u_R.
\]
Le fasi sulle due chiralità sono opposte.
\end{column}
\begin{column}{0.48\textwidth}
Nel termine cinetico si cancellano le matrici unitarie:
\[
 \bar q_L L^\dagger i\!\not\!D\,Lq_L
 =\bar q_L i\!\not\!D\,q_L.
\]
Nel termine misto compare invece
\[
 \bar q_Lq_R\mapsto
 \bar q_LL^\dagger Rq_R.
\]
Per $L\neq R$ questo termine, in generale, cambia.
\end{column}
\end{columns}
\end{frame}

\begin{frame}{Simmetria vettoriale, assiale e rottura esplicita}
\small
\begin{columns}[T,onlytextwidth]
\begin{column}{0.48\textwidth}
Trasformazioni vettoriali:
\[
 L=R=V,\qquad q\mapsto Vq.
\]
Esse costituiscono il sottogruppo diagonale $SU(2)_V$ (isospin).

Per una trasformazione puramente assiale si può scegliere $L=A$, $R=A^\dagger$.
\end{column}
\begin{column}{0.48\textwidth}
Per masse uguali non nulle $M=m_q\mathbf1$:
\[
 \mathcal L_m=-m_q(\bar q_Lq_R+\bar q_Rq_L).
\]
Questo termine impone $L^\dagger R=1$ per conservarne l'invarianza; resta quindi $L=R$.

Se $m_u\neq m_d$, anche l'isospin è leggermente rotto.
\end{column}
\end{columns}
\medskip
La presenza di $M$ realizza una \emph{rottura esplicita}: la lagrangiana stessa non è più invariante sotto tutte le trasformazioni chirali.
\end{frame}

\subsection{Condensato chirale e pioni}

\begin{frame}{Condensato chirale: simmetria del vuoto}
\small
\begin{columns}[T,onlytextwidth]
\begin{column}{0.48\textwidth}
Anche per $M=0$, il vuoto interagente può avere
\[
 \langle\bar q q\rangle_0\neq0.
\]
Introduciamo la matrice nello spazio dei sapori
\[
 \Phi_{ij}=\langle\bar q_{R,j}q_{L,i}\rangle_0,
 \quad i,j=u,d.
\]
Essa trasforma come $\Phi\mapsto L\Phi R^\dagger$.
\end{column}
\begin{column}{0.48\textwidth}
Se il vuoto seleziona
\[
 \Phi_0=v\mathbf1_2,\qquad v\neq0,
\]
una sua simmetria deve soddisfare
\[
 L\Phi_0R^\dagger=\Phi_0
 \quad\Rightarrow\quad L=R.
\]
Pertanto:
\[
 SU(2)_L\times SU(2)_R\longrightarrow SU(2)_V.
\]
\end{column}
\end{columns}
\medskip
Questa è una \emph{rottura spontanea}: la lagrangiana conserva la simmetria completa, ma il vuoto no. Il condensato è un valore di aspettazione di operatori, non una norma spinoriale.
\end{frame}

\begin{frame}{Teorema di Goldstone e tripletto pionico}
\small
\begin{columns}[T,onlytextwidth]
\begin{column}{0.48\textwidth}
Il gruppo iniziale possiede $3+3=6$ generatori; il sottogruppo vettoriale conservato ne possiede $3$.
\[
 N_{\mathrm{rotti}}=6-3=3.
\]
Il teorema di Goldstone associa ai tre generatori assiali rotti tre modi massless nel limite chirale.
\end{column}
\begin{column}{0.48\textwidth}
In QCD i modi sono i pioni, con
\[
 I=1,\qquad J^P=0^-,
\]
\[
 \pi^+,\quad\pi^0,\quad\pi^-.
\]
In presenza di masse correnti $m_u,m_d\neq0$, i pioni acquistano massa: sono \emph{pseudo-Goldstone}.
\end{column}
\end{columns}
\end{frame}

\begin{frame}{Campi pionici come oscillazioni del condensato}
\small
\begin{columns}[T,onlytextwidth]
\begin{column}{0.48\textwidth}
Le orientazioni del condensato sono descritte da
\[
 U(x)=\exp\!\left(\frac{i\tau^a\pi^a(x)}{f_\pi}\right),
\]
con $a=1,2,3$, matrici di Pauli $\tau^a$, campi reali $\pi^a(x)$ e $f_\pi\simeq92\,\mathrm{MeV}$.

Per piccoli campi: $U\simeq1+i\tau^a\pi^a/f_\pi+\cdots$.
\end{column}
\begin{column}{0.48\textwidth}
La lagrangiana efficace inizia con
\[
 \mathcal L_{\mathrm{eff}}=\frac{f_\pi^2}{4}
 \operatorname{Tr}(\partial_\mu U\partial^\mu U^\dagger)+\cdots.
\]
Al secondo ordine:
\[
 \mathcal L_{\mathrm{eff}}=\frac12\partial_\mu\pi^a\partial^\mu\pi^a+\cdots.
\]
Nel limite chirale non compare un termine $m_\pi^2(\pi^a)^2$.
\end{column}
\end{columns}
\medskip
I campi carichi si ottengono da $\pi^\pm=(\pi^1\mp i\pi^2)/\sqrt2$; $\pi^0=\pi^3$.
\end{frame}

\begin{frame}{Relazione GMOR: masse dei quark e del pione}
\small
\begin{columns}[T,onlytextwidth]
\begin{column}{0.48\textwidth}
Al primo ordine in $m_u,m_d$, la relazione di Gell-Mann--Oakes--Renner è
\[
 f_\pi^2m_\pi^2\simeq
 -(m_u+m_d)\,\Sigma_q,
\]
con $\Sigma_q\equiv\langle\bar u u\rangle_0
 =\langle\bar d d\rangle_0<0$ condensato \emph{per sapore}, nel limite chirale.
\end{column}
\begin{column}{0.48\textwidth}
In particolare:
\[
 m_\pi^2\propto m_u+m_d.
\]
Se $m_u,m_d\to0$ mantenendo un condensato non nullo, allora $m_\pi\to0$.

I pioni sono leggeri rispetto ai barioni perché risentono della piccola rottura esplicita della simmetria chirale.
\end{column}
\end{columns}
\end{frame}

\begin{frame}{Condensato chirale e materia nucleare densa}
\small
\begin{columns}[T,onlytextwidth]
\begin{column}{0.48\textwidth}
Il condensato scalare
\[
 \langle\bar q q\rangle
\]
segnala la rottura spontanea della simmetria chirale nel vuoto.

Al crescere di temperatura o densità, il parametro d'ordine può diminuire, favorendo la restaurazione chirale.
\end{column}
\begin{column}{0.48\textwidth}
Il condensato pionico è differente:
\[
 \langle\bar q\,i\gamma^5\tau^a q\rangle\neq0.
\]
Possiede numeri quantici pseudoscalari/isovettoriali e può rappresentare un campo pionico macroscopico nel mezzo.

Il deconfinamento riguarda invece i gradi di libertà adronici rispetto a quark e gluoni.
\end{column}
\end{columns}
\medskip
Restaurazione chirale, deconfinamento e condensazione pionica sono fenomeni distinti, anche quando possono essere collegati nella materia fortemente interagente.
\end{frame}

\section{Materia adronica e stelle compatte: RMF}
\subsection{Mean-Field Theory e modello di Walecka}

\begin{frame}{MFT, RMF e Walecka: livelli della descrizione}
\small
\begin{columns}[T,onlytextwidth]
\begin{column}{0.48\textwidth}
\begin{itemize}
\item MFT: approssimazione di un sistema a molti corpi mediante campi medi autoconsistenti.
\item RMF: descrizione relativistica con campi di Dirac barionici e campi mesonici efficaci.
\item Walecka: modello RMF elementare con scalare $\sigma$ e vettore $\omega^\mu$.
\end{itemize}
\end{column}
\begin{column}{0.48\textwidth}
La QCD ha quark $q$ e gluoni $G_\mu^a$ come gradi di libertà fondamentali; la RMF usa barioni $N,\Lambda,\ldots$ e mesoni efficaci.

$\sigma$ produce attrazione scalare; $\omega^0$ repulsione vettoriale. I parametri sono calibrati su osservabili nucleari.
\end{column}
\end{columns}
\medskip
Per le stelle compatte: campi medi $\to$ composizione e densità $\to$ EOS $P(\varepsilon)$ $\to$ integrazione TOV $\to$ relazioni massa--raggio $M(R)$.
\end{frame}

\begin{frame}{Lagrangiana del modello lineare $\sigma$--$\omega$}
\small
\begin{columns}[T,onlytextwidth]
\begin{column}{0.48\textwidth}
Campo nucleonico $N=(p,n)^T$ (doppietto di Dirac):
\[
 \begin{aligned}
 \mathcal L_N=\bar N\big(&i\gamma^\mu\partial_\mu-m_N
 +g_\sigma\sigma\\
 &-g_\omega\gamma^\mu\omega_\mu\big)N.
 \end{aligned}
\]
$m_N$: massa nucleonica; $g_\sigma,g_\omega$: accoppiamenti adimensionali.
\end{column}
\begin{column}{0.48\textwidth}
Campi mesonici:
\[
 \begin{aligned}
 \mathcal L_\sigma={}&\frac12\partial_\mu\sigma\partial^\mu\sigma
 -\frac12m_\sigma^2\sigma^2,\\
 \mathcal L_\omega={}&-\frac14W_{\mu\nu}W^{\mu\nu}
 +\frac12m_\omega^2\omega_\mu\omega^\mu.
 \end{aligned}
\]
$W_{\mu\nu}=\partial_\mu\omega_\nu-\partial_\nu\omega_\mu$; $m_\sigma,m_\omega$: masse mesoniche.
\end{column}
\end{columns}
\medskip
La lagrangiana totale è $\mathcal L=\mathcal L_N+\mathcal L_\sigma+\mathcal L_\omega$; il modello qui è lineare, con accoppiamenti costanti.
\end{frame}

\begin{frame}{Equazioni del moto prima dell'approssimazione}
\small
Le equazioni di Eulero--Lagrange, variando separatamente i campi, forniscono
\[
 \begin{aligned}
 (i\gamma^\mu\partial_\mu-m_N+g_\sigma\sigma
 -g_\omega\gamma^\mu\omega_\mu)N&=0,\\
 (\Box+m_\sigma^2)\sigma&=g_\sigma\bar NN,\\
 \partial_\nu W^{\nu\mu}+m_\omega^2\omega^\mu
 &=g_\omega\bar N\gamma^\mu N.
 \end{aligned}
\]
\begin{columns}[T,onlytextwidth]
\begin{column}{0.48\textwidth}
$\Box=\partial_\mu\partial^\mu$; $\bar NN$ è la sorgente scalare nucleonica.
\end{column}
\begin{column}{0.48\textwidth}
$\bar N\gamma^\mu N$ è la corrente barionica, sorgente del campo vettoriale.
\end{column}
\end{columns}
\medskip
Le equazioni sono accoppiate: i mesoni modificano il moto dei barioni e le densità barioniche generano i campi mesonici.
\end{frame}

\begin{frame}{Campo medio uniforme e spettro nucleonico}
\small
\begin{columns}[T,onlytextwidth]
\begin{column}{0.48\textwidth}
In materia statica, uniforme e isotropa:
\[
 \sigma\to\bar\sigma,\qquad
 \omega^\mu\to(\bar\omega_0,\mathbf0).
\]
Le sorgenti diventano
\[
 n_s\equiv\langle\bar NN\rangle,\qquad
 n_B\equiv\langle N^\dagger N\rangle.
\]
E dunque
\[
 m_\sigma^2\bar\sigma=g_\sigma n_s,
 \quad m_\omega^2\bar\omega_0=g_\omega n_B.
\]
\end{column}
\begin{column}{0.48\textwidth}
L'equazione di Dirac fornisce
\[
 m_N^*=m_N-g_\sigma\bar\sigma,
 \quad V_N=g_\omega\bar\omega_0,
\]
\[
 E_N(k)=\sqrt{k^2+(m_N^*)^2}+V_N.
\]
$m_N^*$ è la massa efficace di Dirac; $V_N$ è il contributo vettoriale alla singola particella.
\end{column}
\end{columns}
\medskip
Qui i valori medi riguardano i campi mesonici e le sorgenti; non si sostituisce il campo fermionico con un unico spinore medio.
\end{frame}

\subsection{Soluzione autoconsistente ed equazione di stato}

\begin{frame}{Densità vettoriale e scalare a $T=0$}
\small
\begin{columns}[T,onlytextwidth]
\begin{column}{0.48\textwidth}
Per materia nucleare \emph{simmetrica}, spin e isospin producono degenerazione $d=4$:
\[
 n_B=\frac{d}{2\pi^2}\int_0^{k_F}\!k^2dk
 =\frac{dk_F^3}{6\pi^2}.
\]
$k_F$ è il momento di Fermi comune di neutroni e protoni, $n_B=n_n+n_p$.
\end{column}
\begin{column}{0.48\textwidth}
La densità scalare è
\[
 n_s=\frac{d}{2\pi^2}
 \int_0^{k_F}\! dk\,k^2
 \frac{m_N^*}{\sqrt{k^2+(m_N^*)^2}}.
\]
Dipende da $m_N^*$, che a sua volta è determinata da $n_s$.
\end{column}
\end{columns}
\medskip
Esempio: $n_B=0.16\,\mathrm{fm}^{-3}$ implica $k_F=1.333\,\mathrm{fm}^{-1}\simeq263\,\mathrm{MeV}$, con $\hbar c=197.327\,\mathrm{MeV\,fm}$.
\end{frame}

\begin{frame}{Autoconsistenza: problema di punto fisso}
\small
\begin{columns}[T,onlytextwidth]
\begin{column}{0.48\textwidth}
A $n_B$ fissata, $k_F$ è noto e
\[
 \bar\omega_0=\frac{g_\omega n_B}{m_\omega^2}.
\]
L'equazione scalare diventa
\[
 m_N^*=m_N-
 \frac{g_\sigma^2}{m_\sigma^2}n_s(m_N^*,k_F).
\]
Si cerca lo zero della funzione
\[
 F(m^*)=m^*-m_N+
 \frac{g_\sigma^2}{m_\sigma^2}n_s(m^*,k_F).
\]
\end{column}
\begin{column}{0.48\textwidth}
Metodo iterativo con rilassamento:
\[
 \begin{aligned}
 \widetilde m^{*(j+1)}&=m_N-
 \frac{g_\sigma^2}{m_\sigma^2}n_s(m^{*(j)},k_F),\\
 m^{*(j+1)}&=(1-\lambda)m^{*(j)}
 +\lambda\widetilde m^{*(j+1)}.
 \end{aligned}
\]
$0<\lambda\le1$; $j$ è l'indice iterativo.

Si verifica la convergenza anche mediante $|F(m^*)|$; bisezione o Brent sono alternative robuste.
\end{column}
\end{columns}
\end{frame}

\begin{frame}{Esempio numerico di iterazione RMF}
\small
\begin{columns}[T,onlytextwidth]
\begin{column}{0.48\textwidth}
Esempio puramente illustrativo (non una parametrizzazione calibrata):
\[
 \begin{gathered}
 n_B=0.16\,\mathrm{fm}^{-3},\quad m_N=939\,\mathrm{MeV},\\
 g_\sigma=10,\quad m_\sigma=550\,\mathrm{MeV},\\
 g_\omega=12,\quad m_\omega=783\,\mathrm{MeV}.
 \end{gathered}
\]
Con $m^{*(0)}=939\,\mathrm{MeV}$ e $\lambda=0.6$ si itera la relazione della slide precedente.
\end{column}
\begin{column}{0.48\textwidth}
\begin{center}
\begin{tabular}{r r}
\hline
$j$ & $m^{*(j)}$ [MeV]\\ \hline
0 & 939.0\\
1 & 700.7\\
2 & 609.4\\
3 & 575.7\\
5 & 559.3\\
10 & 557.0\\ \hline
\end{tabular}
\end{center}
Si ottiene $m_N^*\simeq557\,\mathrm{MeV}$; il campo vettoriale produce $V_N\simeq289\,\mathrm{MeV}$.
\end{column}
\end{columns}
\medskip
I valori dipendono dagli accoppiamenti scelti e non impongono automaticamente la saturazione a questa densità.
\end{frame}

\begin{frame}{Energia e pressione della materia di Walecka}
\small
Con $E^*(k)=\sqrt{k^2+(m_N^*)^2}$ e degenerazione $d=4$:
\begin{columns}[T,onlytextwidth]
\begin{column}{0.48\textwidth}
Densità di energia:
\[
 \begin{aligned}
 \varepsilon={}&\frac{d}{2\pi^2}\int_0^{k_F}
 dk\,k^2E^*(k)\\
 &+\frac12m_\sigma^2\bar\sigma^2
 +\frac12m_\omega^2\bar\omega_0^2.
 \end{aligned}
\]
Il termine scalare e quello vettoriale contribuiscono positivamente a $\varepsilon$ nella forma on-shell qui adottata.
\end{column}
\begin{column}{0.48\textwidth}
Pressione:
\[
 \begin{aligned}
 P={}&\frac{d}{6\pi^2}\int_0^{k_F}
 dk\,\frac{k^4}{E^*(k)}\\
 &-\frac12m_\sigma^2\bar\sigma^2
 +\frac12m_\omega^2\bar\omega_0^2.
 \end{aligned}
\]
Il campo scalare riduce $P$; quello vettoriale la aumenta.
\end{column}
\end{columns}
\medskip
Queste formule valgono per materia uniforme a $T=0$ nel modello lineare con accoppiamenti costanti e approssimazione \emph{no-sea}.
\end{frame}

\begin{frame}{Controllo termodinamico e costruzione della EOS}
\small
\begin{columns}[T,onlytextwidth]
\begin{column}{0.48\textwidth}
A $T=0$ il potenziale chimico nucleonico è
\[
 \mu_N=E_N(k_F)
 =\sqrt{k_F^2+(m_N^*)^2}+V_N.
\]
L'identità termodinamica fornisce un test:
\[
 P=n_B\mu_N-\varepsilon.
\]
Il risultato deve coincidere con la pressione ottenuta dal tensore energia--impulso.
\end{column}
\begin{column}{0.48\textwidth}
Procedura per costruire la EOS:
\begin{enumerate}
\item Si fissa $n_B$ e si risolvono i campi.
\item Si calcolano $\varepsilon(n_B)$ e $P(n_B)$.
\item Si ripete su una griglia di densità.
\end{enumerate}
Eliminando $n_B$ si ottiene $P(\varepsilon)$, necessario per le equazioni TOV.
\end{column}
\end{columns}
\end{frame}

\subsection{Materia in equilibrio beta e iperoni}

\begin{frame}{Dal sistema simmetrico alla materia $npe\mu$}
\small
\begin{columns}[T,onlytextwidth]
\begin{column}{0.48\textwidth}
Nel core stellare neutroni e protoni hanno momenti di Fermi separati $k_{F,n}$ e $k_{F,p}$.

A densità totale fissata:
\[
 n_B=n_n+n_p.
\]
Neutralità elettrica locale:
\[
 n_p=n_e+n_\mu.
\]
Le densità delle specie dipendono dai rispettivi momenti di Fermi.
\end{column}
\begin{column}{0.48\textwidth}
In equilibrio beta, con neutrini liberi di uscire:
\[
 \mu_n=\mu_p+\mu_e,
\]
\[
 \mu_\mu=\mu_e
 \quad\text{se i muoni sono presenti}.
\]
Si introduce normalmente anche il campo isovettoriale $\rho_\mu$, che distingue i contributi dei neutroni e dei protoni.
\end{column}
\end{columns}
\medskip
Si risolvono insieme equazioni dei campi, composizione, neutralità ed equilibrio chimico a ogni $n_B$.
\end{frame}

\begin{frame}{Iperoni nella RMF: esempio della soglia $\Lambda$}
\small
\begin{columns}[T,onlytextwidth]
\begin{column}{0.48\textwidth}
Per il barione $Y$ in un campo $\sigma$--$\omega$:
\[
 m_Y^*=m_Y-g_{\sigma Y}\bar\sigma,
\]
\[
 E_Y(k)=\sqrt{k^2+(m_Y^*)^2}
 +g_{\omega Y}\bar\omega_0.
\]
Una comparsa continua di $\Lambda$ comincia da $k_{F,\Lambda}=0$.
\end{column}
\begin{column}{0.48\textwidth}
All'equilibrio debole:
\[
 \mu_\Lambda=\mu_n.
\]
La condizione di soglia nel modello schematico è
\[
 \mu_n=m_\Lambda^*+g_{\omega\Lambda}\bar\omega_0.
\]
Spesso si usa $g_{\omega\Lambda}/g_{\omega N}\simeq2/3$ dalla simmetria di quark SU(6): è un'ipotesi sugli accoppiamenti, non una conseguenza della MFT.
\end{column}
\end{columns}
\medskip
Per una EOS realistica occorrono ulteriori interazioni, vincoli ipernucleari e la soluzione autoconsistente di tutte le specie presenti.
\end{frame}

\subsection{Letture introduttive}

\begin{frame}{Riferimenti per Mean-Field Theory e iperoni}
\footnotesize
\begin{columns}[T,onlytextwidth]
\begin{column}{0.48\textwidth}
\textbf{A. Schmitt (2009)}\\
\emph{The Physics of Compact Stars}, sez. III.A--B: modello di Walecka e introduzione degli iperoni.\\
\href{https://hep.itp.tuwien.ac.at/~aschmitt/stars.pdf}{Dispense PDF}.

\medskip
\textbf{A. Schmitt (2010)}\\
\emph{Dense Matter in Compact Stars: A Pedagogical Introduction}, cap. 3.\\
\href{https://arxiv.org/abs/1001.3294}{arXiv:1001.3294}.

\medskip
\textbf{A. Sorensen}\\
Materiale didattico e tesi: campo medio e potenziali di Yukawa.\\
\href{https://agnieszkasorensen.com/teaching.html}{Materiale online}.
\end{column}
\begin{column}{0.48\textwidth}
\textbf{I. Vida\~na (2018)}\\
\emph{Hyperons: The Strange Ingredients of the Nuclear Equation of State}. Interazioni $YN$, $YY$, ipernuclei e hyperon puzzle.\\
\href{https://arxiv.org/abs/1803.00504}{arXiv:1803.00504}.

\medskip
\textbf{D. Chatterjee, I. Vida\~na (2016)}\\
\emph{Do Hyperons Exist in the Interior of Neutron Stars?} Composizione, EOS e masse massime.\\
\href{https://arxiv.org/abs/1510.06306}{arXiv:1510.06306}.

\medskip
Lettura progressiva: Schmitt $\to$ Vida\~na $\to$ Chatterjee--Vida\~na.
\end{column}
\end{columns}
\end{frame}

\section{Dalla lagrangiana di Walecka alla EOS}

\subsection{Interazioni efficaci e campi mesonici}

\begin{frame}{Perche' derivare un potenziale efficace?}
\small
\begin{columns}[T,onlytextwidth]
\begin{column}{0.49\textwidth}
\textbf{Due nucleoni}\par\smallskip
L'interazione e' mediata dai campi mesonici
\(\sigma\) (scalare) e \(\omega^\mu\) (vettoriale).
\[
N+N\ \longrightarrow\ N+N
\]
Lo scambio mesonico conduce, nel limite statico,
a un potenziale \(V_{NN}(r)\).
\end{column}
\begin{column}{0.49\textwidth}
\textbf{Materia nucleare}\par\smallskip
In un sistema a molti corpi i nucleoni generano
campi medi da determinare insieme alle densita'.
\[
(\bar\sigma,\bar\omega_0)
\longleftrightarrow (n_s,n_B)
\]
I campi modificano lo spettro e la EOS.
\end{column}
\end{columns}
\medskip
\[
\mathcal L_W\ \longrightarrow\
V_{NN}(r)\quad\text{oppure}\quad
\varepsilon(n_B),\ P(n_B)
\]
\footnotesize
Sorensen (2021): \S 2.1.1, eq. (2.1)--(2.15), Appendice E.
\end{frame}

\begin{frame}{Lagrangiana lineare di Walecka}
\footnotesize
\[
\begin{aligned}
\mathcal L_W={}&\bar\psi(i\gamma^\mu\partial_\mu-m_N)\psi
+g_\sigma\bar\psi\psi\,\sigma
-g_\omega\bar\psi\gamma^\mu\psi\,\omega_\mu\\
&+\frac12\partial_\mu\sigma\,\partial^\mu\sigma
-\frac12m_\sigma^2\sigma^2
-\frac14F_{\mu\nu}F^{\mu\nu}
+\frac12m_\omega^2\omega_\mu\omega^\mu .
\end{aligned}
\]
\begin{columns}[T,onlytextwidth]
\begin{column}{0.49\textwidth}
\textbf{Nucleoni e campo scalare}\par\smallskip
\(\psi\): campo di Dirac del nucleone; \(m_N\): massa.\par
\(\sigma\): campo scalare di massa \(m_\sigma\).\par
\(g_\sigma\): accoppiamento \(\bar\psi\psi\,\sigma\).\par
\[
n_s(x)\equiv\langle\bar\psi\psi\rangle
\]
\end{column}
\begin{column}{0.49\textwidth}
\textbf{Campo vettoriale}\par\smallskip
\(\omega^\mu\): campo vettoriale di massa \(m_\omega\).\par
\(g_\omega\): accoppiamento alla corrente barionica.\par
\[
F_{\mu\nu}=\partial_\mu\omega_\nu-\partial_\nu\omega_\mu,
\quad j^\mu\equiv\langle\bar\psi\gamma^\mu\psi\rangle.
\]
\end{column}
\end{columns}
\smallskip
Convenzioni: \(\hbar=c=1\), metrica \(g_{\mu\nu}=\mathrm{diag}(1,-1,-1,-1)\).
\end{frame}

\begin{frame}{Equazioni di Eulero--Lagrange dei mesoni}
\footnotesize
\begin{columns}[T,onlytextwidth]
\begin{column}{0.49\textwidth}
\textbf{Campo scalare \(\sigma\)}
\[
\mathcal L_\sigma=\tfrac12(\partial\sigma)^2
-\tfrac12m_\sigma^2\sigma^2+g_\sigma\rho_s\sigma
\]
\[
\frac{\partial\mathcal L}{\partial\sigma}
-\partial_\mu\frac{\partial\mathcal L}
{\partial(\partial_\mu\sigma)}=0
\]
\[
(\Box+m_\sigma^2)\sigma=g_\sigma\rho_s
\]
\(\rho_s=\bar\psi\psi\): sorgente scalare
(nell'approssimazione classica).
\end{column}
\begin{column}{0.49\textwidth}
\textbf{Campo vettoriale \(\omega^\mu\)}
\[
\mathcal L_\omega=-\tfrac14F_{\mu\nu}F^{\mu\nu}
+\tfrac12m_\omega^2\omega_\mu\omega^\mu
-g_\omega j^\mu\omega_\mu
\]
\[
\partial_\nu F^{\nu\mu}+m_\omega^2\omega^\mu
=g_\omega j^\mu
\]
\(j^\mu=\bar\psi\gamma^\mu\psi\): sorgente vettoriale.
Questa e' l'equazione di Proca.
\end{column}
\end{columns}
\medskip
\[
\Box\equiv\partial_t^2-\nabla^2,
\qquad \partial_\mu j^\mu=0
\]
Le equazioni sono accoppiate ai nucleoni attraverso le sorgenti.
\end{frame}

\subsection{Derivazione classica del potenziale statico}

\begin{frame}{Sorgenti statiche: il limite non relativistico}
\small
\begin{columns}[T,onlytextwidth]
\begin{column}{0.49\textwidth}
\textbf{Ipotesi}\par\smallskip
Sorgenti quasi ferme: \(\partial_t\simeq0\).\par
Corrente spaziale trascurabile:
\[
j^\mu=(\rho_B,\mathbf0),\qquad
\omega^\mu=(\omega_0,\mathbf0).
\]
Nel limite non relativistico:
\[
\rho_s\simeq\rho_B.
\]
\end{column}
\begin{column}{0.49\textwidth}
\textbf{Equazioni statiche}\par\smallskip
\[
(-\nabla^2+m_\sigma^2)\sigma(\mathbf r)
=g_\sigma\rho_s(\mathbf r)
\]
\[
(-\nabla^2+m_\omega^2)\omega_0(\mathbf r)
=g_\omega\rho_B(\mathbf r)
\]
Per una sorgente nell'origine:
\[
\rho(\mathbf r)=\delta^{(3)}(\mathbf r).
\]
\end{column}
\end{columns}
\medskip
Entrambe le equazioni hanno la forma
\[
(-\nabla^2+m^2)\phi(\mathbf r)=g\,\delta^{(3)}(\mathbf r).
\]
\end{frame}

\begin{frame}{Funzione di Green: dal campo massivo a Yukawa}
\footnotesize
Definiamo \(G_m\) mediante
\[
(-\nabla^2+m^2)G_m(\mathbf r)=\delta^{(3)}(\mathbf r).
\]
\begin{columns}[T,onlytextwidth]
\begin{column}{0.49\textwidth}
\textbf{Trasformata di Fourier}\par\smallskip
\[
G_m(\mathbf r)=\int\frac{d^3k}{(2\pi)^3}
\widetilde G_m(\mathbf k)e^{i\mathbf k\cdot\mathbf r}
\]
Usando \(-\nabla^2\mapsto\mathbf k^2\):
\[
(\mathbf k^2+m^2)\widetilde G_m=1
\]
\[
\widetilde G_m(\mathbf k)
=\frac1{\mathbf k^2+m^2}.
\]
\end{column}
\begin{column}{0.49\textwidth}
\textbf{Antitrasformata}\par\smallskip
\[
G_m(r)=\frac1{2\pi^2r}
\int_0^\infty\!dk\,\frac{k\sin(kr)}{k^2+m^2}
\]
Poiche' l'integrale vale \(\pi e^{-mr}/2\),
\[
G_m(r)=\frac{e^{-mr}}{4\pi r}.
\]
Il campo della sorgente puntiforme e'
\[
\phi(r)=gG_m(r).
\]
\end{column}
\end{columns}
\medskip
\[
\sigma(r)=\frac{g_\sigma e^{-m_\sigma r}}{4\pi r},
\qquad
\omega_0(r)=\frac{g_\omega e^{-m_\omega r}}{4\pi r}.
\]
\end{frame}

\begin{frame}{Energia del secondo nucleone: segni opposti}
\footnotesize
L'equazione di Dirac del nucleone in campi mesonici statici e'
\[
\left[i\gamma^\mu\partial_\mu-m_N+g_\sigma\sigma
-g_\omega\gamma^0\omega_0\right]\psi=0.
\]
\begin{columns}[T,onlytextwidth]
\begin{column}{0.49\textwidth}
\textbf{Hamiltoniana relativistica}\par\smallskip
\[
H_D=\boldsymbol\alpha\!\cdot\!\mathbf p
+\beta(m_N-g_\sigma\sigma)
+g_\omega\omega_0
\]
\(\alpha_i=\gamma^0\gamma^i\), \(\beta=\gamma^0\).\par
Lo scalare riduce la massa efficace:
\[
m_N^*=m_N-g_\sigma\sigma.
\]
\end{column}
\begin{column}{0.49\textwidth}
\textbf{Limite statico di energia positiva}\par\smallskip
\[
U_\sigma(r)=-g_\sigma\sigma(r)
\]
\[
U_\omega(r)=+g_\omega\omega_0(r)
\]
Pertanto lo scambio scalare attrae,
mentre quello vettoriale respinge.
\end{column}
\end{columns}
\medskip
\[
V_{NN}(r)=
-\frac{g_\sigma^2}{4\pi}\frac{e^{-m_\sigma r}}r
+\frac{g_\omega^2}{4\pi}\frac{e^{-m_\omega r}}r.
\]
\end{frame}

\begin{frame}{Energia di due sorgenti: il fattore $1/2$}
\footnotesize
Eliminando i campi statici mediante le rispettive equazioni del moto,
si ottengono termini non locali nell'energia:
\[
\begin{aligned}
E_\sigma&=-\frac{g_\sigma^2}{2}\int d^3x\,d^3y\,
\rho_s(\mathbf x)G_{m_\sigma}(\mathbf x-\mathbf y)\rho_s(\mathbf y),\\
E_\omega&=+\frac{g_\omega^2}{2}\int d^3x\,d^3y\,
\rho_B(\mathbf x)G_{m_\omega}(\mathbf x-\mathbf y)\rho_B(\mathbf y).
\end{aligned}
\]
\begin{columns}[T,onlytextwidth]
\begin{column}{0.49\textwidth}
\textbf{Due sorgenti puntiformi}\par\smallskip
\[
\rho(\mathbf x)=\delta^{(3)}(\mathbf x-\mathbf r_1)
+\delta^{(3)}(\mathbf x-\mathbf r_2)
\]
\end{column}
\begin{column}{0.49\textwidth}
\textbf{Sviluppo della densita'}\par\smallskip
Due autoenergie \((1,1),(2,2)\)
piu' due termini incrociati \((1,2),(2,1)\).
Il loro fattore \(2\) cancella \(1/2\).
\end{column}
\end{columns}
\smallskip
Omettendo le autoenergie, \(E_{\rm int}=V_{NN}(|\mathbf r_1-\mathbf r_2|)\).
\end{frame}

\begin{frame}{Portata e limiti del potenziale $\sigma$--$\omega$}
\small
\begin{columns}[T,onlytextwidth]
\begin{column}{0.49\textwidth}
\textbf{Portata caratteristica}\par\smallskip
In unita' convenzionali:
\[
R_m=\frac{\hbar c}{m_mc^2}.
\]
Con masse illustrative:
\[
\begin{aligned}
m_\sigma c^2&\simeq550\ \mathrm{MeV},\\
m_\omega c^2&\simeq783\ \mathrm{MeV},
\end{aligned}
\]
\[
R_\sigma\simeq0.36\ \mathrm{fm},\quad
R_\omega\simeq0.25\ \mathrm{fm}.
\]
\end{column}
\begin{column}{0.49\textwidth}
\textbf{Interpretazione}\par\smallskip
Se \(m_\sigma<m_\omega\), l'attrazione scalare
ha una portata maggiore della repulsione vettoriale.\par\smallskip
Un nucleo reale richiede anche effetti di
scambio, spin, isospin, contributi pionici
ed eventuali correlazioni a molti corpi.
\end{column}
\end{columns}
\medskip
Il campo \(\sigma\) della RMF e' un canale scalare efficace;
non va necessariamente identificato con un mesone \(\sigma\)
elementare e stabile.
\end{frame}

\subsection{Dalla sorgente statica ai campi medi autoconsistenti}

\begin{frame}{Due approssimazioni distinte: LDA e MFA}
\small
\begin{columns}[T,onlytextwidth]
\begin{column}{0.49\textwidth}
\textbf{Local Density Approximation (LDA)}\par\smallskip
Trascuriamo le derivate dei campi mesonici
rispetto ai termini di massa:
\[
\begin{aligned}
m_\sigma^2\sigma(x)&\simeq g_\sigma\rho_s(x),\\
m_\omega^2\omega^\mu(x)&\simeq g_\omega j^\mu(x).
\end{aligned}
\]
Le densita' possono ancora variare nello spazio.
\end{column}
\begin{column}{0.49\textwidth}
\textbf{Mean-Field Approximation (MFA)}\par\smallskip
Separiamo valore medio e fluttuazioni:
\[
\sigma=\bar\sigma+\delta\sigma,
\qquad\omega^\mu=\bar\omega^\mu+\delta\omega^\mu.
\]
Trascuriamo le fluttuazioni mesoniche:
\[
\begin{aligned}
m_\sigma^2\bar\sigma&=g_\sigma n_s,\\
m_\omega^2\bar\omega^\mu&=g_\omega j^\mu.
\end{aligned}
\]
\end{column}
\end{columns}
\medskip
Per materia infinita uniforme, i gradienti medi sono nulli esattamente:
la loro assenza non richiede una LDA aggiuntiva.
\end{frame}

\begin{frame}{Materia nucleare uniforme: soluzione autoconsistente}
\footnotesize
A \(T=0\), per materia nucleare simmetrica, la degenerazione e' \(\gamma=4\)
(2 spin \(\times\) 2 isospin). Fissiamo \(n_B\):
\[
n_B=\frac{\gamma k_F^3}{6\pi^2},
\qquad E^*(k)=\sqrt{k^2+(m_N^*)^2}.
\]
\begin{columns}[T,onlytextwidth]
\begin{column}{0.49\textwidth}
\textbf{Sorgente scalare}\par\smallskip
\[
n_s(m_N^*,k_F)=\frac{\gamma}{2\pi^2}
\int_0^{k_F}\!dk\,k^2\frac{m_N^*}{E^*(k)}.
\]
\[
m_N^*=m_N-\frac{g_\sigma^2}{m_\sigma^2}n_s.
\]
\end{column}
\begin{column}{0.49\textwidth}
\textbf{Campo vettoriale}\par\smallskip
\[
\bar\omega_0=\frac{g_\omega}{m_\omega^2}n_B,
\qquad V_N=g_\omega\bar\omega_0.
\]
\[
E_N(k)=E^*(k)+V_N.
\]
\end{column}
\end{columns}
\medskip
La massa efficace determina \(n_s\), che determina nuovamente \(m_N^*\):
si cerca lo zero di \(F(m_N^*)\equiv m_N^*-m_N+g_\sigma^2 n_s/m_\sigma^2\).
\end{frame}

\begin{frame}{Dai campi medi alla EOS}
\footnotesize
Per il modello lineare \(\sigma\)--\(\omega\) a \(T=0\):
\[
\begin{aligned}
\varepsilon&=\frac{\gamma}{2\pi^2}\int_0^{k_F}\!dk\,k^2E^*(k)
+\frac12m_\sigma^2\bar\sigma^2
+\frac12m_\omega^2\bar\omega_0^2,\\
P&=\frac{\gamma}{6\pi^2}\int_0^{k_F}\!dk\,\frac{k^4}{E^*(k)}
-\frac12m_\sigma^2\bar\sigma^2
+\frac12m_\omega^2\bar\omega_0^2.
\end{aligned}
\]
\begin{columns}[T,onlytextwidth]
\begin{column}{0.49\textwidth}
\textbf{Controllo termodinamico}\par\smallskip
\[
\mu_N=E^*(k_F)+V_N,
\qquad P=n_B\mu_N-\varepsilon.
\]
\end{column}
\begin{column}{0.49\textwidth}
\textbf{Applicazione stellare}\par\smallskip
Ripetendo per diversi \(n_B\) otteniamo
\(P(\varepsilon)\). Nelle stelle occorre
estendere il modello a materia beta-equilibrata
con neutroni, protoni e leptoni.
\end{column}
\end{columns}
\end{frame}

\subsection{Derivazione quantistica: scattering e propagatori}

\begin{frame}{Scattering $NN$: scambio di un mesone}
\small
\begin{columns}[T,onlytextwidth]
\begin{column}{0.49\textwidth}
\textbf{Cinematica (Sorensen, App. E.1)}
\[
N(p_1)+N(p_2)\to N(p_3)+N(p_4)
\]
\[
q^\mu=p_1^\mu-p_3^\mu=p_4^\mu-p_2^\mu.
\]
Con nucleoni lenti \(|\mathbf p_i|\ll m_N\):
\[
q^0\simeq\frac{\mathbf p_4^2-\mathbf p_2^2}{2m_N}
\ll|\mathbf q|.
\]
\end{column}
\begin{column}{0.49\textwidth}
\textbf{Approssimazioni}\par\smallskip
Si conserva soltanto il diagramma diretto
con scambio di un mesone (livello tree).\par\smallskip
Si trascurano il termine di scambio dovuto
all'identita' dei fermioni e le correzioni
dipendenti da spin e velocita'.
\end{column}
\end{columns}
\medskip
Per un mesone scalare di massa \(m_s\), il propagatore e'
\[
D_F(q)=\frac{i}{q^2-m_s^2+i\epsilon}
\xrightarrow{\ q^0\simeq0\ }
-\frac{i}{\mathbf q^2+m_s^2}.
\]
\end{frame}

\begin{frame}{Scambio scalare: ampiezza e potenziale di Born}
\footnotesize
Il vertice e' determinato da \(\mathcal L_{I,\sigma}=g_\sigma\bar\psi\psi\sigma\).
\begin{columns}[T,onlytextwidth]
\begin{column}{0.49\textwidth}
\textbf{Ampiezza relativistica}\par\smallskip
Con \(S=1+i\mathcal M+\cdots\), nel limite statico:
\[
\mathcal M_\sigma\simeq
\frac{g_\sigma^2}{\mathbf q^2+m_\sigma^2}
(\bar u_3u_1)(\bar u_4u_2).
\]
Per spinori quasi a riposo \(\bar u u\simeq2m_N\):
\[
\mathcal M_\sigma\simeq
\frac{4m_N^2g_\sigma^2}{\mathbf q^2+m_\sigma^2}.
\]
\end{column}
\begin{column}{0.49\textwidth}
\textbf{Confronto con Born}\par\smallskip
Con la normalizzazione \(\bar u u=2m_N\),
\[
\widetilde V(\mathbf q)
=-\frac{\mathcal M(\mathbf q)}{4m_N^2}.
\]
Segue
\[
\widetilde V_\sigma(\mathbf q)
=-\frac{g_\sigma^2}{\mathbf q^2+m_\sigma^2}.
\]
\end{column}
\end{columns}
\medskip
\[
V_\sigma(r)=\int\!\frac{d^3q}{(2\pi)^3}
\,e^{i\mathbf q\cdot\mathbf r}\widetilde V_\sigma(\mathbf q)
=-\frac{g_\sigma^2}{4\pi}\frac{e^{-m_\sigma r}}r.
\]
\end{frame}

\begin{frame}{Scambio vettoriale: il segno repulsivo}
\footnotesize
Il vertice e' determinato da
\(\mathcal L_{I,\omega}=-g_\omega\bar\psi\gamma^\mu\psi\,\omega_\mu\).
\begin{columns}[T,onlytextwidth]
\begin{column}{0.49\textwidth}
\textbf{Corrente vettoriale}\par\smallskip
Per nucleoni lenti:
\[
\bar u(p')\gamma^0u(p)\simeq2m_N,
\qquad\bar u\gamma^iu\ll2m_N.
\]
Il contributo del propagatore di Proca,
con le correnti conservate, produce
un segno opposto allo scambio scalare.
\end{column}
\begin{column}{0.49\textwidth}
\textbf{Potenziale in spazio dei momenti}\par\smallskip
\[
\mathcal M_\omega\simeq
-\frac{4m_N^2g_\omega^2}{\mathbf q^2+m_\omega^2}
\]
\[
\widetilde V_\omega(\mathbf q)
=+\frac{g_\omega^2}{\mathbf q^2+m_\omega^2}.
\]
\end{column}
\end{columns}
\medskip
\[
V_\omega(r)=+\frac{g_\omega^2}{4\pi}\frac{e^{-m_\omega r}}r,
\qquad V_{NN}=V_\sigma+V_\omega.
\]
\footnotesize
Sorensen, Appendice E.2: nel limite statico le correnti spaziali sono trascurabili.
\end{frame}

\begin{frame}{Confronto fra le due derivazioni}
\small
\begin{columns}[T,onlytextwidth]
\begin{column}{0.49\textwidth}
\textbf{Campi classici}\par\smallskip
Equazioni mesoniche con sorgenti statiche
\(\rightarrow\) funzione di Green \(G_m(r)\)\par
\(\rightarrow\) energia del secondo nucleone.
\[
G_m(r)=\frac{e^{-mr}}{4\pi r}.
\]
Il percorso prosegue naturalmente verso
i campi medi autoconsistenti.
\end{column}
\begin{column}{0.49\textwidth}
\textbf{Scattering quantistico}\par\smallskip
Vertici fermione--mesone e propagatori
\(\rightarrow\) ampiezza \(\mathcal M\)\par
\(\rightarrow\) potenziale di Born.
\[
\widetilde G_m(\mathbf q)=\frac1{\mathbf q^2+m^2}.
\]
Mostra il legame con i modi spinoriali
\(u_s(\mathbf p)\) e la QFT.
\end{column}
\end{columns}
\medskip
\[
V_{NN}(r)=\frac{g_\omega^2e^{-m_\omega r}
-g_\sigma^2e^{-m_\sigma r}}{4\pi r}.
\]
Ambedue richiedono il limite statico; il potenziale non include scambio
fermionico, forze tensoriali e correlazioni a molti corpi.
\end{frame}

\begin{frame}{Lettura guidata di Sorensen (2021)}
\small
\begin{columns}[T,onlytextwidth]
\begin{column}{0.49\textwidth}
\textbf{Capitolo 2, \S 2.1.1}\par\smallskip
\begin{itemize}
\item Lagrangiana di Walecka: eq. (2.1).
\item Potenziale efficace: eq. (2.3).
\item Equazioni dei campi: eq. (2.4)--(2.6).
\item LDA e MFA: eq. (2.7)--(2.12).
\item Massa efficace: eq. (2.14).
\end{itemize}
\end{column}
\begin{column}{0.49\textwidth}
\textbf{Appendice E}\par\smallskip
\begin{itemize}
\item E.1: cinematica, propagatore scalare,
ampiezza e potenziale di Born.
\item E.2: estensione al campo vettoriale
\(\omega^\mu\) e potenziale totale.
\item Per i segni del vertice vettoriale,
assumere le convenzioni della eq. (2.1).
\end{itemize}
\end{column}
\end{columns}
\medskip
\textbf{Ordine didattico suggerito:}
classico \(\rightarrow\) RMF/EOS \(\rightarrow\) verifica mediante scattering.
\end{frame}


\section{QCD e materia di quark (8.1)}

\subsection{Gradi di libert\`a e confinamento}

\begin{frame}{Dal nucleo ai quark: obiettivo del capitolo 8}
\small
Le EOS nucleoniche descrivono la materia attraverso nucleoni, iperoni
ed eventuali condensati mesonici. La QCD utilizza invece quark e gluoni.
\begin{columns}[T,onlytextwidth]
\begin{column}{0.49\textwidth}
\textbf{Descrizione adronica}
\begin{itemize}
\item Gradi di libert\`a: $n,p,\Lambda,\Sigma,\ldots$.
\item Interazioni efficaci tra adroni.
\item Utile a densit\`a nucleari e subnucleari.
\end{itemize}
\end{column}
\begin{column}{0.49\textwidth}
\textbf{Descrizione in QCD}
\begin{itemize}
\item Campi di quark $q_f(x)$ e gluoni $G_\mu^a(x)$.
\item Interazione forte con simmetria $SU(3)_c$.
\item Necessaria per discutere il deconfinamento.
\end{itemize}
\end{column}
\end{columns}
\medskip
La domanda per le stelle compatte e se, a densit\`a sufficientemente
alte, i gradi di libert\`a adronici debbano essere sostituiti da quelli
partonici: in entrambi i casi serve una EOS $P(\varepsilon)$ per le TOV.
\end{frame}

\begin{frame}{Adroni, quark e leptoni}
\small
\begin{columns}[T,onlytextwidth]
\begin{column}{0.49\textwidth}
\textbf{Adroni}: stati composti che partecipano
all'interazione forte.
\begin{itemize}
\item Barioni ordinari: tre quark; $p=uud$, $n=udd$.
\item Mesoni ordinari: un quark e un antiquark;
      es. $\pi^+=u\bar d$.
\item Esistono anche stati esotici: tetraquark,
      pentaquark e candidati glueball.
\end{itemize}
La carica elettrica dei quark \`e frazionaria:
\[
Q_u=\frac23e,\qquad Q_d=Q_s=-\frac13e.
\]
Il numero barionico di ogni quark vale $B_q=1/3$.
\end{column}
\begin{column}{0.49\textwidth}
    \textbf{Campi elementari}
\begin{itemize}
\item Quark di sei sapori: $u,d,s,c,b,t$.
\item Gluoni: mediatori dell'interazione forte.
\item Leptoni: $e,\mu,\tau$ e neutrini;
      non contengono quark.
\end{itemize}
\end{column}
\end{columns}
\begin{figure}[!ht]\includegraphics[trim={0cm 0cm 0 0},clip,keepaspectratio,angle=-90,width=0.95\textwidth]{Schaffner/chargemassquarks}\label{fig:Schaffner_chargemassquarks}\end{figure}
\end{frame}

\begin{frame}{QCD: colore e struttura non abeliana}
\small
La simmetria locale della QCD \`e $SU(3)_c$;
i quark trasportano tre stati di colore.
\begin{columns}[T,onlytextwidth]
\begin{column}{0.49\textwidth}
\textbf{Quark e gluoni}
\begin{itemize}
\item Un quark e nella rappresentazione fondamentale
      di $SU(3)_c$.
\item I gluoni sono otto, associati agli
      $N_c^2-1=8$ generatori.
\item Gli stati adronici osservabili sono singoletti di colore.
\end{itemize}
\end{column}
\begin{column}{0.49\textwidth}
\textbf{Confronto con la QED}
\begin{itemize}
\item Il fotone non porta carica elettrica.
\item I gluoni portano colore e possono
      interagire direttamente fra loro.
\item La QCD \`e non abeliana: contiene
      termini di autointerazione dei gluoni.
\end{itemize}
\end{column}
\end{columns}
\medskip
Questa dinamica contribuisce a confinamento e libert\`a asintotica,
che sono fenomeni di scala energetica molto differenti.
\end{frame}

\begin{frame}{Confinamento e scala caratteristica della QCD}
\small
In condizioni ordinarie non si osservano quark isolati come stati
asintotici: le particelle osservabili sono complessivamente neutre di colore.
\begin{columns}[T,onlytextwidth]
\begin{column}{0.49\textwidth}
\textbf{Scala spaziale}
\[
R_{\mathrm{had}}\sim1\ \mathrm{fm}.
\]
Con $p\sim\hbar/R$:
\[
pc\sim\frac{\hbar c}{R}.
\]
\end{column}
\begin{column}{0.49\textwidth}
\textbf{Stima numerica}
\[
\hbar c\simeq197\ \mathrm{MeV\,fm},
\]
\[
pc\sim\frac{197\ \mathrm{MeV\,fm}}{1\ \mathrm{fm}}
\sim200\ \mathrm{MeV}.
\]
\end{column}
\end{columns}
\medskip
\`E una stima dimensionale, non un calcolo della massa del nucleone.
La dinamica confinante non equivale semplicemente a una forza
fondamentale a corto raggio: le forze \emph{residue tra adroni}
sono invece di raggio nucleare finito.
\end{frame}

\begin{frame}{Perch\'e i nucleoni sono molto pi\`u pesanti dei quark correnti?}
\small
La massa del protone non \`e la somma delle masse correnti dei suoi
quark di valenza.
\begin{columns}[T,onlytextwidth]
\begin{column}{0.49\textwidth}
\textbf{Esempio: il protone $uud$}
\[
m_u\simeq2.2\ \mathrm{MeV},\quad
m_d\simeq4.7\ \mathrm{MeV}.
\]
\[
2m_u+m_d\simeq9.1\ \mathrm{MeV}.
\]
\[
m_pc^2\simeq938\ \mathrm{MeV}.
\]
\end{column}
\begin{column}{0.49\textwidth}
\textbf{Origine dell'energia adronica}
\begin{itemize}
\item Dinamica dei quark e gluoni.
\item Energia cinetica e interazioni.
\item Struttura non perturbativa del vuoto QCD.
\end{itemize}
\[
M_{\mathrm{had}}c^2=E_{\mathrm{tot}}
\]
nel sistema di riposo.
\end{column}
\end{columns}
\medskip
Anche per gli iperoni, ad esempio $\Lambda=uds$ con massa
$\sim1116\ \mathrm{MeV}$, la somma delle masse correnti dei
quark non spiega la massa adronica.
\end{frame}

\subsection{Condensati e modello MIT bag}

\begin{frame}{Vuoto QCD e condensato gluonico}
\small
Il vuoto della QCD interagente possiede una struttura quantistica:
non coincide con un vuoto classico privo di campi.
\begin{columns}[T,onlytextwidth]
\begin{column}{0.49\textwidth}
Una quantit\`a frequentemente studiata e il
\textbf{condensato gluonico}, schematicamente
\[
\left\langle\frac{\alpha_s}{\pi}
G^a_{\mu\nu}G^{a\mu\nu}\right\rangle,
\]
dove $G^a_{\mu\nu}$ \`e il tensore del campo
gluonico e $\alpha_s$ l'accoppiamento forte.
\end{column}
\begin{column}{0.49\textwidth}
La struttura del vuoto contribuisce
alle propriet\`a degli adroni.
\begin{itemize}
\item L'energia adronica non \`e soltanto massa
      a riposo dei quark.
\item La massa e legata al tensore energia-impulso
      della QCD e alla sua anomalia di traccia.
\end{itemize}
\end{column}
\end{columns}
\medskip
Il parametro $B$ del MIT bag model rappresenta \emph{fenomenologicamente}
una differenza di densit\`a di energia del vuoto; non va identificato
in modo universale con un singolo condensato gluonico.
\end{frame}

\begin{frame}{MIT bag model: energia e raggio dell'adrone}
\small
Un adrone viene rappresentato da una regione di volume $V$ che
contiene i quark, con condizioni al contorno che li confinano.
\begin{columns}[T,onlytextwidth]
\begin{column}{0.49\textwidth}
\textbf{Energia del modello}
\[
E_{\mathrm{bag}}=BV,
\]
per una sfera $V=4\pi R^3/3$.
In un modello schematico:
\[
E(R)\simeq\frac{C}{R}+\frac{4\pi}{3}BR^3,
\]
con $C>0$ dipendente dai modi dei quark.
\end{column}
\begin{column}{0.49\textwidth}
\textbf{Competizione dei contributi}
\begin{itemize}
\item Energia cinetica dei quark:
      aumenta quando $R$ diminuisce.
\item Energia del volume di bag:
      aumenta quando $R$ cresce.
\end{itemize}
\[
\frac{dE}{dR}=0
\]
seleziona un raggio di equilibrio.
\end{column}
\end{columns}
\medskip
\`E un modello efficace del confinamento, non una soluzione esatta
della QCD a bassissima energia.
\end{frame}

\begin{frame}{Pressione negativa del vuoto: il parametro $B$}
\small
Il contributo $E_{\mathrm{vac}}=BV$ implica
\[
P_{\mathrm{vac}}=-\left(\frac{\partial E_{\mathrm{vac}}}
{\partial V}\right)_{S,N}=-B.
\]
\begin{columns}[T,onlytextwidth]
\begin{column}{0.49\textwidth}
In alternativa, dalla relazione termodinamica
\[
\varepsilon=-P+Ts+\mu n,
\]
per un contributo di vuoto indipendente da
$T,n$: $s=n=0$ e quindi
\[
\varepsilon_{\mathrm{vac}}=-P_{\mathrm{vac}}=B.
\]
\end{column}
\begin{column}{0.49\textwidth}
\textbf{Intervallo citato nel testo}
\[
B^{1/4}=145\text{--}235\ \mathrm{MeV}.
\]
Usando $\hbar c\simeq197.327$ MeV fm:
\[
B[\mathrm{MeV/fm^3}]
=\frac{(B^{1/4}[\mathrm{MeV}])^4}
{(\hbar c)^3}.
\]
\[
B\simeq57.5\text{--}397\ \mathrm{MeV/fm^3}.
\]
\end{column}
\end{columns}
\medskip
Energia positiva del vuoto e pressione negativa non sono in
contraddizione: $B$ \`e una \emph{densit\`a di energia} positiva.
\end{frame}

\begin{frame}{Chiralita: simmetria per quark leggeri massless}
\small
Consideriamo il doppietto di sapore $q=(u,d)^T$ e i proiettori
\[
P_L=\frac{1-\gamma^5}{2},\qquad
P_R=\frac{1+\gamma^5}{2},\qquad q_{L,R}=P_{L,R}q.
\]
\begin{columns}[T,onlytextwidth]
\begin{column}{0.49\textwidth}
\textbf{Limite chirale}
\[
m_u=m_d=0.
\]
La lagrangiana cinetica si separa:
\[
\mathcal L_q=\bar q_Li\!\not\!D q_L+
\bar q_Ri\!\not\!D q_R.
\]
\end{column}
\begin{column}{0.49\textwidth}
\textbf{Trasformazioni globali}
\[
q_L\longrightarrow Lq_L,
\quad q_R\longrightarrow Rq_R,
\]
con $L\in SU(2)_L$ e $R\in SU(2)_R$
indipendenti. Agiscono sul sapore,
non sugli indici di Dirac.
\end{column}
\end{columns}
\medskip
Il termine di massa $-m_q\bar q q=-m_q(\bar q_Lq_R+\bar q_Rq_L)$
collega le chiralita e rompe esplicitamente la simmetria assiale.
\end{frame}

\begin{frame}{Condensato chirale, Goldstone e pioni}
\small
Nel vuoto QCD si forma un condensato chirale:
\[
\langle\bar q q\rangle\neq0,
\qquad\bar q q=\bar q_Lq_R+\bar q_Rq_L.
\]
\begin{columns}[T,onlytextwidth]
\begin{column}{0.49\textwidth}
La lagrangiana massless e invariante sotto
$SU(2)_L\times SU(2)_R$, ma il vuoto
conserva soltanto il sottogruppo vettoriale:
\[
SU(2)_L\times SU(2)_R\to SU(2)_V.
\]
I generatori rotti sono $6-3=3$.
\end{column}
\begin{column}{0.49\textwidth}
Il teorema di Goldstone prevede tre
modi massless nel limite chirale,
identificati con $\pi^+,\pi^0,\pi^-$.
\medskip
Nella QCD reale, $m_u,m_d\ne0$:
i pioni sono \emph{pseudo-Goldstone}
con massa piccola ma non nulla.
\end{column}
\end{columns}
\medskip
Al primo ordine nelle masse leggere, la relazione GMOR e
\[
f_\pi^2m_\pi^2\simeq-(m_u+m_d)\langle\bar q q\rangle_0,
\]
dove il condensato e inteso per singolo sapore nel limite chirale.
\end{frame}

\begin{frame}{Massa corrente, massa costituente e simmetria chirale}
\small
Non confondere le masse dei quark nella lagrangiana
con quelle efficaci nei modelli a quark costituenti.
\begin{columns}[T,onlytextwidth]
\begin{column}{0.49\textwidth}
\textbf{Massa corrente}
\[
m_{u,d}\sim\mathrm{pochi\ MeV}.
\]
\`E il parametro del termine di massa
$-\bar qMq$ della QCD, dipendente
dalla convenzione di rinormalizzazione.
\end{column}
\begin{column}{0.49\textwidth}
\textbf{Massa costituente}
\[
m_{u,d}^{\mathrm{const}}\sim300\ \mathrm{MeV}.
\]
\`E una massa efficace fenomenologica
associata anche alla dinamica di rottura
spontanea chirale.
\end{column}
\end{columns}
\medskip
La restaurazione chirale nel mezzo non implica automaticamente
quark liberi: \emph{deconfinamento} e \emph{restaurazione chirale}
sono nozioni distinte, anche se collegate.
\end{frame}

\begin{frame}{Condensati e regime termico della QCD}
\small
La temperatura modifica le propriet\`a del vuoto e i gradi di
libert\`a termodinamici della materia fortemente interagente.
\begin{columns}[T,onlytextwidth]
\begin{column}{0.49\textwidth}
\textbf{A temperatura crescente}
\begin{itemize}
\item Il condensato chirale diminuisce.
\item La descrizione in termini di adroni
      lascia gradualmente spazio a quark e gluoni.
\item A masse fisiche e $\mu_B\simeq0$, si osserva
      un crossover attorno a $T_c\sim150$ MeV.
\end{itemize}
\end{column}
\begin{column}{0.49\textwidth}
\textbf{Nelle stelle compatte}
\begin{itemize}
\item Regime approssimativamente $T\simeq0$.
\item Potenziale chimico barionico $\mu_B$
      molto elevato.
\item Natura e posizione di eventuali
      transizioni ad alta densit\`a sono incerte.
\end{itemize}
\end{column}
\end{columns}
\medskip
Un crossover termico a $\mu_B\simeq0$ non dimostra una
transizione di primo ordine nel core stellare.
\end{frame}

\subsection{libert\`a asintotica}

\begin{frame}{Accoppiamento forte e scala energetica}
\small
La QCD possiede un accoppiamento che dipende dalla scala
energetica $Q$ del processo: \emph{running coupling}.
\begin{columns}[T,onlytextwidth]
\begin{column}{0.49\textwidth}
Si definisce
\[
\alpha_s(Q)=\frac{g_s^2(Q)}{4\pi},
\]
dove $g_s$ e l'accoppiamento
quark-gluone rinormalizzato.
In QED, per confronto,
\[
\alpha_{\mathrm{EM}}\simeq\frac1{137}.
\]
Il comportamento nasce dalle auto-interazioni dei gluoni
ed e qualitativamente diverso dal running in QED.
\end{column}
\begin{column}{0.49\textwidth}
    \begin{figure}[!ht]\includegraphics[trim={0cm 0cm 0 0},clip,keepaspectratio,width=0.95\textwidth]{Schaffner/QCD-alphas_energyscale}\label{fig:Schaffner_QCD-alphas_energyscale}\end{figure}
    \textbf{Dati illustrativi del libro}
\[
\alpha_s(1.8\ \mathrm{GeV})\sim0.3,
\]
\[
\alpha_s(100\ \mathrm{GeV})\sim0.1.
\]
A scale alte $\alpha_s$ diminuisce:
si apre un regime perturbativo.
\end{column}
\end{columns}
\end{frame}

\begin{frame}{Funzione beta della QCD}
\small
Nello schema di rinormalizzazione $\overline{\mathrm{MS}}$:
\[
\frac{d\alpha_s}{d\ln\bar\Lambda^2}
=\beta(\alpha_s)
=-b_0\alpha_s^2-b_1\alpha_s^3-\cdots.
\]
\begin{columns}[T,onlytextwidth]
\begin{column}{0.49\textwidth}
Il numero $N_f$ indica i sapori
attivi alla scala considerata.
\[
b_0=\frac{33-2N_f}{12\pi},
\]
\[
b_1=\frac{153-19N_f}{24\pi^2}.
\]
\end{column}
\begin{column}{0.49\textwidth}
Per $N_f=3$:
\[
b_0=\frac{27}{12\pi}>0.
\]
Quindi, al primo ordine,
\[
\frac{d\alpha_s}{d\ln\bar\Lambda^2}<0:
\]
$\alpha_s$ diminuisce se cresce la scala.
\end{column}
\end{columns}
\medskip
La scala $\bar\Lambda$ \`e quella di rinormalizzazione;
non \`e il campo iperonico $\Lambda$.
\end{frame}

\begin{frame}{Soluzione al primo ordine: libert\`a asintotica}
\small
Trascurando i termini oltre il primo nella funzione beta:
\[
\frac{d\alpha_s}{d\ln\bar\Lambda^2}=-b_0\alpha_s^2.
\]
\begin{columns}[T,onlytextwidth]
\begin{column}{0.49\textwidth}
\textbf{Separazione delle variabili}
\[
\frac{d\alpha_s}{\alpha_s^2}
=-b_0\,d\ln\bar\Lambda^2.
\]
Integrando:
\[
\frac1{\alpha_s}=b_0\ln\bar\Lambda^2+C.
\]
La costante $C$ introduce una scala
caratteristica $\Lambda_{\overline{\mathrm{MS}}}$.
\end{column}
\begin{column}{0.49\textwidth}
\textbf{Soluzione}
\[
\alpha_s(\bar\Lambda)=
\frac1{b_0\ln(\bar\Lambda^2/
\Lambda_{\overline{\mathrm{MS}}}^2)}.
\]
Per $\bar\Lambda\to\infty$:
\[
\alpha_s\to0.
\]
La crescita formale a bassa scala segnala
il fallimento dell'approssimazione perturbativa.
\end{column}
\end{columns}
\end{frame}

\subsection{Gas relativistico di quark e gluoni}

\begin{frame}{Dal plasma interagente al gas ideale}
\small
A temperature elevate la libert\`a asintotica suggerisce
come primo riferimento un gas di particelle relativistiche.
\begin{columns}[T,onlytextwidth]
\begin{column}{0.49\textwidth}
\textbf{Energia e occupazione}
\[
E(k)=\sqrt{k^2+m^2},
\]
\[
f(k,T)=\frac1{e^{E(k)/T}\pm1},
\]
con $+$ per fermioni e $-$ per bosoni;
qui il potenziale chimico e nullo.
\end{column}
\begin{column}{0.49\textwidth}
\textbf{Pressione cinetica}
\[
P=\gamma\int\frac{d^3k}{(2\pi)^3}
\frac{k^2}{3E(k)}f(k,T).
\]
$\gamma$ \`e la degenerazione interna:
spin, colore, sapore e, se necessario,
particelle/antiparticelle.
\end{column}
\end{columns}
\medskip
Per isotropia l'integrale angolare porta a
\[
P=\frac{\gamma}{6\pi^2}\int_0^\infty dk\,
\frac{k^4}{E(k)}f(k,T).
\]
\end{frame}

\begin{frame}{Limite ultrarelativistico: la legge $T^4$}
\small
Se $m\ll T$, allora $E(k)\simeq k$. Ponendo $x=k/T$:
\[
P=\frac{\gamma T^4}{6\pi^2}
\int_0^\infty\frac{x^3\,dx}{e^x\pm1}.
\]
\begin{columns}[T,onlytextwidth]
\begin{column}{0.49\textwidth}
\textbf{Bosoni}
\[
\int_0^\infty\frac{x^3\,dx}{e^x-1}
=\frac{\pi^4}{15}.
\]
\[
P_b=\gamma\frac{\pi^2}{90}T^4.
\]
\end{column}
\begin{column}{0.49\textwidth}
\textbf{Fermioni}
\[
\int_0^\infty\frac{x^3\,dx}{e^x+1}
=\frac{7\pi^4}{120}.
\]
\[
P_f=\frac78\gamma\frac{\pi^2}{90}T^4.
\]
\end{column}
\end{columns}
\medskip
Il fattore $7/8$ dipende dalla statistica di Fermi--Dirac,
non dall'interazione forte.
\end{frame}

\begin{frame}{Conta dei gradi di libert\`a: gluoni e quark}
\small
In QCD abbiamo $N_c=3$ colori e $N_f$ sapori
sufficientemente leggeri da essere trattati come ultrarelativistici.
\begin{columns}[T,onlytextwidth]
\begin{column}{0.49\textwidth}
\textbf{Gluoni}
\begin{itemize}
\item 2 polarizzazioni fisiche;
\item $N_c^2-1=8$ tipi di gluone.
\end{itemize}
\[
\gamma_g=2(N_c^2-1)=16.
\]
\end{column}
\begin{column}{0.49\textwidth}
\textbf{Quark e antiquark}
\begin{itemize}
\item 2 spin, $N_c=3$ colori;
\item $N_f$ sapori, fattore 2 per $q,\bar q$.
\end{itemize}
\[
\gamma_q=2\cdot3\cdot N_f\cdot2=12N_f.
\]
\end{column}
\end{columns}
\medskip
Nel limite ideale, il fattore fermionico trasforma
$12N_f$ in $(7/8)12N_f=10.5N_f$ gradi di libert\`a
efficaci per la pressione.
\end{frame}

\begin{frame}{Limite di Stefan--Boltzmann per il QGP}
\small
La pressione totale del gas ideale ultrarelativistico \`e
\[
P_{\mathrm{SB}}=\frac{\pi^2T^4}{90}
\left(\sum_{\mathrm{bosoni}}\gamma_i+
\frac78\sum_{\mathrm{fermioni}}\gamma_i\right).
\]
\begin{columns}[T,onlytextwidth]
\begin{column}{0.49\textwidth}
Per il plasma QCD:
\[
P_{\mathrm{SB}}=\frac{\pi^2}{90}
(16+10.5N_f)T^4.
\]
Se contribuiscono $u,d$:
\[
N_f=2\Rightarrow g_{\mathrm{eff}}=37.
\]
\end{column}
\begin{column}{0.49\textwidth}
Includendo anche il quark strange:
\[
N_f=3\Rightarrow g_{\mathrm{eff}}=47.5.
\]
Con $m_s\sim95$ MeV, trattarlo come
massless presso $T_c$ resta
un'approssimazione.
\end{column}
\end{columns}
\medskip
$g_{\mathrm{eff}}$ conta i gradi di libert\`a termici,
non il numero delle specie adroniche in un gas reale.
\end{frame}

\begin{frame}{Confronto dei gradi di libert\`a: adroni e plasma}
\small
A basse temperature, la componente hadronica pi\`u leggera
comprende il tripletto pionico $\pi^+,\pi^0,\pi^-$.
\begin{columns}[T,onlytextwidth]
\begin{column}{0.49\textwidth}
\textbf{Gas ideale di pioni massless}
\[
\gamma_\pi=3,
\]
\[
P_{\pi,\mathrm{SB}}=3\frac{\pi^2}{90}T^4.
\]
Questo modello e soltanto illustrativo:
i pioni reali hanno massa $m_\pi\ne0$.
\end{column}
\begin{column}{0.49\textwidth}
\textbf{QGP con $N_f=2$}
\[
\frac{g_{\mathrm{QGP}}}{g_{\pi}}=
\frac{37}{3}\simeq12.3.
\]
Al crescere di $T$, partecipano pero
molti altri adroni e risonanze,
come nel modello HRG.
\end{column}
\end{columns}
\medskip
Il grande aumento dei gradi di libert\`a spiega
qualitativamente la crescita della pressione attraverso
la regione del crossover, ma non ne determina da solo
la curva quantitativa.
\end{frame}

\subsection{Confronto con lattice QCD e stelle compatte}

\begin{frame}{Lattice QCD: pressione rispetto al limite ideale}
\begin{columns}[T,onlytextwidth]
\begin{column}{0.49\textwidth}
\small
La figura 8.2 di Schaffner confronta diversi metodi tramite $\frac{P(T)}{P_{\mathrm{SB}}(T)}$.
\textbf{Gas ideale}
\[
P/P_{\mathrm{SB}}=1.
\]
\textbf{Lattice QCD}\par
Calcolo numerico non perturbativo:
il rapporto cresce vicino al crossover
ma rimane sotto 1 anche a $T>T_c$.
    \textbf{pQCD}\par
Espansione perturbativa con correzioni
termiche: pi\`u affidabile per $T\gg T_c$.
$P/P_{\mathrm{SB}}<1$ indica che
le interazioni sono ancora importanti
nel plasma deconfinato.
Deconfinamento non significa quindi gas esattamente libero
di quark e gluoni.
\end{column}
\begin{column}{0.49\textwidth}
    \begin{figure}[!ht]\includegraphics[trim={0cm 0cm 0 0},clip,keepaspectratio,width=0.95\textwidth]{Schaffner/pressure_QDC_lattice_pQCD}\label{fig:Schaffner_pressure_QDC_lattice_pQCD}\end{figure}
\end{column}
\end{columns}
\end{frame}

\begin{frame}{Perch\'e il bag model differisce dalla lattice QCD?}
\small
Nel modello pi\`u semplice, il termine di vuoto aggiunge
una pressione negativa costante:
\[
P_{\mathrm{bag}}(T)=P_{\mathrm{SB}}(T)-B.
\]
\begin{columns}[T,onlytextwidth]
\begin{column}{0.49\textwidth}
Dividendo per la pressione ideale:
\[
\frac{P_{\mathrm{bag}}}{P_{\mathrm{SB}}}
=1-\frac{B}{P_{\mathrm{SB}}}.
\]
Poiche $P_{\mathrm{SB}}\propto T^4$:
\[
\frac{P_{\mathrm{bag}}}{P_{\mathrm{SB}}}
\longrightarrow1
\quad(T\to\infty).
\]
\end{column}
\begin{column}{0.49\textwidth}
La figura 8.2 mostra che il bag model
semplice si avvicina molto rapidamente
al limite ideale.
\medskip
La lattice QCD registra invece
deviazioni persistenti dovute a
interazioni del plasma e alla sua
dinamica non perturbativa.
\end{column}
\end{columns}
\medskip
Il parametro $B$ da solo non riproduce la dipendenza reale
$P(T)$ in prossimita del crossover.
\end{frame}

\begin{frame}{Crossover termico e pseudotemperatura critica}
\small
Per masse fisiche dei quark e potenziale chimico barionico
quasi nullo, la trasformazione termica \`e un \emph{crossover}.
\begin{columns}[T,onlytextwidth]
\begin{column}{0.49\textwidth}
\textbf{Cosa significa?}
\begin{itemize}
\item Le grandezze termodinamiche
      cambiano rapidamente.
\item Non e richiesta una discontinuita
      di primo ordine.
\item Non e presente una singolarita
      critica di secondo ordine.
\end{itemize}
\end{column}
\begin{column}{0.49\textwidth}
\textbf{Temperatura pseudocritica}
\[
T_c\sim150\ \mathrm{MeV}.
\]
Si identifica, per esempio, da
massimi di suscettivita o dalla
rapidita di variazione del
condensato chirale.
\end{column}
\end{columns}
\medskip
Il capitolo riporta una crescita rapida della densit\`a di
energia nella regione $T\simeq145$--$163$ MeV:
$\varepsilon\simeq0.18$--$0.5\ \mathrm{GeV/fm^3}$.
\end{frame}

\begin{frame}{Regimi della QCD: plasma caldo e materia stellare}
\small
La fisica del plasma caldo e quella dei core stellari
riguardano regioni diverse del diagramma di fase $(T,\mu_B)$.
\begin{columns}[T,onlytextwidth]
\begin{column}{0.49\textwidth}
\textbf{QGP caldo}
\begin{itemize}
\item $T\gtrsim T_c$, $\mu_B\simeq0$.
\item Occupazioni termiche di Bose/Fermi.
\item Riferimento ideale:
\end{itemize}
\[
P_{\mathrm{SB}}\propto T^4.
\]
\end{column}
\begin{column}{0.49\textwidth}
\textbf{Core di stelle compatte}
\begin{itemize}
\item $T\simeq0$, $\mu_B$ elevato.
\item Fermioni fortemente degeneri.
\item Per quark massless liberi:
\end{itemize}
\[
P\propto\mu_q^4,
\]
con $\mu_q$ potenziale chimico di quark.
\end{column}
\end{columns}
\medskip
Per costruire una EOS stellare servono inoltre
neutralita elettrica, equilibrio beta e interazioni.
Non si puo applicare direttamente la formula $P\propto T^4$.
\end{frame}

\begin{frame}{Dal capitolo 8.1 alla EOS della materia di quark}
\small
Il percorso del capitolo e:
\[
\begin{gathered}
\text{QCD: quark, gluoni e confinamento}\\
\Downarrow\\
\text{Condensati, simmetria chirale e MIT bag model}\\
\Downarrow\\
\text{libert\`a asintotica e pressione del gas ideale}\\
\Downarrow\\
\text{Confronto con il risultato non perturbativo lattice}\\
\Downarrow\\
\text{EOS della materia di quark fredda e densa}.
\end{gathered}
\]
\medskip
Per le stelle compatte, la quantit\`a finale e $P(\varepsilon)$:
la presenza di quark non garantisce di per se una EOS rigida.
\end{frame}

\section{Materia di quark strange libera}

\subsection{Potenziali chimici ed equilibrio debole}

\begin{frame}[t]{Dal capitolo sulla QCD alla materia di quark fredda}
\small
\begin{columns}[T,onlytextwidth]
\begin{column}{0.49\textwidth}
La libert\'a asintotica suggerisce che, a scale energetiche elevate, l'interazione forte diventi pi\'u debole.

Nella materia fredda la scala rilevante non \'e la temperatura, ma il momento di Fermi:
\[
Q\sim k_{F,q}\sim\mu_q.
\]
Si studia quindi un gas degenere di quark \emph{liberi}.
\end{column}
\begin{column}{0.49\textwidth}
Ipotesi iniziali:
\[
T=0,\qquad m_u=m_d=m_s=0.
\]
Specie: tre sapori $u,d,s$; elettroni solo se richiesti dalla neutralit\`a.

Obiettivo: determinare
\[
P,\quad \varepsilon,\quad n_B,
\]
in funzione dei potenziali chimici, imponendo equilibrio debole e neutralit\`a elettrica.
\end{column}
\end{columns}
\medskip
Le formule seguenti usano $\hbar=c=1$, salvo quando si ripristinano le unit\`a.
\end{frame}

\begin{frame}[t]{Cariche conservate e potenziali chimici}
\small
\begin{columns}[T,onlytextwidth]
\begin{column}{0.49\textwidth}
Se neutrini e antineutrini sfuggono, la materia in equilibrio debole conserva globalmente:
\begin{itemize}
\item numero barionico $B$;
\item carica elettrica $Q$.
\end{itemize}
Per una specie $i$:
\[
\mu_i=B_i\mu_B+Q_i\mu_Q.
\]
$\mu_B$ e $\mu_Q$: potenziali chimici delle due cariche conservate.
\end{column}
\begin{column}{0.49\textwidth}
\centering
\begin{tabular}{c|cc}
Specie $i$ & $B_i$ & $Q_i$ \\ \hline
$u$ & $1/3$ & $+2/3$ \\
$d$ & $1/3$ & $-1/3$ \\
$s$ & $1/3$ & $-1/3$ \\
$e^-$ & $0$ & $-1$
\end{tabular}
\par\medskip
Ciascun quark porta $1/3$ di numero barionico.

La stranezza non \`e una carica indipendentemente conservata quando sono in equilibrio le reazioni deboli.
\end{column}
\end{columns}
\end{frame}

\begin{frame}[t]{Espressioni dei potenziali chimici}
\small
\begin{columns}[T,onlytextwidth]
\begin{column}{0.49\textwidth}
Applicando $\mu_i=B_i\mu_B+Q_i\mu_Q$:
\[
\begin{aligned}
\mu_u&=\tfrac13\mu_B+\tfrac23\mu_Q,\\
\mu_d&=\tfrac13\mu_B-\tfrac13\mu_Q,\\
\mu_s&=\tfrac13\mu_B-\tfrac13\mu_Q,\\
\mu_e&=-\mu_Q.
\end{aligned}
\]
Il libro indica $\mu_Q$ anche con $\mu_Z$.
\end{column}
\begin{column}{0.49\textwidth}
Introduciamo
\[
\mu_q\equiv\frac{\mu_B}{3},\qquad
\mu_Q=-\mu_e.
\]
Allora
\[
\begin{aligned}
\mu_u&=\mu_q-\tfrac23\mu_e,\\
\mu_d&=\mu_q+\tfrac13\mu_e,\\
\mu_s&=\mu_q+\tfrac13\mu_e.
\end{aligned}
\]
$\mu_q$ \`e un parametro comune, non necessariamente il potenziale chimico di ciascun sapore.
\end{column}
\end{columns}
\end{frame}

\begin{frame}[t]{Equilibrio debole: le condizioni sono automatiche}
\small
\begin{columns}[T,onlytextwidth]
\begin{column}{0.49\textwidth}
Per il processo
\[
u+e^-\rightleftharpoons d+\nu_e
\]
l'equilibrio richiede
\[
\mu_u+\mu_e=\mu_d+\mu_{\nu_e}.
\]
Se i neutrini non sono intrappolati:
\[
\mu_{\nu_e}=0
\quad\Rightarrow\quad
\mu_d=\mu_u+\mu_e.
\]
Processi deboli con stranezza impongono
\[
\mu_s=\mu_d.
\]
\end{column}
\begin{column}{0.49\textwidth}
Verifichiamo le condizioni mediante le cariche conservate:
\[
\begin{aligned}
\mu_u+\mu_e
&=\left(\mu_q-\tfrac23\mu_e\right)+\mu_e\\
&=\mu_q+\tfrac13\mu_e\\
&=\mu_d=\mu_s.
\end{aligned}
\]
L'uguaglianza $\mu_d=\mu_s$ deriva dall'identit\`a delle cariche $B_i,Q_i$, \emph{non} da uguaglianza di massa o densit\`a.
\end{column}
\end{columns}
\end{frame}

\subsection{Gas relativistico degenere: pressione e densit\`a}

\begin{frame}[t]{Superficie di Fermi e degenerazione dei quark}
\small
\begin{columns}[T,onlytextwidth]
\begin{column}{0.49\textwidth}
Per un fermione libero:
\[
E_i(k)=\sqrt{k^2+m_i^2}.
\]
La distribuzione di Fermi--Dirac:
\[
f_i(k,T)=\frac{1}{e^{(E_i(k)-\mu_i)/T}+1}.
\]
A $T=0$ si riduce a
\[
f_i(k,0)=\Theta(k_{F,i}-k).
\]
$\Theta$ \`e la funzione gradino; $k_{F,i}$ \`e il momento di Fermi.
\end{column}
\begin{column}{0.49\textwidth}
La degenerazione per \emph{ciascun} sapore di quark \`e
\[
\gamma_i=(2s+1)N_c=2\cdot3=6,
\]
con spin $s=1/2$ e numero di colori $N_c=3$.

Per $m_i=0$:
\[
E_i(k)=k,\qquad \mu_i=k_{F,i}.
\]
Gli antiquark, a $T=0$ e $\mu_i>0$, non sono occupati nel gas considerato: non si aggiunge un fattore 2 a $\gamma_i$.
\end{column}
\end{columns}
\end{frame}

\begin{frame}[t]{Densit\`a numerica: integrazione nella sfera di Fermi}
\small
\begin{columns}[T,onlytextwidth]
\begin{column}{0.49\textwidth}
Contiamo gli stati occupati per unit\`a di volume:
\[
 n_i=\gamma_i\int_{|\mathbf k|<k_{F,i}}
\frac{d^3k}{(2\pi)^3}.
\]
Per isotropia, $d^3k=4\pi k^2\,dk$:
\[
\begin{aligned}
n_i&=\frac{\gamma_i}{2\pi^2}
\int_0^{k_{F,i}}k^2\,dk\\
&=\frac{\gamma_i}{2\pi^2}
\left[\frac{k^3}{3}\right]_0^{k_{F,i}}.
\end{aligned}
\]
\end{column}
\begin{column}{0.49\textwidth}
Dunque
\[
 n_i=\frac{\gamma_i k_{F,i}^3}{6\pi^2}.
\]
Con $\gamma_i=6$:
\[
 n_i=\frac{k_{F,i}^3}{\pi^2}.
\]
Per quark senza massa ($k_{F,i}=\mu_i$):
\[
\boxed{n_i=\frac{\mu_i^3}{\pi^2}}.
\]
Il colore aumenta la densit\`a di stati rispetto a un fermione con sola degenerazione di spin.
\end{column}
\end{columns}
\end{frame}

\begin{frame}[t]{Pressione: dall'integrale cinetico a $\mu_i^4$}
\small
\begin{columns}[T,onlytextwidth]
\begin{column}{0.49\textwidth}
Per un gas isotropo:
\[
P_i=\frac{\gamma_i}{3}\int\frac{d^3k}{(2\pi)^3}
\frac{k^2}{E_i(k)}f_i(k).
\]
A $T=0$, per particelle massless $E_i(k)=k$:
\[
\begin{aligned}
P_i&=\frac{\gamma_i}{6\pi^2}
\int_0^{k_{F,i}}\frac{k^4}{k}\,dk\\
&=\frac{\gamma_i}{6\pi^2}
\int_0^{k_{F,i}}k^3\,dk.
\end{aligned}
\]
\end{column}
\begin{column}{0.49\textwidth}
Integrando:
\[
\begin{aligned}
P_i&=\frac{\gamma_i}{6\pi^2}
\left[\frac{k^4}{4}\right]_0^{k_{F,i}}\\
&=\frac{\gamma_i k_{F,i}^4}{24\pi^2}.
\end{aligned}
\]
Per $\gamma_i=6$ e $\mu_i=k_{F,i}$:
\[
\boxed{P_i=\frac{\mu_i^4}{4\pi^2}}.
\]
Pressione complessiva dei tre sapori:
\[
P=\frac{\mu_u^4+\mu_d^4+\mu_s^4}{4\pi^2}.
\]
\end{column}
\end{columns}
\end{frame}

\begin{frame}[t]{Controllo termodinamico: $n_i=\partial P/\partial\mu_i$}
\small
\begin{columns}[T,onlytextwidth]
\begin{column}{0.49\textwidth}
Nel formalismo gran-canonico:
\[
\Omega=-PV.
\]
Per variabili $\mu_i$ formalmente indipendenti:
\[
n_i=\left(\frac{\partial P}{\partial\mu_i}\right)_{T,\mu_{j\ne i}}.
\]
Da
\[
P_i=\frac{\mu_i^4}{4\pi^2}
\]
segue
\[
\frac{\partial P_i}{\partial\mu_i}
=\frac{4\mu_i^3}{4\pi^2}.
\]
\end{column}
\begin{column}{0.49\textwidth}
Otteniamo
\[
\boxed{n_i=\frac{\mu_i^3}{\pi^2}},
\]
come gi\`a ricavato mediante il volume della sfera di Fermi.

Le $\mu_i$ possono successivamente essere espresse mediante $\mu_B$ e $\mu_Q$: queste ultime sono le variabili chimiche indipendenti dopo aver imposto l'equilibrio debole.
\end{column}
\end{columns}
\end{frame}

\begin{frame}[t]{Densit\`a di energia e relazione $\varepsilon=3P$}
\small
\begin{columns}[T,onlytextwidth]
\begin{column}{0.49\textwidth}
A temperatura nulla la relazione termodinamica \`e
\[
\varepsilon=-P+\sum_i\mu_i n_i.
\]
Per una specie massless:
\[
\begin{aligned}
\mu_i n_i&=\mu_i\frac{\mu_i^3}{\pi^2}\\
&=\frac{\mu_i^4}{\pi^2}=4P_i.
\end{aligned}
\]
Quindi
\[
\varepsilon_i=-P_i+4P_i=3P_i.
\]
\end{column}
\begin{column}{0.49\textwidth}
Sommando i tre sapori:
\[
\boxed{\varepsilon=3P},
\qquad
\boxed{P=\frac\varepsilon3}.
\]
Lo stesso risultato segue direttamente da
\[
\varepsilon_i=\frac{\gamma_i}{2\pi^2}
\int_0^{k_{F,i}}k^2E_i(k)\,dk,
\]
che per $E_i=k$ vale $\gamma_i k_{F,i}^4/(8\pi^2)$.
\medskip
Nel gas ideale massless:
\[
c_s^2=\frac{dP}{d\varepsilon}=\frac13.
\]
\end{column}
\end{columns}
\end{frame}

\subsection{Neutralit\`a e soluzione a tre sapori massless}

\begin{frame}[t]{Neutralit\`a elettrica: quark e leptoni}
\small
\begin{columns}[T,onlytextwidth]
\begin{column}{0.49\textwidth}
La densit\`a di carica elettrica (in unit\`a di $e$) \`e
\[
n_Q=\sum_iQ_in_i.
\]
Per $u,d,s,e^-,\mu^-$:
\[
\begin{aligned}
n_Q={}&\tfrac23n_u-\tfrac13n_d-\tfrac13n_s\\
&-n_e-n_\mu.
\end{aligned}
\]
La neutralit\`a locale impone
\[
\boxed{n_Q=0}.
\]
\end{column}
\begin{column}{0.49\textwidth}
Per la prima soluzione, poniamo $n_e=n_\mu=0$:
\[
2n_u-n_d-n_s=0.
\]
Una soluzione immediata \`e
\[
 n_u=n_d=n_s\equiv n.
\]
Infatti
\[
\left(\tfrac23-\tfrac13-\tfrac13\right)n=0.
\]
La densit\`a di carica si annulla grazie alle popolazioni dei tre sapori.
\end{column}
\end{columns}
\end{frame}

\begin{frame}[t]{Densit\`a uguali implicano $\mu_e=0$ nel caso massless}
\small
\begin{columns}[T,onlytextwidth]
\begin{column}{0.49\textwidth}
Per quark liberi senza massa:
\[
n_i=\frac{\mu_i^3}{\pi^2}.
\]
Se $n_u=n_d=n_s$ e $\mu_i>0$, allora
\[
\mu_u=\mu_d=\mu_s.
\]
L'equilibrio debole richiede
\[
\mu_d=\mu_u+\mu_e.
\]
Dunque
\[
\boxed{\mu_e=0}.
\]
\end{column}
\begin{column}{0.49\textwidth}
Poich\'e $\mu_Q=-\mu_e$, anche
\[
\mu_Q=0.
\]
Le formule dei potenziali chimici danno
\[
\boxed{\mu_u=\mu_d=\mu_s=\mu_q},
\qquad \mu_q=\frac{\mu_B}{3}.
\]
A $T=0$, $\mu_e=0$ non produce una popolazione elettronica.

La soluzione \`e quindi sia neutra sia in equilibrio debole.
\end{column}
\end{columns}
\end{frame}

\begin{frame}[t]{Densit\`a barionica dei tre sapori}
\small
\begin{columns}[T,onlytextwidth]
\begin{column}{0.49\textwidth}
Il numero barionico di ciascun quark \`e $1/3$:
\[
n_B=\frac13(n_u+n_d+n_s).
\]
Nella soluzione simmetrica:
\[
n_u=n_d=n_s=\frac{\mu_q^3}{\pi^2}.
\]
Perci\`o
\[
\begin{aligned}
n_B&=\frac13\left(3\frac{\mu_q^3}{\pi^2}\right)\\
&=\frac{\mu_q^3}{\pi^2}.
\end{aligned}
\]
\end{column}
\begin{column}{0.49\textwidth}
La conversione da unit\`a naturali a $\mathrm{fm^{-3}}$ introduce $(\hbar c)^{-3}$:
\[
n_B=\frac{\mu_q^3}{\pi^2(\hbar c)^3},
\]
con $\hbar c=197.327\,\mathrm{MeV\,fm}$.

Quindi
\[
\boxed{
_B\simeq0.36\,\mathrm{fm^{-3}}
\left(\frac{\mu_q}{300\,\mathrm{MeV}}\right)^3}.
\]
\end{column}
\end{columns}
\end{frame}

\begin{frame}[t]{Pressione ed energia per tre sapori massless}
\small
\begin{columns}[T,onlytextwidth]
\begin{column}{0.49\textwidth}
La pressione di ciascun sapore \`e
\[
P_i=\frac{\mu_i^4}{4\pi^2}.
\]
Poich\'e $\mu_i=\mu_q$:
\[
P=\sum_{i=u,d,s}P_i
=\frac{3\mu_q^4}{4\pi^2}.
\]
Dal risultato $\varepsilon=3P$:
\[
\boxed{\varepsilon=\frac{9\mu_q^4}{4\pi^2}}.
\]
\end{column}
\begin{column}{0.49\textwidth}
La densit\`a di energia fisica richiede ancora $(\hbar c)^{-3}$:
\[
\varepsilon_{\rm fis}=
\frac{9\mu_q^4}{4\pi^2(\hbar c)^3}.
\]
Numericamente:
\[
\boxed{
\varepsilon_{\rm fis}\simeq
240\,\mathrm{MeV\,fm^{-3}}
\left(\frac{\mu_q}{300\,\mathrm{MeV}}\right)^4}.
\]
Si ha $n_B\propto\mu_q^3$ ed $\varepsilon\propto\mu_q^4$.
\end{column}
\end{columns}
\end{frame}

\begin{frame}[t]{Perch\'e $\varepsilon\propto n_B^{4/3}$?}
\small
\begin{columns}[T,onlytextwidth]
\begin{column}{0.49\textwidth}
Dalla densit\`a barionica:
\[
n_B=\frac{\mu_q^3}{\pi^2}
\quad\Rightarrow\quad
\mu_q=(\pi^2n_B)^{1/3}.
\]
Sostituendo nell'energia:
\[
\begin{aligned}
\varepsilon
&=\frac{9}{4\pi^2}\mu_q^4\\
&=\frac{9}{4\pi^2}(\pi^2n_B)^{4/3}.
\end{aligned}
\]
\end{column}
\begin{column}{0.49\textwidth}
Semplificando:
\[
\boxed{\varepsilon=\frac94\pi^{2/3}n_B^{4/3}}
\]
(in unit\`a naturali).

Quindi
\[
\boxed{\varepsilon\propto n_B^{4/3}}.
\]
L'energia cresce pi\`u rapidamente del numero barionico, perch\'e aumentano sia gli stati occupati sia il momento medio dei quark.
\end{column}
\end{columns}
\end{frame}

\begin{frame}[t]{Numeri di riferimento: densit\`a delle stelle compatte}
\small
\begin{columns}[T,onlytextwidth]
\begin{column}{0.49\textwidth}
\centering
\begin{tabular}{r|r|r}
$\mu_q$ & $n_B$ & $\varepsilon$ \\
(MeV) & ($\mathrm{fm^{-3}}$) & ($\mathrm{MeV/fm^3}$)\\ \hline
300 & 0.36 & 240 \\
400 & 0.84 & 760 \\
500 & 1.6 & 1900
\end{tabular}
\par\medskip
Il potenziale chimico barionico \`e
\[
\mu_B=3\mu_q
=900,1200,1500\ \mathrm{MeV}.
\]
\end{column}
\begin{column}{0.49\textwidth}
Per $n_0\simeq0.15\,\mathrm{fm^{-3}}$:
\[
\frac{n_B}{n_0}\simeq2.4,\ 5.6,\ 11.
\]
Stima della densit\`a media di energia stellare, in unit\`a naturali:
\[
\bar\varepsilon\simeq\frac{3M}{4\pi R^3}.
\]
Per $M=2M_\odot$, $R=10\,\mathrm{km}$:
\[
\bar\varepsilon\simeq270\,\mathrm{MeV\,fm^{-3}}.
\]
Il confronto riguarda le scale, non dimostra la presenza di quark deconfinati.
\end{column}
\end{columns}
\end{frame}

\subsection{Massa dello strange e neutralit\`a}

\begin{frame}[t]{La massa del quark strange modifica la composizione}
\small
\begin{columns}[T,onlytextwidth]
\begin{column}{0.49\textwidth}
Per $u,d$ si usa ancora
\[
m_u,m_d\ll\mu_q,
\quad E_{u,d}(k)\simeq k.
\]
Per il quark strange, $m_s\sim95\,\mathrm{MeV}$:
\[
E_s(k)=\sqrt{k^2+m_s^2}.
\]
A $T=0$:
\[
\mu_s=\sqrt{k_{F,s}^2+m_s^2}.
\]
\end{column}
\begin{column}{0.49\textwidth}
Risolvendo rispetto al momento di Fermi:
\[
k_{F,s}=\sqrt{\mu_s^2-m_s^2}
\qquad(\mu_s>m_s).
\]
Dunque
\[
\boxed{n_s=\frac{(\mu_s^2-m_s^2)^{3/2}}{\pi^2}}.
\]
Mentre
\[
n_d=\frac{\mu_d^3}{\pi^2},
\quad \mu_s=\mu_d.
\]
A uguale potenziale chimico, $n_s<n_d$.
\end{column}
\end{columns}
\end{frame}

\begin{frame}[t]{Neutralit\`a con $m_s>0$ ed elettroni}
\small
\begin{columns}[T,onlytextwidth]
\begin{column}{0.49\textwidth}
La riduzione di $n_s$ altera la compensazione fra cariche positive e negative.

Senza muoni, la neutralit\`a \`e
\[
\tfrac23n_u-\tfrac13n_d-\tfrac13n_s-n_e=0,
\]
cio\`e
\[
\boxed{n_e=\tfrac23n_u-\tfrac13n_d-\tfrac13n_s}.
\]
\end{column}
\begin{column}{0.49\textwidth}
Usiamo
\[
\begin{aligned}
\mu_u&=\mu_q-\tfrac23\mu_e,\\
\mu_d&=\mu_s=\mu_q+\tfrac13\mu_e.
\end{aligned}
\]
Ci proponiamo di trovare $\mu_e$ supponendo
\[
\mu_q\gg m_s\gg\mu_e\gg m_e.
\]
Servir\`a uno sviluppo in $m_s/\mu_q$ e $\mu_e/\mu_q$.
\end{column}
\end{columns}
\end{frame}

\begin{frame}[t]{Sviluppo di $n_u$ al primo ordine in $\mu_e$}
\small
\begin{columns}[T,onlytextwidth]
\begin{column}{0.49\textwidth}
La densit\`a massless del quark up \`e
\[
n_u=\frac{\mu_u^3}{\pi^2}.
\]
Poich\'e $\mu_u=\mu_q-\frac23\mu_e$:
\[
 n_u=\frac1{\pi^2}
\left(\mu_q-\frac23\mu_e\right)^3.
\]
Poniamo $a=\mu_q$, $b=2\mu_e/3$ e usiamo
\[
(a-b)^3=a^3-3a^2b+O(b^2).
\]
\end{column}
\begin{column}{0.49\textwidth}
Il termine lineare \`e
\[
3\mu_q^2\frac{2\mu_e}{3}
=2\mu_q^2\mu_e.
\]
Pertanto
\[
\boxed{
_u\simeq\frac1{\pi^2}
\left(\mu_q^3-2\mu_q^2\mu_e\right)}.
\]
Gli errori trascurati sono di ordine $\mu_q\mu_e^2$ o superiore.
\end{column}
\end{columns}
\end{frame}

\begin{frame}[t]{Sviluppo di $n_d$ al primo ordine in $\mu_e$}
\small
\begin{columns}[T,onlytextwidth]
\begin{column}{0.49\textwidth}
Per il quark down:
\[
n_d=\frac{\mu_d^3}{\pi^2}
=\frac1{\pi^2}
\left(\mu_q+\frac13\mu_e\right)^3.
\]
Usiamo
\[
(a+b)^3=a^3+3a^2b+O(b^2).
\]
\end{column}
\begin{column}{0.49\textwidth}
Con $a=\mu_q$, $b=\mu_e/3$:
\[
3a^2b=\mu_q^2\mu_e.
\]
Segue
\[
\boxed{
n_d\simeq\frac1{\pi^2}
\left(\mu_q^3+\mu_q^2\mu_e\right)}.
\]
La variazione di $n_d$ \`e opposta a quella di $n_u$: $\mu_d$ aumenta quando aumenta $\mu_e$.
\end{column}
\end{columns}
\end{frame}

\begin{frame}[t]{Sviluppo della densit\`a strange: primo passaggio}
\small
\begin{columns}[T,onlytextwidth]
\begin{column}{0.49\textwidth}
Per lo strange massivo:
\[
n_s=\frac1{\pi^2}
(\mu_s^2-m_s^2)^{3/2},
\]
con $\mu_s=\mu_q+\frac13\mu_e$.

Espandiamo il quadrato:
\[
\mu_s^2=\mu_q^2+
\frac23\mu_q\mu_e+
\frac19\mu_e^2.
\]
Al primo ordine in $\mu_e$:
\[
\mu_s^2-m_s^2\simeq
\mu_q^2+\tfrac23\mu_q\mu_e-m_s^2.
\]
\end{column}
\begin{column}{0.49\textwidth}
Raccogliamo $\mu_q^2$:
\[
\begin{aligned}
\mu_s^2-m_s^2\simeq{}&\mu_q^2
\left(1+\frac{2\mu_e}{3\mu_q}
-\frac{m_s^2}{\mu_q^2}\right).
\end{aligned}
\]
Quindi
\[
 n_s\simeq\frac{\mu_q^3}{\pi^2}
\left(1+\frac{2\mu_e}{3\mu_q}
-\frac{m_s^2}{\mu_q^2}\right)^{3/2}.
\]
\end{column}
\end{columns}
\end{frame}

\begin{frame}[t]{Sviluppo della densit\`a strange: secondo passaggio}
\small
\begin{columns}[T,onlytextwidth]
\begin{column}{0.49\textwidth}
Per $|x|\ll1$:
\[
(1+x)^{3/2}=1+\frac32x+O(x^2).
\]
Poniamo
\[
x=\frac{2\mu_e}{3\mu_q}-\frac{m_s^2}{\mu_q^2}.
\]
Allora
\[
\begin{aligned}
(1+x)^{3/2}\simeq 1
&+\frac{\mu_e}{\mu_q}\\
&-\frac32\frac{m_s^2}{\mu_q^2}.
\end{aligned}
\]
\end{column}
\begin{column}{0.49\textwidth}
Moltiplicando per $\mu_q^3/\pi^2$:
\[
\boxed{
_s\simeq\frac1{\pi^2}
\left(\mu_q^3+\mu_q^2\mu_e
-\frac32\mu_qm_s^2\right)}.
\]
I termini omessi includono $\mu_e^2$, $m_s^4$ e prodotti di ordine superiore.

La massa strange riduce $n_s$ di un termine proporzionale a $m_s^2$.
\end{column}
\end{columns}
\end{frame}

\begin{frame}[t]{Sostituzione nell'equazione di neutralit\`a}
\footnotesize
Partiamo da
\[
0=\frac23n_u-\frac13n_d-\frac13n_s-n_e.
\]
\begin{columns}[T,onlytextwidth]
\begin{column}{0.49\textwidth}
Inseriamo i tre sviluppi:
\[
\begin{aligned}
\pi^2n_u&\simeq\mu_q^3-2\mu_q^2\mu_e,\\
\pi^2n_d&\simeq\mu_q^3+\mu_q^2\mu_e,\\
\pi^2n_s&\simeq\mu_q^3+\mu_q^2\mu_e
-\tfrac32\mu_qm_s^2.
\end{aligned}
\]
\end{column}
\begin{column}{0.49\textwidth}
Perci\`o, al primo ordine,
\[
\begin{aligned}
0\simeq {}&\frac1{\pi^2}
\left[\left(\tfrac23-\tfrac13-\tfrac13\right)\mu_q^3\right.\\
&\left.\quad-\left(\tfrac43+\tfrac13+\tfrac13\right)
\mu_q^2\mu_e
+\tfrac12\mu_qm_s^2\right]-n_e.
\end{aligned}
\]
I termini $\mu_q^3$ si cancellano.
\end{column}
\end{columns}
\end{frame}

\begin{frame}[t]{Neutralit\`a e potenziale chimico elettronico}
\small
\begin{columns}[T,onlytextwidth]
\begin{column}{0.49\textwidth}
Dopo la cancellazione:
\[
0\simeq
-\frac{2\mu_q^2\mu_e}{\pi^2}
+\frac{\mu_qm_s^2}{2\pi^2}
-n_e.
\]
Per elettroni ultrarelativistici, con degenerazione 2:
\[
n_e=\frac{\mu_e^3}{3\pi^2}.
\]
Poich\'e $\mu_e/\mu_q\ll1$, questo termine \`e di ordine cubico e non contribuisce all'ordine considerato.
\end{column}
\begin{column}{0.49\textwidth}
Trascurando $n_e$ soltanto \emph{nell'equazione al primo ordine}:
\[
2\mu_q^2\mu_e\simeq\frac12\mu_qm_s^2.
\]
Dividiamo per $2\mu_q^2$:
\[
\boxed{\mu_e\simeq\frac{m_s^2}{4\mu_q}}.
\]
Il potenziale chimico elettronico \`e positivo, ma piccolo rispetto a $\mu_q$.
\end{column}
\end{columns}
\end{frame}

\begin{frame}[t]{Verifica numerica: $\mu_e$ e gerarchia delle scale}
\small
\begin{columns}[T,onlytextwidth]
\begin{column}{0.49\textwidth}
Assumiamo $m_s=95\,\mathrm{MeV}$:
\[
\mu_e\simeq\frac{95^2}{4\mu_q}
\;\mathrm{MeV}.
\]
\centering
\begin{tabular}{r|r}
$\mu_q$ (MeV) & $\mu_e$ (MeV)\\ \hline
300 & 7.52\\
400 & 5.64\\
500 & 4.51
\end{tabular}
\end{column}
\begin{column}{0.49\textwidth}
La gerarchia \`e coerente:
\[
\mu_e\ll m_s\ll\mu_q.
\]
I valori di $\mu_e$ sono dell'ordine di $5$--$8\,\mathrm{MeV}$.

L'effetto della massa strange sulla composizione dei quark \`e presente gi\`a all'ordine $m_s^2/\mu_q^2$; la densit\`a elettronica resta molto pi\`u piccola.
\end{column}
\end{columns}
\end{frame}

\begin{frame}[t]{Frazione elettronica: derivazione del sesto ordine}
\small
\begin{columns}[T,onlytextwidth]
\begin{column}{0.49\textwidth}
Per elettroni ultrarelativistici:
\[
n_e\simeq\frac{\mu_e^3}{3\pi^2}.
\]
Al primo ordine, per i quark:
\[
n_B\simeq\frac{\mu_q^3}{\pi^2}.
\]
Quindi
\[
\begin{aligned}
\frac{n_e}{n_B}
&\simeq\frac{\mu_e^3/(3\pi^2)}
{\mu_q^3/\pi^2}\\
&=\frac13\left(\frac{\mu_e}{\mu_q}\right)^3.
\end{aligned}
\]
\end{column}
\begin{column}{0.49\textwidth}
Sostituiamo $\mu_e\simeq m_s^2/(4\mu_q)$:
\[
\begin{aligned}
\frac{n_e}{n_B}
&\simeq\frac13
\left(\frac{m_s^2}{4\mu_q^2}\right)^3\\
&=\frac{m_s^6}{192\mu_q^6}.
\end{aligned}
\]
Equivalentemente:
\[
\boxed{\frac{n_e}{n_B}\simeq
\frac13\left(\frac{m_s}{2\mu_q}\right)^6}.
\]
La frazione elettronica \`e soppressa alla sesta potenza.
\end{column}
\end{columns}
\end{frame}

\begin{frame}[t]{Elettroni rarissimi e assenza di muoni}
\small
\begin{columns}[T,onlytextwidth]
\begin{column}{0.49\textwidth}
Per $m_s=95\,\mathrm{MeV}$:
\begin{center}
\begin{tabular}{r|r}
$\mu_q$ (MeV) & $n_e/n_B$\\ \hline
300 & $5.2\times10^{-6}$\\
400 & $9.3\times10^{-7}$\\
500 & $2.5\times10^{-7}$
\end{tabular}
\end{center}
L'elettrone compensa una piccola carica residua, ma contribuisce pochissimo alla densit\`a barionica (alla quale comunque non partecipa).
\end{column}
\begin{column}{0.49\textwidth}
Se l'equilibrio fra leptoni impone
\[
\mu_\mu=\mu_e,
\]
i muoni liberi possono comparire a $T=0$ solo per
\[
\mu_\mu\ge m_\mu\simeq105.7\,\mathrm{MeV}.
\]
Poich\'e $\mu_e\sim5$--$8\,\mathrm{MeV}$:
\[
\boxed{n_\mu=0}.
\]
La quasi assenza di leptoni \`e una propriet\`a della materia $uds$ quasi neutra.
\end{column}
\end{columns}
\end{frame}

\subsection{Termodinamica e collegamento con le stelle compatte}

\begin{frame}[t]{Identit\`a termodinamica in materia neutra}
\small
\begin{columns}[T,onlytextwidth]
\begin{column}{0.49\textwidth}
A $T=0$:
\[
\varepsilon+P=\sum_i\mu_i n_i.
\]
Poich\'e
\[
\mu_i=B_i\mu_B+Q_i\mu_Q,
\]
segue
\[
\begin{aligned}
\sum_i\mu_i n_i
&=\mu_B\sum_iB_in_i\\
&\quad+\mu_Q\sum_iQ_in_i.
\end{aligned}
\]
\end{column}
\begin{column}{0.49\textwidth}
Definiamo
\[
n_B=\sum_iB_in_i,\qquad
n_Q=\sum_iQ_in_i.
\]
Dunque
\[
\varepsilon+P=\mu_B n_B+\mu_Q n_Q.
\]
Per materia elettricamente neutra ($n_Q=0$):
\[
\boxed{\varepsilon+P=\mu_Bn_B},
\quad
\boxed{\mu_B=\frac{\varepsilon+P}{n_B}}.
\]
\end{column}
\end{columns}
\end{frame}

\begin{frame}[t]{Controllo nel gas $uds$ massless}
\small
\begin{columns}[T,onlytextwidth]
\begin{column}{0.49\textwidth}
Nella soluzione simmetrica:
\[
P=\frac{3\mu_q^4}{4\pi^2},
\quad
\varepsilon=3P,
\quad
n_B=\frac{\mu_q^3}{\pi^2}.
\]
La relazione generale fornisce
\[
\mu_B=\frac{\varepsilon+P}{n_B}
=\frac{4P}{n_B}.
\]
\end{column}
\begin{column}{0.49\textwidth}
Sostituendo $P$ e $n_B$:
\[
\begin{aligned}
\mu_B
&=\frac{4(3\mu_q^4/4\pi^2)}
{\mu_q^3/\pi^2}\\
&=3\mu_q.
\end{aligned}
\]
\[
\boxed{\mu_B=3\mu_q}.
\]
La relazione termodinamica conferma la definizione chimica iniziale.
\end{column}
\end{columns}
\end{frame}

\begin{frame}[t]{Dal gas libero al MIT bag model}
\small
\begin{columns}[T,onlytextwidth]
\begin{column}{0.49\textwidth}
Finora:
\[
P_{\rm free}=\frac{3\mu_q^4}{4\pi^2},
\qquad\varepsilon_{\rm free}=3P_{\rm free}.
\]
Per modellizzare un contributo costante di energia del vuoto, introduciamo il parametro $B>0$:
\[
\begin{aligned}
P&=P_{\rm free}-B,\\
\varepsilon&=\varepsilon_{\rm free}+B.
\end{aligned}
\]
$B$ ha le dimensioni di densit\`a di energia.
\end{column}
\begin{column}{0.49\textwidth}
Dalla seconda equazione:
\[
\varepsilon-B=3P_{\rm free}.
\]
E dalla prima:
\[
P_{\rm free}=P+B.
\]
Quindi
\[
\varepsilon-B=3(P+B)
\]
ossia
\[
\boxed{P=\frac{\varepsilon-4B}{3}}.
\]
Questo vale nella versione massless con $B$ costante.
\end{column}
\end{columns}
\end{frame}

\begin{frame}[t]{Superficie a pressione nulla e autolegamento}
\small
\begin{columns}[T,onlytextwidth]
\begin{column}{0.49\textwidth}
Nel bag model massless:
\[
P=\frac{\varepsilon-4B}{3}.
\]
Alla superficie di una configurazione autolegata:
\[
P=0\quad\Longrightarrow\quad
\boxed{\varepsilon_{\rm surf}=4B}.
\]
La pressione si annulla quindi a una densit\`a di energia finita.
\end{column}
\begin{column}{0.49\textwidth}
Questo \`e il meccanismo schematico per una stella di quark \emph{self-bound}.

Non basta, per\`o, a dimostrare che la materia strange sia lo stato fondamentale della materia: occorre confrontare l'energia per barione a pressione nulla con quella della materia adronica.

Le interazioni forti e la massa $m_s$ introducono ulteriori correzioni alla EOS.
\end{column}
\end{columns}
\end{frame}

\begin{frame}[t]{Che cosa resta valido quando $m_s\ne0$?}
\small
\begin{columns}[T,onlytextwidth]
\begin{column}{0.49\textwidth}
Con $m_s>0$:
\[
\mu_d=\mu_s=\mu_u+\mu_e,
\]
ma $n_s\ne n_d$ perch\'e
\[
k_{F,s}=\sqrt{\mu_s^2-m_s^2}.
\]
La neutralit\`a richiede una piccola popolazione elettronica:
\[
\mu_e\simeq\frac{m_s^2}{4\mu_q}.
\]
\end{column}
\begin{column}{0.49\textwidth}
Invece la semplice relazione
\[
P=\frac\varepsilon3
\]
non \`e esatta per particelle massive.

La pressione del sapore strange \`e data da
\[
P_s=\frac1{\pi^2}\int_0^{k_{F,s}}
\frac{k^4}{\sqrt{k^2+m_s^2}}\,dk.
\]
La piccola frazione elettronica non implica quindi correzioni altrettanto piccole alla EOS.
\end{column}
\end{columns}
\end{frame}

\begin{frame}[t]{Sintesi e limiti dell'approssimazione}
\small
\begin{columns}[T,onlytextwidth]
\begin{column}{0.49\textwidth}
Tre sapori liberi e massless, $T=0$:
\[
\begin{gathered}
\mu_u=\mu_d=\mu_s=\mu_q,\\
n_u=n_d=n_s,\qquad n_e=0,\\
n_B=\mu_q^3/\pi^2,\\
P=3\mu_q^4/(4\pi^2),\\
\varepsilon=3P.
\end{gathered}
\]
\end{column}
\begin{column}{0.49\textwidth}
Per strange massivo, al primo ordine:
\[
\begin{gathered}
\mu_e\simeq m_s^2/(4\mu_q),\\
\displaystyle\frac{n_e}{n_B}\simeq
\frac13\left(\frac{m_s}{2\mu_q}\right)^6.
\end{gathered}
\]
Il gas libero \`e un riferimento asintotico; per stelle compatte realistiche servono contributi dell'interazione forte e, eventualmente, del vuoto.
\end{column}
\end{columns}
\medskip
Nel libro il passo successivo \`e capire se una EOS della materia di quark pu\`o descrivere stelle compatte autolegate.
\end{frame}

\section{Stelle autolegate: materia strange}
\subsection{Equazione di stato del MIT bag model}

% 01
\begin{frame}{Dal gas di quark libero alle stelle autolegate}
\small
\begin{columns}[T,onlytextwidth]
\begin{column}{0.49\textwidth}
Nel \S8.2: materia fredda di quark $u,d,s$,
assunti inizialmente massless e non interagenti.
\[
P_q=\frac{\varepsilon_q}{3}.
\]
$P_q$ ed $\varepsilon_q$ sono pressione e densita'
di energia del solo gas di quark.
\medskip
Nel MIT bag model si aggiunge l'energia
di vuoto $B>0$, supposta costante.
\end{column}
\begin{column}{0.49\textwidth}
La pressione del vuoto e' negativa:
\[
\varepsilon_{\rm vac}=B,\qquad
P_{\rm vac}=-B.
\]
Obiettivi del \S8.3:
\begin{itemize}
\item ricavare l'EOS $P(\varepsilon)$;
\item distinguere autolegamento e stabilita' assoluta;
\item studiare $M(R)$ e i vincoli su $B$.
\end{itemize}
\end{column}
\end{columns}
\medskip
Useremo $\hbar=c=1$; nei frame sulle TOV porremo anche $G=1$.
\end{frame}

% 02
\begin{frame}{MIT bag model: derivazione dell'EOS (8.35)--(8.37)}
\small
\begin{columns}[T,onlytextwidth]
\begin{column}{0.49\textwidth}
La pressione totale somma i contributi:
\[
P=P_q+P_{\rm vac}=P_q-B.
\]
La densita' di energia totale e'
\[
\varepsilon=\varepsilon_q+B=3P_q+B.
\]
Isoliamo la pressione dei quark:
\[
P_q=\frac{\varepsilon-B}{3}.
\]
\end{column}
\begin{column}{0.49\textwidth}
Sostituiamo nella pressione totale:
\begin{align*}
P&=\frac{\varepsilon-B}{3}-B\\
 &=\frac{\varepsilon-B-3B}{3}\\
 &=\frac{\varepsilon-4B}{3}.
\end{align*}
Quindi
\[
P(\varepsilon)=\frac13(\varepsilon-4B).
\]
\end{column}
\end{columns}
\end{frame}

% 03
\begin{frame}{Una superficie a pressione nulla e densita' finita}
\small
\begin{columns}[T,onlytextwidth]
\begin{column}{0.49\textwidth}
Al bordo di una configurazione senza
pressione esterna imponiamo $P(R)=0$:
\[
0=\frac{\varepsilon(R)-4B}{3}.
\]
Pertanto la densita' superficiale e'
\[
\varepsilon_s\equiv\varepsilon(R)=4B>0.
\]
La densita' di massa e' $\rho_s=\varepsilon_s/c^2$.
\end{column}
\begin{column}{0.49\textwidth}
Per confronto, il gas libero ha
$P=\varepsilon/3$ e dunque $P=0$
solo per $\varepsilon=0$.
\medskip
Il MIT bag model consente invece
una superficie a densita' finita.
\medskip
Per questa EOS lineare:
\[
c_s^2=\frac{dP}{d\varepsilon}=\frac13
\]
in unita' $c=1$.
\end{column}
\end{columns}
\medskip
La proprieta' $\varepsilon_s>0$ non dimostra da sola che la materia
strange sia piu' stabile della materia nucleare.
\end{frame}

% 04
\begin{frame}{Energia per barione della materia $uds$ massless}
\small
\begin{columns}[T,onlytextwidth]
\begin{column}{0.49\textwidth}
In equilibrio debole e neutralita':
\[
\mu_u=\mu_d=\mu_s=\mu_q,
\]
\[
n_u=n_d=n_s=\frac{\mu_q^3}{\pi^2}.
\]
Ogni quark ha numero barionico $1/3$:
\begin{align*}
n_B&=\frac{n_u+n_d+n_s}{3}\\
&=\frac{\mu_q^3}{\pi^2}.
\end{align*}
\end{column}
\begin{column}{0.49\textwidth}
La densita' di energia e'
\[
\varepsilon=\frac{9\mu_q^4}{4\pi^2}+B.
\]
L'energia per barione $E/A=\varepsilon/n_B$:
\begin{align*}
\frac EA&=\frac{9\mu_q^4/(4\pi^2)+B}
{\mu_q^3/\pi^2}\\
&=\frac94\mu_q+\frac{\pi^2B}{\mu_q^3}.
\end{align*}
\end{column}
\end{columns}
\end{frame}

% 05
\begin{frame}{Il minimo di $E/A$ si trova a $P=0$}
\small
\begin{columns}[T,onlytextwidth]
\begin{column}{0.49\textwidth}
Deriviamo rispetto al potenziale chimico
$\mu_q$ a $B$ costante:
\[
\frac EA=\frac94\mu_q+\pi^2B\mu_q^{-3}.
\]
\begin{align*}
\frac{d(E/A)}{d\mu_q}
&=\frac94-3\pi^2B\mu_q^{-4}.
\end{align*}
Al minimo:
\[
\frac94=\frac{3\pi^2B}{\mu_q^4},
\qquad
B=\frac{3\mu_q^4}{4\pi^2}.
\]
\end{column}
\begin{column}{0.49\textwidth}
La pressione del gas $uds$ e'
\[
P_q=\frac{3\mu_q^4}{4\pi^2}.
\]
La condizione di minimo diventa $B=P_q$:
\[
P=P_q-B=0.
\]
In generale a $T=0$, per materia neutra
in equilibrio:
\[
P=n_B^2\frac{d}{dn_B}
\left(\frac{\varepsilon}{n_B}\right).
\]
Quindi un minimo di $E/A$ a $n_B>0$
coincide con la pressione nulla.
\end{column}
\end{columns}
\end{frame}

\subsection{Struttura delle stelle autolegate}

% 06
\begin{frame}{Ramo di piccola massa: relazione massa--raggio}
\small
\begin{columns}[T,onlytextwidth]
\begin{column}{0.49\textwidth}
Per masse piccole, la gravita' comprime
poco la materia autolegata:
\[
\varepsilon(r)\simeq\varepsilon_s=4B.
\]
A $c$ esplicito, $\rho\simeq4B/c^2$.
\begin{align*}
M&=\int_0^R 4\pi r^2\rho(r)\,dr\\
&\simeq\frac{4\pi}{3}\rho R^3\\
&=\frac{16\pi B}{3c^2}R^3.
\end{align*}
\end{column}
\begin{column}{0.49\textwidth}
La legge corretta e'
\[
M\propto R^3,\qquad R\propto M^{1/3}.
\]
Quindi, idealmente,
\[
M\to0\quad\Rightarrow\quad R\to0.
\]
\medskip
\textit{Nota sul testo:} l'equazione (8.38)
riporta $M\propto R^{-3}$, incompatibile
con la derivazione da densita' costante:
e' un apparente refuso.
\end{column}
\end{columns}
\end{frame}

% 07
\begin{frame}{Equazioni TOV per un'EOS autolegata}
\footnotesize
\begin{columns}[T,onlytextwidth]
\begin{column}{0.49\textwidth}
In unita' $G=c=1$:
\[
\frac{dm}{dr}=4\pi r^2\varepsilon,
\]
\[
\frac{dP}{dr}=-(\varepsilon+P)
\frac{m+4\pi r^3P}{r(r-2m)}.
\]
$m(r)$: massa gravitazionale racchiusa
entro il raggio $r$.
\medskip
Le condizioni centrali sono
$m(0)=0$, $P(0)=P_c$; la superficie
si trova quando $P(R)=0$.
\end{column}
\begin{column}{0.49\textwidth}
Per il MIT bag model massless:
\[
P=\frac{\varepsilon-4B}{3}.
\]
Introduciamo grandezze adimensionali:
\[
\widetilde\varepsilon=\frac\varepsilon{4B},
\quad\widetilde P=\frac P{4B},
\]
\[
\widetilde r=\sqrt{4B}\,r,
\quad\widetilde m=\sqrt{4B}\,m.
\]
Allora
\[
\widetilde P=\frac{\widetilde\varepsilon-1}{3}.
\]
\end{column}
\end{columns}
\end{frame}

% 08
\begin{frame}{Perche' massa e raggio scalano come $B^{-1/2}$?}
\small
\begin{columns}[T,onlytextwidth]
\begin{column}{0.49\textwidth}
Sostituiamo le variabili senza dimensioni
nelle equazioni TOV.
\medskip
Per esempio:
\begin{align*}
\frac{dm}{dr}
&=\frac{d\widetilde m}{d\widetilde r}\\
&=4\pi\widetilde r^{\,2}
\widetilde\varepsilon.
\end{align*}
Anche nell'equazione di $dP/dr$
il fattore $B$ scompare.
\end{column}
\begin{column}{0.49\textwidth}
Le soluzioni adimensionali sono le stesse
per qualunque $B$ fissata la pressione
centrale adimensionale.
\medskip
Ritornando alle unita' fisiche:
\[
R\propto B^{-1/2},\qquad M\propto B^{-1/2}.
\]
In particolare:
\[
M_{\max},R_{\max}\propto B^{-1/2}.
\]
Ponendo $b=B^{1/4}$:
\[
M_{\max},R_{\max}\propto b^{-2}.
\]
\end{column}
\end{columns}
\end{frame}

% 09
\begin{frame}{Due leggi massa--raggio da non confondere}
\small
\begin{columns}[T,onlytextwidth]
\begin{column}{0.49\textwidth}
\textbf{Una EOS fissata; stelle leggere}
\[
M\simeq\frac{16\pi B}{3c^2}R^3.
\]
Qui $B$ rimane costante e variamo
il numero di barioni nella stella.
\medskip
La densita' e' quasi uniforme
$\varepsilon\simeq4B$.
\end{column}
\begin{column}{0.49\textwidth}
\textbf{EOS diverse; stessa soluzione adimensionale}
\[
M,R\propto B^{-1/2}.
\]
Qui cambiamo $B$ e confrontiamo
configurazioni TOV corrispondenti.
\medskip
La massa massima diminuisce
se aumenta la densita' caratteristica $B$.
\end{column}
\end{columns}
\medskip
Il ramo $M\propto R^3$ cessa di essere accurato per stelle massicce,
quando la compressione gravitazionale diventa importante.
\end{frame}


\begin{frame}{Origine dell'energia di legame nucleare}
\small
\begin{columns}[T,onlytextwidth]
\begin{column}{0.49\textwidth}

L'energia di legame di un nucleo
con $A$ nucleoni e $Z$ protoni è
\[
E_b=Zm_p+(A-Z)m_n-M(A,Z).
\]

La formula semiempirica delle masse:
\[
\begin{aligned}
E_b\simeq{}&
a_v A-a_s A^{2/3}\\
&-a_c\frac{Z(Z-1)}{A^{1/3}}\\
&-a_{\rm sym}\frac{(A-2Z)^2}{A}
+\delta .
\end{aligned}
\]

Il coefficiente di volume è
\[
a_v\simeq 15.5\text{--}16\,\mathrm{MeV}.
\]

\end{column}
\begin{column}{0.49\textwidth}

Per materia nucleare infinita,
simmetrica e senza Coulomb:
\[
A\to\infty,\qquad Z/A=1/2.
\]

I contributi di superficie e pairing
per nucleone svaniscono;
quello di simmetria è nullo.

Rimane il contributo di volume:
\[
\frac{E_b}{A}\simeq a_v.
\]

Quindi
\[
\frac EA\simeq m_N-16\,\mathrm{MeV}
\simeq923\,\mathrm{MeV}.
\]

Il minimo si trova alla densità
di saturazione:
\[
n_0\simeq0.16\,\mathrm{fm}^{-3},
\qquad P(n_0)=0.
\]

\end{column}
\end{columns}
\end{frame}


% 10
\begin{frame}{Perche' i nuclei sono autolegati ma le stelle di neutroni no?}
\footnotesize
\begin{columns}[T,onlytextwidth]
\begin{column}{0.49\textwidth}
La materia nucleare \emph{simmetrica} a $T=0$
ha un minimo dell'energia per nucleone:
\[
\left.\frac EA\right|_{n_0}
\simeq m_N-16\,\mathrm{MeV},
\]
\[
n_0\simeq0.15\text{--}0.16\,\mathrm{fm}^{-3},
\qquad P(n_0)=0.
\]
E' la base fenomenologica dell'autolegamento
dei nuclei, senza bisogno della gravita'.
\end{column}
\begin{column}{0.49\textwidth}
La materia stellare deve essere
neutra ed in equilibrio beta:
\[
n_p=n_e+n_\mu,
\quad\mu_n=\mu_p+\mu_e.
\]
E' molto ricca di neutroni: non coincide
con materia nucleare simmetrica.
\medskip
Nelle EOS nucleoniche standard non e'
autolegata: la gravita' ne assicura
l'equilibrio globale.
\end{column}
\end{columns}
\medskip
Materia $uds$ puo' invece essere quasi neutra senza leptoni,
perche' $2/3-1/3-1/3=0$.
\end{frame}

\subsection{Ipotesi di Bodmer--Witten e vincoli di stabilita'}

% 11
\begin{frame}{Ipotesi di Bodmer--Witten}
\small
\begin{columns}[T,onlytextwidth]
\begin{column}{0.49\textwidth}
La materia strange uniforme $uds$
potrebbe essere il vero stato fondamentale
della materia barionica fredda.
\medskip
Si deve confrontare l'energia per barione
a \emph{pressione nulla} con quella dei nuclei:
\[
\left.\frac EA\right|_{uds,P=0}
\stackrel{?}{<}
\left.\frac EA\right|_{\rm nuclei}.
\]
\end{column}
\begin{column}{0.49\textwidth}
Un nucleo ordinario puo' restare
metastabile anche se esistesse
uno stato strange piu' favorevole.
\medskip
La produzione di quark $s$ richiede
processi deboli; possono esserci
barriere di nucleazione e correzioni
di dimensione finita.
\medskip
\textit{Stabilita' termodinamica} e
\textit{rapidita' della conversione} sono
due questioni differenti.
\end{column}
\end{columns}
\end{frame}

% 12
\begin{frame}{I due requisiti di stabilita' (8.39)--(8.40)}
\small
\begin{columns}[T,onlytextwidth]
\begin{column}{0.49\textwidth}
\textbf{Due sapori} $N_f=2$: materia $ud$.
\[
\left.\frac EA\right|_{ud,P=0}
>934\,\mathrm{MeV}.
\]
La materia $ud$ non deve essere
piu' stabile dei nuclei ordinari.
\medskip
Il valore $934\,\mathrm{MeV}$ include
una correzione fenomenologica di
superficie di circa $4\,\mathrm{MeV}$
rispetto al confronto con $^{56}\mathrm{Fe}$.
\end{column}
\begin{column}{0.49\textwidth}
\textbf{Tre sapori} $N_f=3$: materia $uds$.
\[
\left.\frac EA\right|_{uds,P=0}
<930\,\mathrm{MeV}.
\]
La materia strange in bulk deve risultare
piu' stabile della materia nucleare.
\medskip
$N_f$ e' il numero di sapori di quark
presenti. Entrambi i vincoli riguardano
materia macroscopica a $T=0$ e $P=0$.
\end{column}
\end{columns}
\end{frame}

% 13
\begin{frame}{Perche' $E/A=\mu_B$ quando $P=0$?}
\small
\begin{columns}[T,onlytextwidth]
\begin{column}{0.49\textwidth}
L'identita' termodinamica a $T=0$ e'
\[
\varepsilon+P=\sum_i\mu_i n_i.
\]
Per specie in equilibrio debole:
\[
\mu_i=B_i\mu_B+Q_i\mu_Q.
\]
$B_i$ e $Q_i$ sono numero barionico
e carica elettrica della specie $i$.
\end{column}
\begin{column}{0.49\textwidth}
Sommando sulle specie:
\begin{align*}
\sum_i\mu_i n_i
&=\mu_B\sum_iB_i n_i
+\mu_Q\sum_iQ_i n_i\\
&=\mu_Bn_B+\mu_Qn_Q.
\end{align*}
La neutralita' impone $n_Q=0$:
\[
\varepsilon+P=\mu_Bn_B.
\]
Quindi, se $P=0$,
\[
\frac EA\equiv\frac\varepsilon{n_B}=\mu_B.
\]
\end{column}
\end{columns}
\end{frame}

\subsection{Vincolo a due sapori}

% 14
\begin{frame}{Materia $ud$: cariche, densita' e approssimazioni}
\small
\begin{columns}[T,onlytextwidth]
\begin{column}{0.49\textwidth}
Numeri quantici:
\[
(B_u,Q_u)=\left(\tfrac13,\tfrac23\right),
\quad(B_d,Q_d)=\left(\tfrac13,-\tfrac13\right).
\]
Per due quark massless:
\[
n_u=\frac{\mu_u^3}{\pi^2},
\quad n_d=\frac{\mu_d^3}{\pi^2}.
\]
Le degenerazioni di colore e spin
sono gia' incluse: $\gamma=3\cdot2=6$.
\end{column}
\begin{column}{0.49\textwidth}
Il testo trascura il contributo elettronico
alla neutralita' e alla pressione:
\[
\frac23n_u-\frac13n_d\simeq0.
\]
Gli elettroni sono in realta' presenti
in equilibrio beta, perche'
\[
\mu_e=\mu_d-\mu_u\ne0.
\]
Dunque si tratta di un'approssimazione,
non della soluzione esatta $ude$.
\end{column}
\end{columns}
\end{frame}

% 15
\begin{frame}{Materia $ud$: dalla neutralita' ai potenziali chimici}
\small
\begin{columns}[T,onlytextwidth]
\begin{column}{0.49\textwidth}
Dalla neutralita' approssimata:
\begin{align*}
\frac23n_u-\frac13n_d&=0,\\
2n_u&=n_d.
\end{align*}
Usando $n_i=\mu_i^3/\pi^2$:
\[
2\mu_u^3=\mu_d^3.
\]
Quindi
\[
\mu_d=2^{1/3}\mu_u.
\]
\end{column}
\begin{column}{0.49\textwidth}
Le cariche conservate danno
\[
\mu_u=\frac13\mu_B-\frac23\mu_e,
\]
\[
\mu_d=\frac13\mu_B+\frac13\mu_e.
\]
Combinandole:
\begin{align*}
\mu_u+2\mu_d
&=\mu_B
+\left(-\tfrac23+\tfrac23\right)\mu_e\\
&=\mu_B.
\end{align*}
Perciò $\mu_B=(1+2^{4/3})\mu_u$.
\end{column}
\end{columns}
\end{frame}

% 16
\begin{frame}{Materia $ud$: pressione nulla (8.41)}
\small
\begin{columns}[T,onlytextwidth]
\begin{column}{0.49\textwidth}
Il gas libero di due sapori esercita
\[
P_q=\frac{\mu_u^4+\mu_d^4}{4\pi^2}.
\]
Il MIT bag model fornisce
\[
P=P_q-B.
\]
In assenza di pressione esterna
($P=0$) deve valere $P_q=B$.
\end{column}
\begin{column}{0.49\textwidth}
Con $\mu_d=2^{1/3}\mu_u$:
\begin{align*}
B&=\frac{\mu_u^4+\mu_d^4}{4\pi^2}\\
&=\frac{1+2^{4/3}}{4\pi^2}\mu_u^4.
\end{align*}
Dunque
\[
\mu_u^4=\frac{4\pi^2B}{1+2^{4/3}}.
\]
Si tratta della relazione (8.41).
\end{column}
\end{columns}
\end{frame}

% 17
\begin{frame}{Materia $ud$: energia per barione (8.42)}
\small
\begin{columns}[T,onlytextwidth]
\begin{column}{0.49\textwidth}
A pressione nulla:
\[
\frac EA=\mu_B=(1+2^{4/3})\mu_u.
\]
Dalla pressione nulla:
\[
\mu_u=\left(
\frac{4\pi^2B}{1+2^{4/3}}
\right)^{1/4}.
\]
Sostituiamo l'espressione di $\mu_u$.
\end{column}
\begin{column}{0.49\textwidth}
\begin{align*}
\frac EA
&=(1+2^{4/3})
\left(\frac{4\pi^2B}{1+2^{4/3}}\right)^{1/4}\\[2pt]
&=\left[4\pi^2(1+2^{4/3})^3B\right]^{1/4}.
\end{align*}
Vincolo nucleare:
\[
\frac EA>934\,\mathrm{MeV}.
\]
Quindi
\[
B^{1/4}>\frac{934\,\mathrm{MeV}}
{[4\pi^2(1+2^{4/3})^3]^{1/4}}
\simeq145\,\mathrm{MeV}.
\]
\end{column}
\end{columns}
\end{frame}

\subsection{Vincolo a tre sapori}

% 18
\begin{frame}{Materia $uds$ massless: soluzione neutra}
\small
\begin{columns}[T,onlytextwidth]
\begin{column}{0.49\textwidth}
Cariche dei tre sapori:
\[
Q_u=+\tfrac23,\quad
Q_d=Q_s=-\tfrac13.
\]
La soluzione simmetrica e'
\[
n_u=n_d=n_s,\quad n_e=0.
\]
Infatti
\[
n_Q=\frac23n_u-\frac13n_d-\frac13n_s=0.
\]
\end{column}
\begin{column}{0.49\textwidth}
Per quark massless:
\[
\mu_u=\mu_d=\mu_s=\mu_q,
\]
\[
n_i=\frac{\mu_q^3}{\pi^2}.
\]
Il numero barionico totale e'
\[
n_B=\frac13(n_u+n_d+n_s)
=\frac{\mu_q^3}{\pi^2}.
\]
Il potenziale barionico e' $\mu_B=3\mu_q$.
\end{column}
\end{columns}
\end{frame}

% 19
\begin{frame}{Materia $uds$: pressione nulla e $E/A$ (8.43)}
\small
\begin{columns}[T,onlytextwidth]
\begin{column}{0.49\textwidth}
La pressione del gas massless e'
\[
P_q=\frac{3\mu_q^4}{4\pi^2}.
\]
Alla superficie autolegata:
\[
P=P_q-B=0,
\]
quindi
\[
B=\frac{3\mu_q^4}{4\pi^2},
\qquad \mu_q^4=\frac{4\pi^2B}{3}.
\]
\end{column}
\begin{column}{0.49\textwidth}
Per $P=0$, $E/A=\mu_B=3\mu_q$:
\begin{align*}
\left(\frac EA\right)^4
&=3^4\mu_q^4\\
&=3^4\frac{4\pi^2B}{3}\\
&=3^3\,4\pi^2B.
\end{align*}
Pertanto
\[
\frac EA=(3^3\,4\pi^2B)^{1/4}.
\]
\end{column}
\end{columns}
\end{frame}

% 20
\begin{frame}{Materia $uds$: vincolo di stabilita' assoluta}
\small
\begin{columns}[T,onlytextwidth]
\begin{column}{0.49\textwidth}
Condizione di Bodmer--Witten:
\[
\left.\frac EA\right|_{uds}<930\,\mathrm{MeV}.
\]
Sostituendo il risultato precedente:
\[
(3^3\,4\pi^2B)^{1/4}
<930\,\mathrm{MeV}.
\]
\end{column}
\begin{column}{0.49\textwidth}
Isoliamo la scala del bag:
\[
B^{1/4}<\frac{930\,\mathrm{MeV}}
{(3^3\,4\pi^2)^{1/4}}.
\]
Numericamente:
\[
B^{1/4}<163\,\mathrm{MeV}.
\]
Questa e' una soglia approssimata
nel limite di tre sapori massless.
\end{column}
\end{columns}
\end{frame}

% 21
\begin{frame}{La finestra di stabilita' di materia strange}
\small
\begin{columns}[T,onlytextwidth]
\begin{column}{0.49\textwidth}
Combinando i due vincoli:
\[
\begin{aligned}
N_f=2:&\quad B^{1/4}>145\,\mathrm{MeV},\\
N_f=3:&\quad B^{1/4}<163\,\mathrm{MeV}.
\end{aligned}
\]
La finestra e'
\[
145<B^{1/4}<163\,\mathrm{MeV}.
\]
I valori dipendono dalle ipotesi
su masse, interazioni e superfici.
\end{column}
\begin{column}{0.49\textwidth}
Attenzione alle unita': $B$ e' una
densita' di energia; in unita' naturali
$[B]=\mathrm{MeV}^4$.
\[
B[\mathrm{MeV\,fm^{-3}}]
=\frac{(B^{1/4}[\mathrm{MeV}])^4}
{(\hbar c)^3}.
\]
Con $\hbar c=197.327\,\mathrm{MeV\,fm}$:
\[
57.5<B<91.9\,\mathrm{MeV\,fm^{-3}}.
\]
\end{column}
\end{columns}
\end{frame}

% 22
\begin{frame}{Esempio: energia di legame della materia $uds$}
\small
\begin{columns}[T,onlytextwidth]
\begin{column}{0.49\textwidth}
Alla soglia inferiore $B^{1/4}=145\,\mathrm{MeV}$:
\[
\left.\frac EA\right|_{uds}
=(3^3\,4\pi^2)^{1/4}B^{1/4}.
\]
Il coefficiente e' circa $5.714$:
\[
\frac EA\simeq5.714\times145\,\mathrm{MeV}
\simeq829\,\mathrm{MeV}.
\]
\end{column}
\begin{column}{0.49\textwidth}
Rispetto al riferimento nucleare:
\[
\Delta e\simeq930-829
\simeq101\,\mathrm{MeV/barione}.
\]
La materia $uds$ potrebbe dunque essere
molto legata, \emph{all'interno del modello}.
\medskip
Questo non dimostra sperimentalmente
l'esistenza di strange quark matter
assolutamente stabile.
\end{column}
\end{columns}
\end{frame}

\subsection{Strangelet, crosta ed osservazioni}

% 23
\begin{frame}{Dalla materia infinita agli strangelet}
\small
\begin{columns}[T,onlytextwidth]
\begin{column}{0.49\textwidth}
Per una goccia con numero barionico $A$:
\[
E(A)\simeq Ae_{\rm bulk}
+E_{\rm surf}+E_{\rm curv}+E_{\rm Coul}.
\]
Se la densita' e' quasi costante:
\[
R\propto A^{1/3}.
\]
La superficie scala come
$E_{\rm surf}\propto R^2\propto A^{2/3}$.
\end{column}
\begin{column}{0.49\textwidth}
Pertanto le correzioni per barione sono
\[
\frac{E_{\rm surf}}A\propto A^{-1/3},
\qquad
\frac{E_{\rm curv}}A\propto A^{-2/3}.
\]
Per gocce piccole queste correzioni
possono rendere $E/A$ maggiore
che nella materia strange infinita.
\medskip
La stabilita' in bulk non garantisce
la stabilita' di ogni strangelet.
\end{column}
\end{columns}
\end{frame}

% 24
\begin{frame}{Perche' i nuclei non devono decadere rapidamente?}
\small
\begin{columns}[T,onlytextwidth]
\begin{column}{0.49\textwidth}
La conversione $ud\to uds$ richiede
una variazione della stranezza:
\[
\Delta S\ne0.
\]
Le interazioni forti conservano
la stranezza; quelle deboli possono
modificarla.
\medskip
Una fase con energia libera inferiore
non e' necessariamente raggiunta
in tempi brevi.
\end{column}
\begin{column}{0.49\textwidth}
Altre barriere possibili:
\begin{itemize}
\item costo di nucleazione di una goccia;
\item energia superficiale e di curvatura;
\item contributi coulombiani;
\item dimensione minima del germe stabile.
\end{itemize}
La sopravvivenza dei nuclei ordinari
e' quindi compatibile, in linea di principio,
con una fase strange piu' stabile.
\end{column}
\end{columns}
\end{frame}

% 25
\begin{frame}{Stella strange nuda o con crosta nucleare}
\small
\begin{columns}[T,onlytextwidth]
\begin{column}{0.49\textwidth}
Se $m_s>0$, la popolazione dei quark
strange si riduce a parita' di $\mu_s$.
\medskip
La neutralita' richiede una piccola
popolazione elettronica, che puo'
estendersi oltre la superficie dei quark.
\medskip
Il campo elettrostatico puo' sostenere
una crosta di nuclei ordinari.
\end{column}
\begin{column}{0.49\textwidth}
Possibili strutture:
\begin{itemize}
\item \emph{stella nuda}: superficie del core $uds$;
\item \emph{stella con crosta}: core $uds$,
strato elettronico e nuclei.
\end{itemize}
Una crosta ordinaria elettrostaticamente
sostenuta resta in genere sotto
la soglia di neutron drip: i neutroni
liberi non sono sostenuti dalla
barriera coulombiana.
\end{column}
\end{columns}
\end{frame}

% 26
\begin{frame}{Massa massima: dipendenza da $B$}
\small
\begin{columns}[T,onlytextwidth]
\begin{column}{0.49\textwidth}
Per la EOS lineare:
\[
M_{\max}\propto B^{-1/2}
\propto(B^{1/4})^{-2}.
\]
Il testo riporta indicativamente
\[
M_{\max}(145\,\mathrm{MeV})
\simeq2.01\,M_\odot,
\]
dove il numero $145\,\mathrm{MeV}$
e' $B^{1/4}$, non $B$.
\end{column}
\begin{column}{0.49\textwidth}
Al limite alto $B^{1/4}=163\,\mathrm{MeV}$:
\begin{align*}
M_{\max}(163)
&\simeq2.01M_\odot
\left(\frac{145}{163}\right)^2\\
&\simeq1.59M_\odot.
\end{align*}
Dunque il bag model piu' semplice
e' in tensione con stelle compatte
di massa superiore a circa $2M_\odot$.
\end{column}
\end{columns}
\end{frame}

% 27
\begin{frame}{Due punti da verificare nel testo del \S8.3}
\footnotesize
\begin{columns}[T,onlytextwidth]
\begin{column}{0.49\textwidth}
\textbf{Equazione (8.38).}
\medskip
Il libro riporta $M\propto R^{-3}$ per
stelle autolegate di bassa massa.
\medskip
Con $\varepsilon\simeq4B$ costante:
\[
M=\frac{4\pi}{3}\rho R^3.
\]
Quindi l'esponente fisicamente
coerente e' $+3$.
\end{column}
\begin{column}{0.49\textwidth}
\textbf{Intervallo per il bag parameter.}
\medskip
Il testo usa anche $B=145$--$235\,
\mathrm{MeV\,fm^{-3}}$, ma altrove
scrive $B^{1/4}=145\,\mathrm{MeV}$.
\medskip
Sono grandezze diverse:
\[
B^{1/4}=145\,\mathrm{MeV}
\Longleftrightarrow
B\simeq57.5\,\mathrm{MeV\,fm^{-3}}.
\]
L'intervallo iniziale necessita
quindi di verifica delle unita'.
\end{column}
\end{columns}
\end{frame}

% 28
\begin{frame}{Limiti del modello e conclusione fisica}
\small
\begin{columns}[T,onlytextwidth]
\begin{column}{0.49\textwidth}
Ipotesi semplificative fin qui:
\begin{itemize}
\item quark $u,d,s$ massless;
\item gas non interagente;
\item bag constant $B$ indipendente
 dalla densita';
\item confronto con nuclei e correzioni
 finite-size approssimate.
\end{itemize}
La massa $m_s$, l'interazione forte e
l'eventuale pairing possono cambiare
la finestra di stabilita' e l'EOS.
\end{column}
\begin{column}{0.49\textwidth}
Risultati da conservare:
\[
P=\frac{\varepsilon-4B}{3},
\quad \varepsilon_s=4B,
\]
\[
M\propto R^3\quad\text{(piccole masse)},
\]
\[
M_{\max}\propto B^{-1/2}.
\]
Nel modello massless semplificato:
\[
145<B^{1/4}<163\,\mathrm{MeV}.
\]
\end{column}
\end{columns}
\medskip
Autolegamento ($P=0$ a $n_B>0$) e stabilita' assoluta
($E/A$ inferiore a quello nucleare) sono condizioni distinte.
\end{frame}

\section{Materia di quark interagente: pQCD}
\subsection{Dal gas libero alla pressione perturbativa}

% 01
\begin{frame}{Dal gas libero alla materia di quark interagente}
\small
\begin{columns}[T,onlytextwidth]
\begin{column}{0.49\textwidth}
\textbf{Nei paragrafi precedenti}
\[
P_{\rm free}=\frac{N_f}{4\pi^2}\mu_q^4,
\]
\[
P_{\rm MIT}=P_{\rm free}-B.
\]
Abbiamo considerato un gas degenere di quark liberi
ed eventualmente un contributo costante del vuoto.
\end{column}
\begin{column}{0.49\textwidth}
\textbf{Nuova domanda}

In che modo lo scambio di gluoni modifica
la pressione e l'equazione di stato?
\medskip
Il \S8.4 usa la QCD perturbativa (pQCD)
a $T=0$, assumendo $N_f=3$ sapori massless.
\medskip
Distingueremo $P(\mu_q)$ da $P(\varepsilon)$.
\end{column}
\end{columns}
\medskip
Per tre sapori massless neutri: $\mu_u=\mu_d=\mu_s=\mu_q$ e
$\mu_B=3\mu_q$.
\end{frame}

% 02
\begin{frame}{Pressione ideale a temperatura nulla}
\small
\begin{columns}[T,onlytextwidth]
\begin{column}{0.49\textwidth}
Per ogni sapore massless:
\[
g_q=2\;({\rm spin})\times3\;({\rm colori})=6.
\]
A $T=0$, $k_F=\mu_q$ e
\[
P_f=\frac{g_q}{6\pi^2}
\int_0^{k_F}\frac{k^4}{E_k}\,dk.
\]
Poiche' $E_k=k$:
\[
P_f=\frac6{6\pi^2}\int_0^{\mu_q}k^3\,dk.
\]
\end{column}
\begin{column}{0.49\textwidth}
Integrando:
\[
P_f=\frac{\mu_q^4}{4\pi^2}.
\]
Sommando $N_f$ sapori con uguale $\mu_q$:
\[
P_{\rm SB}(\mu_q)=
\frac{N_f\mu_q^4}{4\pi^2}.
\]
Qui ``Stefan--Boltzmann'' indica
il \emph{limite ideale a $T=0$},
non un plasma termico di gluoni.
\end{column}
\end{columns}
\end{frame}

% 03
\begin{frame}{Due diversi limiti di Stefan--Boltzmann}
\small
\begin{columns}[T,onlytextwidth]
\begin{column}{0.49\textwidth}
\textbf{\S8.1: alta temperatura, $\mu=0$}
\[
P_{\rm SB}(T)=\frac{\pi^2}{90}
\left(16+\frac{21}{2}N_f\right)T^4.
\]
\begin{itemize}
\item Gluoni termici: $8\times2=16$ stati.
\item Quark e antiquark: $12N_f$ stati fermionici.
\item La pressione cresce come $T^4$.
\end{itemize}
\end{column}
\begin{column}{0.49\textwidth}
\textbf{\S8.4: materia fredda, $T=0$}
\[
P_{\rm SB}(\mu_q)=
\frac{N_f}{4\pi^2}\mu_q^4.
\]
\begin{itemize}
\item Quark degeneri con $\mu_q>0$.
\item Nessun gas termico di gluoni liberi.
\item Gluoni presenti negli scambi virtuali.
\end{itemize}
\end{column}
\end{columns}
\medskip
In Figura 8.3 il rapporto mostrato e' $P(\mu_B)/P_{\rm SB}(\mu_B)$.
\end{frame}

% 04
\begin{frame}{pQCD: la pressione fino all'ordine $\alpha_s^2$}
\footnotesize
Per $N_f$ quark massless con uguale potenziale $\mu_q$:
\[
\begin{aligned}
P(\mu_q)=\frac{N_f\mu_q^4}{4\pi^2}
\Bigg\{1
&-2\frac{\alpha_s}{\pi}\\[-1mm]
&-\left[
G+N_f\ln\!\left(\frac{\alpha_s}{\pi}\right)
+\left(11-\frac23N_f\right)
\ln\!\left(\frac{\bar\Lambda}{\mu_q}\right)
\right]
\left(\frac{\alpha_s}{\pi}\right)^2\Bigg\}.
\end{aligned}
\]
\[
G=10.376-0.536N_f+N_f\ln N_f,
\qquad \alpha_s=\alpha_s(\bar\Lambda).
\]
\begin{columns}[T,onlytextwidth]
\begin{column}{0.49\textwidth}
$\mu_q$: potenziale chimico dei quark.

$\alpha_s=g_s^2/(4\pi)$: accoppiamento forte.
\end{column}
\begin{column}{0.49\textwidth}
$\bar\Lambda$: scala di rinormalizzazione (energia).

$G$: coefficiente del calcolo perturbativo.
\end{column}
\end{columns}
\end{frame}

% 05
\begin{frame}{La struttura dell'espansione perturbativa}
\small
Definiamo $a=\alpha_s/\pi$ e $c=11-2N_f/3$:
\[
\frac{P}{P_{\rm SB}}
=1-2a-\biggl[G+N_f\ln a+c\ln\!\left(\frac{\bar\Lambda}{\mu_q}\right)\biggr]a^2.
\]
\begin{columns}[T,onlytextwidth]
\begin{column}{0.49\textwidth}
\textbf{Ordine zero}
\[
\frac{P}{P_{\rm SB}}=1.
\]
Gas ideale di quark massless.
\medskip
\textbf{Primo ordine}
\[
\frac{P}{P_{\rm SB}}\simeq1-2a.
\]
La correzione e' negativa.
\end{column}
\begin{column}{0.49\textwidth}
\textbf{Secondo ordine}
\[
-\left[G+N_f\ln a+c\ln\frac{\bar\Lambda}{\mu_q}\right]a^2.
\]
Compaiono $\alpha_s^2$, $\alpha_s^2\ln\alpha_s$
e il logaritmo della scala di rinormalizzazione.
\medskip
Non e' una semplice serie polinomiale in $\alpha_s$.
\end{column}
\end{columns}
\end{frame}

% 06
\begin{frame}{Primo ordine: riduzione della pressione ideale}
\small
\[
\frac{P}{P_{\rm SB}}\simeq1-\frac{2\alpha_s}{\pi}.
\]
\begin{columns}[T,onlytextwidth]
\begin{column}{0.49\textwidth}
A parita' di potenziale chimico, l'interazione
riduce la pressione rispetto al gas libero:
\[
\Delta P\simeq-
\frac{2\alpha_s}{\pi}P_{\rm SB}.
\]
Si tratta della \emph{prima} correzione:
non contiene ancora i termini $O(\alpha_s^2)$.
\end{column}
\begin{column}{0.49\textwidth}
\begin{tabular}{c|c|c}
$\alpha_s$ & $P/P_{\rm SB}$ & Riduzione\\ \hline
$0.1$ & $0.936$ & $6.4\%$\\
$0.3$ & $0.809$ & $19.1\%$\\
$0.5$ & $0.682$ & $31.8\%$
\end{tabular}
\medskip
Per $\alpha_s$ grande, gli ordini trascurati
possono essere importanti.
\end{column}
\end{columns}
\end{frame}

% 07
\begin{frame}{Tre sapori e origine dei termini logaritmici}
\small
\begin{columns}[T,onlytextwidth]
\begin{column}{0.49\textwidth}
Per $N_f=3$:
\[
c=11-\frac23\cdot3=9,
\]
\[
G=10.376-1.608+3\ln3
\simeq12.064.
\]
La correzione contiene
\[
\alpha_s^2\ln\alpha_s
\quad\text{e}\quad
\alpha_s^2\ln(\bar\Lambda/\mu_q).
\]
\end{column}
\begin{column}{0.49\textwidth}
Il termine $\ln\alpha_s$ riflette
la struttura non banale della teoria perturbativa
in materia densa.
\medskip
A temperatura finita alcune risommazioni
producono anche potenze dispari di $g_s$.
\medskip
La forma qui riportata e' specifica
dell'espansione a $T=0$ fino a $\alpha_s^2$.
\end{column}
\end{columns}
\end{frame}

\subsection{Significato della scala di rinormalizzazione in QCD}

% 08
\begin{frame}{Perche' serve rinormalizzare la QCD?}
\small
\begin{columns}[T,onlytextwidth]
\begin{column}{0.49\textwidth}
I diagrammi quantistici con loop possono
contenere integrali ultravioletti divergenti.
\medskip
Si introduce una regolarizzazione,
per esempio la regolarizzazione dimensionale:
\[
d=4-2\epsilon.
\]
I termini divergenti compaiono come poli
$1/\epsilon$ e vengono assorbiti nei
parametri e campi rinormalizzati.
\end{column}
\begin{column}{0.49\textwidth}
Il coupling che compare nei calcoli non e'
una costante universale fissata una volta
per tutte:
\[
g_s=g_s(\bar\Lambda),
\qquad
\alpha_s(\bar\Lambda)=\frac{g_s^2(\bar\Lambda)}{4\pi}.
\]
$\bar\Lambda$ e' una scala di energia introdotta
nella \emph{definizione del coupling rinormalizzato}.
\end{column}
\end{columns}
\end{frame}

% 09
\begin{frame}{Che cos'e' precisamente $\bar\Lambda$?}
\small
\begin{columns}[T,onlytextwidth]
\begin{column}{0.49\textwidth}
\textbf{Significato operativo}

$\bar\Lambda$ e' la scala rispetto alla quale
si definiscono i parametri rinormalizzati,
per esempio in schema $\overline{\mathrm{MS}}$.
\medskip
Ha dimensioni di energia; si puo' scegliere
$\bar\Lambda=\kappa\mu_q$.
\medskip
Non e' un nuovo campo, ne' la densita'
fisica del sistema; non e' nemmeno
un taglio fisico sui momenti dei gluoni.
\end{column}
\begin{column}{0.49\textwidth}
\textbf{Perche' compare nei logaritmi?}

I risultati rinormalizzati contengono rapporti
fra la scala introdotta e le scale fisiche:
\[
\ln\frac{\bar\Lambda}{Q},
\qquad Q\sim\mu_q.
\]
Scegliere $\bar\Lambda$ dell'ordine di $Q$
evita logaritmi inutilmente grandi e spesso
migliora l'organizzazione della serie.
\end{column}
\end{columns}
\end{frame}

% 10
\begin{frame}{Tre scale da non confondere}
\small
\begin{columns}[T,onlytextwidth]
\begin{column}{0.49\textwidth}
\textbf{Scala fisica: $\mu_q$}
\medskip
Fissata dallo stato della materia:
quali livelli fermionici sono occupati.
\[
\mu_B=3\mu_q.
\]
\textbf{Scala di rinormalizzazione: $\bar\Lambda$}
\medskip
Scelta convenzionale della descrizione
perturbativa; il risultato esatto non deve
dipenderne.
\end{column}
\begin{column}{0.49\textwidth}
\textbf{Scala caratteristica della QCD:
$\Lambda_{\overline{\rm MS}}$}
\medskip
Parametro dimensionale generato dal running
dell'accoppiamento in uno schema fissato.
\medskip
Non e' la stessa cosa di $\bar\Lambda$:
\[
\bar\Lambda\ \text{si varia},
\qquad
\Lambda_{\overline{\rm MS}}\ \text{e' fissata}.
\]
\end{column}
\end{columns}
\end{frame}

% 11
\begin{frame}{Come nasce $\Lambda_{\overline{\rm MS}}$ dal running?}
\small
\begin{columns}[T,onlytextwidth]
\begin{column}{0.49\textwidth}
Al primo ordine:
\[
\frac{d\alpha_s}{d\ln\bar\Lambda^2}
=-b_0\alpha_s^2,
\qquad
b_0=\frac{11-2N_f/3}{4\pi}.
\]
Separiamo le variabili:
\[
\frac{d\alpha_s}{\alpha_s^2}
=-b_0\,d\ln\bar\Lambda^2.
\]
Integrando:
\[
\frac1{\alpha_s(\bar\Lambda)}
=b_0\ln\frac{\bar\Lambda^2}
{\Lambda_{\overline{\rm MS}}^2}.
\]
\end{column}
\begin{column}{0.49\textwidth}
Quindi
\[
\alpha_s(\bar\Lambda)\simeq
\frac{4\pi}{(11-2N_f/3)
\ln(\bar\Lambda^2/\Lambda_{\overline{\rm MS}}^2)}.
\]
$\Lambda_{\overline{\rm MS}}$ e' una costante
di integrazione (schema e $N_f$ fissati).
\medskip
Il running e' una previsione della QCD;
la scelta di $\bar\Lambda$ non cambia
la fisica esatta.
\end{column}
\end{columns}
\end{frame}

% 12
\begin{frame}{Liberta' asintotica e scelta della scala}
\small
\begin{columns}[T,onlytextwidth]
\begin{column}{0.49\textwidth}
Per $N_f<16.5$, al primo ordine
\[
b_0>0,
\quad
\frac{d\alpha_s}{d\ln\bar\Lambda^2}<0.
\]
Quindi $\alpha_s$ diminuisce quando
$\bar\Lambda$ aumenta.
\medskip
E' la \textbf{liberta' asintotica}.
\end{column}
\begin{column}{0.49\textwidth}
In materia fredda e densa, una scelta naturale e'
\[
\bar\Lambda_0=2\mu_q.
\]
Per stimare la dipendenza residua:
\[
\mu_q\leq\bar\Lambda\leq4\mu_q.
\]
Il fattore $2$ e' una convenzione di calcolo,
non una relazione fisica esatta.
\end{column}
\end{columns}
\end{frame}

% 13
\begin{frame}{Perche' una grandezza fisica non deve dipendere da $\bar\Lambda$?}
\small
\begin{columns}[T,onlytextwidth]
\begin{column}{0.49\textwidth}
Una previsione \emph{esatta} e' indipendente
dalla scala ausiliaria:
\[
\left.\frac{dP_{\rm esatta}}{d\ln\bar\Lambda^2}
\right|_{\mu_q}=0.
\]
La derivata e' \emph{totale}:
include sia i logaritmi espliciti sia
la dipendenza implicita
$\alpha_s=\alpha_s(\bar\Lambda)$.
\end{column}
\begin{column}{0.49\textwidth}
Se tronchiamo la serie a $\alpha_s^2$,
la cancellazione funziona fino
all'ordine calcolato.
\medskip
Gli ordini omessi lasciano una dipendenza
residua:
\[
\frac{dP_{\rm pQCD}}{d\ln\bar\Lambda^2}
=O(\alpha_s^3)\,P_{\rm SB}.
\]
Questo rende utile variare $\bar\Lambda$.
\end{column}
\end{columns}
\end{frame}

\subsection{Cancellazione della dipendenza dalla scala di rinormalizzazione}

% 14
\begin{frame}{Preparazione della derivata: notazione compatta}
\small
Definiamo
\[
t=\ln\bar\Lambda^2,
\qquad a=\frac{\alpha_s}{\pi},
\qquad c=11-\frac23N_f=4\pi b_0.
\]
\[
F\equiv\frac{P}{P_{\rm SB}}
=1-2a-\left[G+N_f\ln a+c\ln(\bar\Lambda/\mu_q)\right]a^2.
\]
\begin{columns}[T,onlytextwidth]
\begin{column}{0.49\textwidth}
Manteniamo $\mu_q$ costante:
\[
\frac{d}{dt}\ln\frac{\bar\Lambda}{\mu_q}
=\frac12.
\]
\end{column}
\begin{column}{0.49\textwidth}
Dal running a un loop:
\[
\frac{da}{dt}=-\pi b_0a^2+O(a^3).
\]
\end{column}
\end{columns}
\end{frame}

% 15
\begin{frame}{Derivata del termine al primo ordine}
\small
\begin{columns}[T,onlytextwidth]
\begin{column}{0.49\textwidth}
La correzione al primo ordine e'
\[
F_1=-2a=-\frac{2\alpha_s}{\pi}.
\]
Derivando rispetto a $t=\ln\bar\Lambda^2$:
\[
\frac{dF_1}{dt}
=-\frac2\pi\frac{d\alpha_s}{dt}.
\]
\end{column}
\begin{column}{0.49\textwidth}
Poiche'
\[
\frac{d\alpha_s}{dt}
=-b_0\alpha_s^2+O(\alpha_s^3),
\]
si ha
\[
\frac{dF_1}{dt}
=\frac{2b_0}{\pi}\alpha_s^2
+O(\alpha_s^3).
\]
La derivata ha \textbf{segno positivo}.
\end{column}
\end{columns}
\end{frame}

% 16
\begin{frame}{Derivata del logaritmo esplicito}
\small
\begin{columns}[T,onlytextwidth]
\begin{column}{0.49\textwidth}
Il termine esplicito e'
\[
F_{2,\ln}=-c\ln\frac{\bar\Lambda}{\mu_q}\,a^2.
\]
Derivando:
\[
\begin{aligned}
\frac{dF_{2,\ln}}{dt}
={}&-c\frac{d\ln(\bar\Lambda/\mu_q)}{dt}a^2\\
&-c\ln\frac{\bar\Lambda}{\mu_q}\frac{da^2}{dt}.
\end{aligned}
\]
\end{column}
\begin{column}{0.49\textwidth}
Usiamo
\[
\frac{d\ln(\bar\Lambda/\mu_q)}{dt}=\frac12
\]
e
\[
\frac{da^2}{dt}=2a\frac{da}{dt}=O(a^3).
\]
Fino a $O(a^2)$:
\[
\frac{dF_{2,\ln}}{dt}
=-\frac c2a^2+O(a^3).
\]
\end{column}
\end{columns}
\end{frame}

% 17
\begin{frame}{Le altre derivate di ordine $\alpha_s^2$}
\small
\begin{columns}[T,onlytextwidth]
\begin{column}{0.49\textwidth}
Il termine senza logaritmi:
\[
F_{2,G}=-Ga^2.
\]
Poiche' $da/dt=O(a^2)$,
\[
\frac{dF_{2,G}}{dt}
=-2Ga\frac{da}{dt}=O(a^3).
\]
Non contribuisce all'ordine $a^2$.
\end{column}
\begin{column}{0.49\textwidth}
Per il termine $-N_fa^2\ln a$:
\[
\begin{aligned}
\frac{d}{dt}(a^2\ln a)
&=(2a\ln a+a)\frac{da}{dt}\\
&=O(a^3\ln a).
\end{aligned}
\]
Anche questo contributo e' di ordine superiore
rispetto a $a^2$.
\end{column}
\end{columns}
\end{frame}

% 18
\begin{frame}{Cancellazione esplicita della dipendenza dalla scala}
\small
Sommiamo i contributi fino all'ordine $\alpha_s^2$:
\[
\begin{aligned}
\frac{d(P/P_{\rm SB})}{dt}
&=\frac{2b_0}{\pi}\alpha_s^2
-\frac c2\left(\frac{\alpha_s}{\pi}\right)^2
+O(\alpha_s^3\ln\alpha_s)\\
&=\left(2\pi b_0-\frac c2\right)
\left(\frac{\alpha_s}{\pi}\right)^2
+O(\alpha_s^3\ln\alpha_s).
\end{aligned}
\]
Ma $c=4\pi b_0$. Pertanto
\[
2\pi b_0-\frac c2
=2\pi b_0-2\pi b_0=0
\]
e, all'ordine considerato,
\[
\frac{d(P/P_{\rm SB})}{d\ln\bar\Lambda^2}
=O(\alpha_s^3\ln\alpha_s).
\]
\end{frame}

% 19
\begin{frame}{Nota sui segni nelle equazioni (8.46)--(8.48)}
\small
\begin{columns}[T,onlytextwidth]
\begin{column}{0.49\textwidth}
Dalla formula (8.44):
\[
-2\frac{\alpha_s}{\pi}
\quad\Longrightarrow\quad
+\frac{2b_0}{\pi}\alpha_s^2
\]
quando deriviamo rispetto a
$\ln\bar\Lambda^2$.
\medskip
Il running ha infatti segno negativo:
\[
\frac{d\alpha_s}{d\ln\bar\Lambda^2}
=-b_0\alpha_s^2.
\]
\end{column}
\begin{column}{0.49\textwidth}
Sempre dalla (8.44), il logaritmo esplicito
contribuisce
\[
-\frac12\left(11-\frac23N_f\right)
\left(\frac{\alpha_s}{\pi}\right)^2.
\]
I due contributi hanno segno opposto.
\medskip
\emph{Nella stampa allegata, i segni di
(8.46)--(8.48) appaiono invertiti; la
cancellazione finale (8.49) e' corretta.}
\end{column}
\end{columns}
\end{frame}

% 20
\begin{frame}{La banda di scala: che cosa misura?}
\small
\begin{columns}[T,onlytextwidth]
\begin{column}{0.49\textwidth}
Si ripete il calcolo di $P$ usando
\[
\bar\Lambda=\mu_q,
\quad2\mu_q,
\quad4\mu_q.
\]
Ogni scelta modifica il valore di
$\alpha_s(\bar\Lambda)$ e dei logaritmi
nella pressione troncata.
\medskip
La famiglia di risultati produce
una fascia, non una curva unica.
\end{column}
\begin{column}{0.49\textwidth}
\textbf{Banda stretta:}
limitata sensibilita' agli ordini omessi.
\medskip
\textbf{Banda larga:}
maggiore incertezza perturbativa;
possibile scarsa convergenza.
\medskip
La variazione di scala \emph{non} e' una stima rigorosa di tutti gli errori teorici.
\end{column}
\end{columns}
\end{frame}

\subsection{Figura 8.3: confronto fra pQCD e bag model}

% 21
\begin{frame}{Figura 8.3: pressione normalizzata a $T=0$}
\small
\begin{columns}[T,onlytextwidth]
\begin{column}{0.59\textwidth}
    \begin{figure}[!ht]\includegraphics[trim={0cm 0cm 0 0},clip,keepaspectratio,width=0.95\textwidth]{Schaffner/pressure_quarkgluon_SB}\label{fig:Schaffner_pressure_quarkgluon_SB}\end{figure}
\end{column}
\begin{column}{0.39\textwidth}
Asse orizzontale:
\[
\mu_B=3\mu_q.
\]
Asse verticale:
\[
P/P_{\rm SB}.
\]
\textbf{Curve:}
\begin{itemize}
\item gas libero: rapporto $1$;
\item bag model: linea tratteggiata;
\item pQCD: fascia grigia.
\end{itemize}
\end{column}
\end{columns}
\medskip
La fascia pQCD deriva dalla variazione di $\bar\Lambda$
intorno alla scelta $2\mu_q$.
\end{frame}

% 22
\begin{frame}{Lettura quantitativa della Figura 8.3}
\small
\begin{columns}[T,onlytextwidth]
\begin{column}{0.49\textwidth}
Per $\mu_B\simeq6\,\mathrm{GeV}$:
\[
P_{\rm pQCD}/P_{\rm SB}\simeq0.8.
\]
Riduzione di circa $20\%$; fascia relativamente stretta.
\medskip
Per $\mu_B\simeq3\,\mathrm{GeV}$:
\[
P_{\rm pQCD}/P_{\rm SB}
\simeq0.6\text{--}0.8.
\]
Riduzione del $20\text{--}40\%$.
\end{column}
\begin{column}{0.49\textwidth}
A $\mu_B$ piu' basso la fascia
si allarga molto e la pressione
estrapolata puo' arrivare a zero.
\medskip
La pQCD, per il calcolo mostrato,
e' ritenuta relativamente controllata
solo intorno a
\[
\mu_B\gtrsim3\,\mathrm{GeV}.
\]
Nelle stelle compatte si incontrano
anche $\mu_B\sim1\text{--}1.5\,\mathrm{GeV}$.
\end{column}
\end{columns}
\end{frame}

% 23
\begin{frame}{Perche' il punto $P=0$ interessa la stabilita'?}
\small
\begin{columns}[T,onlytextwidth]
\begin{column}{0.49\textwidth}
Per materia fredda neutra:
\[
\varepsilon+P=\mu_B n_B.
\]
Alla superficie autolegata, $P=0$:
\[
\frac EA=\frac{\varepsilon}{n_B}=\mu_B.
\]
Una materia $uds$ in bulk sarebbe
assolutamente stabile rispetto ai nuclei
se, indicativamente,
\[
\mu_B^{(0)}<930\,\mathrm{MeV}.
\]
\end{column}
\begin{column}{0.49\textwidth}
La Figura 8.3 comprende estrapolazioni
pQCD fino a piccole $\mu_B$.
\medskip
Ma proprio qui l'espansione non e'
controllata: il possibile zero di $P$
non dimostra la stabilita' della materia strange.
\medskip
La forma pQCD puo' essere usata
come modello efficace: non come predizione
sicura alle densita' stellari.
\end{column}
\end{columns}
\end{frame}

% 24
\begin{frame}{Il bag model si avvicina troppo rapidamente al gas libero}
\small
\begin{columns}[T,onlytextwidth]
\begin{column}{0.49\textwidth}
Nel modello MIT:
\[
P_{\rm MIT}=P_{\rm SB}-B.
\]
Dividendo per $P_{\rm SB}\propto\mu_q^4$:
\[
\frac{P_{\rm MIT}}{P_{\rm SB}}
=1-\frac{B}{P_{\rm SB}}
=1-O(\mu_q^{-4}).
\]
La correzione relativa del bag
si annulla come $\mu_q^{-4}$.
\end{column}
\begin{column}{0.49\textwidth}
La pQCD, invece, ha
\[
\frac{P_{\rm pQCD}}{P_{\rm SB}}
=1-\frac{2\alpha_s(\bar\Lambda)}{\pi}+\cdots.
\]
La libert\`a asintotica fa diminuire
$\alpha_s$ solo logaritmicamente
alle alte scale.
\medskip
Perci\`o la riduzione relativa della pressione
pQCD rimane apprezzabile piu' a lungo.
\end{column}
\end{columns}
\end{frame}

% 25
\begin{frame}{Bag model: esempio con $B^{1/4}=150\,\mathrm{MeV}$}
\small
\begin{columns}[T,onlytextwidth]
\begin{column}{0.49\textwidth}
La pressione e' nulla quando
\[
\frac{3\mu_q^4}{4\pi^2}=B.
\]
Quindi
\[
\mu_B^{(0)}
=3\left(\frac{4\pi^2B}{3}\right)^{1/4}
\simeq0.857\,\mathrm{GeV}.
\]
\end{column}
\begin{column}{0.49\textwidth}
Possiamo allora scrivere
\[
\frac{P_{\rm MIT}}{P_{\rm SB}}
=1-\left(\frac{\mu_B^{(0)}}{\mu_B}\right)^4.
\]
A $\mu_B=2\,\mathrm{GeV}$:
\[
\frac{P_{\rm MIT}}{P_{\rm SB}}
\simeq1-\left(\frac{0.857}{2}\right)^4
\simeq0.966.
\]
Riduzione di circa $3.4\%$ soltanto.
\end{column}
\end{columns}
\end{frame}

\subsection{Modello efficace e identita' termodinamiche}

% 26
\begin{frame}{Un MIT bag model corretto dalle interazioni}
\small
Il libro introduce la parametrizzazione efficace
\[
P(\mu_q)=\frac{N_f}{4\pi^2}a_{\rm eff}\mu_q^4-B_{\rm eff}.
\]
\begin{columns}[T,onlytextwidth]
\begin{column}{0.49\textwidth}
\textbf{$a_{\rm eff}$}

Fattore adimensionale che riduce
il coefficiente della pressione ideale:
\[
a_{\rm eff}=1\quad\text{(gas libero)}.
\]
Valore illustrativo del testo:
\[
a_{\rm eff}\simeq0.63.
\]
\end{column}
\begin{column}{0.49\textwidth}
\textbf{$B_{\rm eff}$}

Contributo energetico costante,
analogo al bag originale.
\medskip
Dimensionalmente
\[
[B_{\rm eff}]={\rm energia}^4.
\]
Se si cita un valore in MeV, si intende
$B_{\rm eff}^{1/4}$ (per esempio $200$ MeV).
\end{column}
\end{columns}
\end{frame}

% 27
\begin{frame}{Come interpretare $a_{\rm eff}$?}
\small
\begin{columns}[T,onlytextwidth]
\begin{column}{0.49\textwidth}
Al primo ordine perturbativo:
\[
P_{\rm pQCD}\simeq P_{\rm SB}
\left(1-\frac{2\alpha_s}{\pi}\right).
\]
Il modello efficace usa:
\[
P_{\rm eff}\simeq a_{\rm eff}P_{\rm SB}
\quad(B_{\rm eff}\text{ trascurato}).
\]
Quindi, qualitativamente,
\[
a_{\rm eff}\sim1-\frac{2\alpha_s}{\pi}.
\]
\end{column}
\begin{column}{0.49\textwidth}
Ma $a_{\rm eff}$ e' assunto costante.
\medskip
Nella pQCD $\alpha_s$ dipende dalla scala
e compaiono ordini superiori.
\medskip
Dunque $a_{\rm eff}$ e' un parametro
fenomenologico, non una misura diretta
del numero di colori, spin o sapori.
\end{column}
\end{columns}
\end{frame}

% 28
\begin{frame}{Identita' termodinamiche a $T=0$}
\small
\begin{columns}[T,onlytextwidth]
\begin{column}{0.49\textwidth}
Per un sistema omogeneo con un solo
potenziale chimico indipendente $\mu_q$:
\[
dP=n_q\,d\mu_q
\quad(T=0).
\]
Pertanto
\[
n_q=\frac{\partial P}{\partial\mu_q}.
\]
$n_q$ e' la densita' totale del numero di quark.
\end{column}
\begin{column}{0.49\textwidth}
Dalla trasformazione di Legendre:
\[
\varepsilon=-P+\mu_qn_q.
\]
Ogni quark trasporta $1/3$ di numero barionico:
\[
n_B=\frac{n_q}{3},
\qquad \mu_B=3\mu_q.
\]
Quindi $\mu_qn_q=\mu_Bn_B$.
\end{column}
\end{columns}
\medskip
Le due relazioni sono indispensabili per passare da $P(\mu_q)$ a $P(\varepsilon)$.
\end{frame}

% 29
\begin{frame}{Densita' dei quark dal modello efficace}
\small
Partiamo da
\[
P(\mu_q)=\frac{N_fa_{\rm eff}}{4\pi^2}\mu_q^4-B_{\rm eff},
\]
con $a_{\rm eff}$ e $B_{\rm eff}$ \emph{costanti}.
\begin{columns}[T,onlytextwidth]
\begin{column}{0.49\textwidth}
Deriviamo:
\[
\begin{aligned}
n_q&=\frac{\partial P}{\partial\mu_q}\\
&=\frac{N_fa_{\rm eff}}{4\pi^2}\,4\mu_q^3\\
&=\frac{N_fa_{\rm eff}}{\pi^2}\mu_q^3.
\end{aligned}
\]
\end{column}
\begin{column}{0.49\textwidth}
Passando alla densita' barionica:
\[
\begin{aligned}
n_B&=\frac{n_q}{3}\\
&=\frac{N_fa_{\rm eff}}{3\pi^2}\mu_q^3.
\end{aligned}
\]
Il parametro $a_{\rm eff}$ rimane
in $n_B(\mu_q)$.
\end{column}
\end{columns}
\end{frame}

% 30
\begin{frame}{Densita' di energia dal modello efficace}
\small
Usiamo $\varepsilon=-P+\mu_qn_q$:
\[
\begin{aligned}
\varepsilon
&=-\left(\frac{N_fa_{\rm eff}}{4\pi^2}\mu_q^4-B_{\rm eff}\right)
+\mu_q\frac{N_fa_{\rm eff}}{\pi^2}\mu_q^3\\
&=\left(-\frac14+1\right)
\frac{N_fa_{\rm eff}}{\pi^2}\mu_q^4+B_{\rm eff}\\
&=\frac{3N_fa_{\rm eff}}{4\pi^2}\mu_q^4+B_{\rm eff}.
\end{aligned}
\]
Il contributo del gas alla densita' di energia e' tre volte
quello alla pressione, per $a_{\rm eff}$ costante.
\end{frame}

% 31
\begin{frame}{Eliminare $\mu_q$: derivazione di $P(\varepsilon)$}
\small
\begin{columns}[T,onlytextwidth]
\begin{column}{0.49\textwidth}
Abbiamo
\[
P=\frac{N_fa_{\rm eff}}{4\pi^2}\mu_q^4-B_{\rm eff},
\]
\[
\varepsilon=
\frac{3N_fa_{\rm eff}}{4\pi^2}\mu_q^4+B_{\rm eff}.
\]
Dalla prima formula:
\[
\frac{N_fa_{\rm eff}}{4\pi^2}\mu_q^4
=P+B_{\rm eff}.
\]
\end{column}
\begin{column}{0.49\textwidth}
Sostituiamo nella seconda:
\[
\begin{aligned}
\varepsilon
&=3(P+B_{\rm eff})+B_{\rm eff}\\
&=3P+4B_{\rm eff}.
\end{aligned}
\]
Ricaviamo infine
\[
P(\varepsilon)
=\frac13(\varepsilon-4B_{\rm eff}).
\]
\end{column}
\end{columns}
\medskip
La forma e' identica alla EOS del MIT bag model con $B\to B_{\rm eff}$.
\end{frame}

% 32
\begin{frame}{Perche' $a_{\rm eff}$ scompare dall'EOS?}
\small
\begin{columns}[T,onlytextwidth]
\begin{column}{0.49\textwidth}
Indichiamo
\[
A=\frac{N_fa_{\rm eff}}{4\pi^2}.
\]
Allora
\[
P=A\mu_q^4-B_{\rm eff},
\]
\[
\varepsilon=3A\mu_q^4+B_{\rm eff}.
\]
Il fattore $A$ e' comune.
\end{column}
\begin{column}{0.49\textwidth}
Sottraendo $3P$ da $\varepsilon$:
\[
\varepsilon-3P=4B_{\rm eff}.
\]
Quindi $a_{\rm eff}$ si cancella.
\medskip
La velocita' del suono e'
\[
c_s^2=\frac{dP}{d\varepsilon}=\frac13.
\]
\end{column}
\end{columns}
\medskip
Questa cancellazione \emph{richiede} che $a_{\rm eff}$ sia costante.
\end{frame}

% 33
\begin{frame}{Il parametro $a_{\rm eff}$ resta fisicamente importante}
\small
\begin{columns}[T,onlytextwidth]
\begin{column}{0.49\textwidth}
La densita' barionica dipende da $a_{\rm eff}$:
\[
n_B=\frac{N_fa_{\rm eff}}{3\pi^2}\mu_q^3.
\]
Due modelli con lo stesso $P(\varepsilon)$
possono associare densita' barioniche
diverse alla medesima pressione.
\end{column}
\begin{column}{0.49\textwidth}
Imponendo $P=0$:
\[
\mu_q^{(0)}=
\left(\frac{4\pi^2B_{\rm eff}}
{N_fa_{\rm eff}}\right)^{1/4},
\]
\[
\mu_B^{(0)}=3\mu_q^{(0)}.
\]
A $B_{\rm eff}$ fisso, diminuire
$a_{\rm eff}$ aumenta $\mu_B^{(0)}$.
\end{column}
\end{columns}
\medskip
Si modificano i vincoli di stabilita' $E/A=\mu_B^{(0)}$
del \S8.3, anche se l'EOS $P(\varepsilon)$ non cambia.
\end{frame}

% 34
\begin{frame}{Se il coupling corre, il coefficiente non e' costante}
\small
Consideriamo una parametrizzazione piu' generale:
\[
P(\mu_q)=C\mu_q^4 f(\mu_q)-B,
\]
con $C$ costante e $f(\mu_q)$ fattore dipendente dalle interazioni.
\begin{columns}[T,onlytextwidth]
\begin{column}{0.49\textwidth}
Derivando:
\[
\begin{aligned}
n_q&=\frac{dP}{d\mu_q}\\
&=C\left[4\mu_q^3f(\mu_q)
+\mu_q^4f'(\mu_q)\right].
\end{aligned}
\]
\end{column}
\begin{column}{0.49\textwidth}
Il nuovo contributo
\[
C\mu_q^4f'(\mu_q)
\]
compare perche' il coefficiente
cambia con $\mu_q$.
\medskip
Nella QCD questo accade quando
$\bar\Lambda\propto\mu_q$ e
$\alpha_s=\alpha_s(\bar\Lambda)$.
\end{column}
\end{columns}
\end{frame}

% 35
\begin{frame}{Conseguenza del running per $P(\varepsilon)$}
\small
\begin{columns}[T,onlytextwidth]
\begin{column}{0.49\textwidth}
Dalla relazione $\varepsilon=-P+\mu_qn_q$:
\[
\begin{aligned}
\varepsilon
={}&-C\mu_q^4f+B\\
&+C\left(4\mu_q^4f+\mu_q^5f'\right)\\
={}&3C\mu_q^4f+C\mu_q^5f'+B.
\end{aligned}
\]
\end{column}
\begin{column}{0.49\textwidth}
Poiche'
\[
3P=3C\mu_q^4f-3B,
\]
segue
\[
\varepsilon-3P=4B+C\mu_q^5f'.
\]
Quindi
\[
P=\frac{\varepsilon-4B}{3}
-\frac{C\mu_q^5}{3}f'(\mu_q).
\]
\end{column}
\end{columns}
\medskip
Se $f'=0$ si recupera esattamente il bag model;
se $f'\ne0$, l'EOS $P(\varepsilon)$ cambia.
\end{frame}


\section{Relazioni massa--raggio delle stelle di quark}

\subsection{EOS autolegata e regime di piccola massa}

% 01
\begin{frame}{Dal \S8.4 alla struttura delle stelle di quark}
\small
\begin{columns}[T,onlytextwidth]
\begin{column}{0.49\textwidth}
\textbf{Risultato del \S8.4}
\[
P(\varepsilon)=\frac13(\varepsilon-4B_{\rm eff}).
\]
La forma dell'EOS coincide con quella del MIT bag model, se il parametro efficace $a_{\rm eff}$ e' costante.
\medskip
Al bordo di una stella \emph{senza crosta}:
\[
P(R)=0 \quad\Longrightarrow\quad
\varepsilon_s=4B_{\rm eff}>0.
\]
\end{column}
\begin{column}{0.49\textwidth}
\textbf{Obiettivi del \S8.5}
\begin{itemize}
\item Ricavare il comportamento $M(R)$ delle configurazioni autolegate.
\item Comprendere l'effetto di una crosta nucleare.
\item Interpretare i limiti relativistici di compattezza.
\item Confrontare stelle strange e stelle di neutroni.
\end{itemize}
\end{column}
\end{columns}
\end{frame}

% 02
\begin{frame}{Il segno dell'esponente nel libro: un refuso}
\small
\begin{columns}[T,onlytextwidth]
\begin{column}{0.49\textwidth}
Nei \S\S8.3 e 8.5 compare la scrittura
\[
M\propto R^{-3}.
\]
Ma il testo del \S8.5 afferma anche
\[
\frac{M}{R^3}=\text{costante}.
\]
Le due affermazioni non sono compatibili.
\end{column}
\begin{column}{0.49\textwidth}
Per densita' media costante:
\[
M=\rho\,V=\rho\,\frac{4\pi}{3}R^3.
\]
Segue necessariamente
\[
M\propto R^3, \qquad \frac{M}{R^3}=\text{costante}.
\]
L'esponente $-3$ e' quindi un refuso ripetuto; la derivazione fisica del testo resta corretta.
\end{column}
\end{columns}
\end{frame}

% 03
\begin{frame}{Stella autolegata leggera: densita' quasi uniforme}
\small
\begin{columns}[T,onlytextwidth]
\begin{column}{0.49\textwidth}
Se la massa e' piccola, la gravita' comprime poco il fluido e
\[
P(r)\ll B_{\rm eff}.
\]
Dall'EOS ricaviamo
\[
\varepsilon(r)=3P(r)+4B_{\rm eff}
\simeq4B_{\rm eff}.
\]
Passando alla densita' di massa:
\[
\rho(r)=\frac{\varepsilon(r)}{c^2}
\simeq\frac{4B_{\rm eff}}{c^2}.
\]
\end{column}
\begin{column}{0.49\textwidth}
L'equazione di massa e'
\[
\frac{dm}{dr}=4\pi r^2\rho(r).
\]
Integrando fino a $R$:
\[
\begin{aligned}
M&=4\pi\int_0^R\rho(r)r^2\,dr\\
&\simeq\frac{16\pi B_{\rm eff}}{c^2}
\int_0^R r^2\,dr\\
&=\frac{16\pi B_{\rm eff}}{3c^2}R^3.
\end{aligned}
\]
\end{column}
\end{columns}
\end{frame}

% 04
\begin{frame}{La legge massa--raggio del ramo autolegato}
\small
\[
\frac{M}{R^3}\simeq\frac{16\pi B_{\rm eff}}{3c^2}
=\text{costante},\qquad M\propto R^3.
\]
\begin{columns}[T,onlytextwidth]
\begin{column}{0.49\textwidth}
\textbf{Esempio a densita' fissata}
\[
R' = \frac{R}{2}
\quad\Longrightarrow\quad
M' = \frac{M}{8}.
\]
La massa diminuisce come il volume, non cresce con $R^{-3}$.
\end{column}
\begin{column}{0.49\textwidth}
\textbf{Limite di piccola massa}
\[
M\to0 \quad\Longrightarrow\quad R\to0.
\]
La sequenza macroscopica delle stelle strange senza crosta puo' iniziare nell'origine del diagramma $M$--$R$.
\end{column}
\end{columns}
\medskip
La legge non e' esatta per stelle massicce: la pressione e la densita' diventano non uniformi.
\end{frame}

% 05
\begin{frame}{Autolegamento e dimensioni microscopiche}
\small
\begin{columns}[T,onlytextwidth]
\begin{column}{0.49\textwidth}
\textbf{Stelle strange macroscopiche}
\[
P(R)=0,\qquad \varepsilon_s=4B_{\rm eff}>0.
\]
Il fluido puo' restare aggregato senza che la gravita' fornisca il meccanismo primario di legame.
\medskip
L'EOS e' una descrizione continua di materia \emph{bulk}.
\end{column}
\begin{column}{0.49\textwidth}
\textbf{Strangelet microscopici}
\[
E(A)\simeq A e_{\rm bulk}+E_{\rm surf}
+E_{\rm curv}+E_{\rm Coul}.
\]
Quando $A$ e' piccolo, superficie, curvatura e Coulomb contano molto.
\medskip
Quindi $M\to0$ lungo la curva macroscopica non dimostra la stabilita' di una goccia con pochi quark.
\end{column}
\end{columns}
\end{frame}

\subsection{Equazioni TOV e massa massima}

% 06
\begin{frame}{Equilibrio idrostatico relativistico: equazioni TOV}
\small
In un fluido statico, sferico e perfetto, con $G=c=1$:
\[
\frac{dm}{dr}=4\pi r^2\varepsilon(r),
\qquad
\frac{dP}{dr}=-\frac{[\varepsilon(r)+P(r)]
[m(r)+4\pi r^3P(r)]}{r[r-2m(r)]}.
\]
\begin{columns}[T,onlytextwidth]
\begin{column}{0.49\textwidth}
Dati $P(0)=P_c$ e $m(0)=0$, si integrano le TOV con l'EOS:
\[
P=\frac{\varepsilon-4B_{\rm eff}}{3}.
\]
\end{column}
\begin{column}{0.49\textwidth}
La superficie del \emph{core nudo} e' definita da $P(R)=0$.
\[
M=m(R),\qquad \varepsilon(R)=4B_{\rm eff}.
\]
Variando $P_c$ si costruisce una sequenza $M(R)$.
\end{column}
\end{columns}
\end{frame}

% 07
\begin{frame}{Perche' la massa massima e' finita?}
\small
\begin{columns}[T,onlytextwidth]
\begin{column}{0.49\textwidth}
\textbf{Regime di piccola massa}
\begin{itemize}
\item $P_c$ modesta rispetto a $B_{\rm eff}$.
\item $\varepsilon(r)\simeq4B_{\rm eff}$.
\item $M\propto R^3$.
\end{itemize}
\end{column}
\begin{column}{0.49\textwidth}
\textbf{Regime di grande massa}
\begin{itemize}
\item La gravita' aumenta $P_c$.
\item La densita' varia significativamente con $r$.
\item La curva TOV raggiunge un massimo di massa lungo il ramo stabile.
\end{itemize}
\end{column}
\end{columns}
\medskip
L'EOS ha velocita' del suono costante:
\[
\frac{c_s^2}{c^2}=\frac{dP}{d\varepsilon}=\frac13.
\]
La pressione non cresce abbastanza rapidamente da sostenere configurazioni stabili arbitrariamente massicce.
\end{frame}

% 08
\begin{frame}{Come la costante del bag fissa la scala della stella}
\small
In unita' geometriche, definiamo variabili adimensionali
\[
\varepsilon=B_{\rm eff}\widetilde\varepsilon,
\quad P=B_{\rm eff}\widetilde P,
\quad r=B_{\rm eff}^{-1/2}\widetilde r,
\quad m=B_{\rm eff}^{-1/2}\widetilde m.
\]
\begin{columns}[T,onlytextwidth]
\begin{column}{0.49\textwidth}
Dall'EOS:
\[
\widetilde P=\frac{\widetilde\varepsilon-4}{3}.
\]
Le TOV in queste variabili non contengono piu' $B_{\rm eff}$.
\end{column}
\begin{column}{0.49\textwidth}
Configurazioni corrispondenti hanno percio' scale
\[
M,R\propto B_{\rm eff}^{-1/2}.
\]
Posto $b=B_{\rm eff}^{1/4}$:
\[
M_{\max},R_{\max}\propto b^{-2}.
\]
\end{column}
\end{columns}
\end{frame}

% 09
\begin{frame}{Due leggi di scala, due confronti distinti}
\small
\begin{columns}[T,onlytextwidth]
\begin{column}{0.49\textwidth}
\textbf{Lungo una stessa sequenza}
\[
M\propto R^3.
\]
\begin{itemize}
\item $B_{\rm eff}$ fissato.
\item Cambia la massa / il numero di barioni.
\item Valido solo se la stella e' poco compressa.
\end{itemize}
\end{column}
\begin{column}{0.49\textwidth}
\textbf{Fra equazioni di stato diverse}
\[
M_{\max},R_{\max}\propto B_{\rm eff}^{-1/2}.
\]
\begin{itemize}
\item Si modifica $B_{\rm eff}$.
\item Si confrontano soluzioni TOV corrispondenti.
\item Vale come riscalamento dell'intera famiglia ideale.
\end{itemize}
\end{column}
\end{columns}
\medskip
Nel libro compare $B_{\rm eff}=145\,\mathrm{MeV}$: dimensionalmente va letto come $B_{\rm eff}^{1/4}=145\,\mathrm{MeV}$.
\end{frame}

\subsection{Figura 8.4: core strange e crosta nucleare}

% 11
\begin{frame}{Come leggere la Figura 8.4}
\small
\begin{columns}[T,onlytextwidth]
\begin{column}{0.49\textwidth}
\textbf{Asse e curve}
\begin{itemize}
\item Ascissa: raggio $R$ in km.
\item Ordinata: $M/M_\odot$.
\item \emph{No crust}: superficie direttamente di materia strange.
\item \emph{With crust}: un involucro nucleare modifica soprattutto il raggio.
\end{itemize}
La zona tratteggiata combina un limite di causalita' con la richiesta di sostenere configurazioni di circa $2M_\odot$.
\end{column}
\begin{column}{0.49\textwidth}
    \begin{figure}[!ht]\includegraphics[trim={0cm 0cm 0 0},clip,keepaspectratio,width=0.95\textwidth]{Schaffner/MR-quarkstar}\label{fig:Schaffner_MR-quarkstar}\end{figure}
    \textbf{Vincoli mostrati}
\begin{itemize}
\item Retta dell'orizzonte $\mathcal C=1/2$.
\item Retta del limite di Buchdahl $\mathcal C=4/9$.
\item Retta e curva di causalita'.
\item Massa pulsar vicina a $2M_\odot$.
\end{itemize}
\end{column}
\end{columns}
\end{frame}

% 12
\begin{frame}{Perche' una stella strange puo' avere una crosta?}
\small
\begin{columns}[T,onlytextwidth]
\begin{column}{0.49\textwidth}
\textbf{Core di quark}
\[
P(R)=0,\qquad \varepsilon_s=4B_{\rm eff}>0.
\]
La materia strange e' autolegata nel modello.
\medskip
Con $m_s\ne0$, un piccolo eccesso di carica positiva nel core e uno strato elettronico possono generare un campo elettrico superficiale.
\end{column}
\begin{column}{0.49\textwidth}
\textbf{Crosta nucleare}
\begin{itemize}
\item Puo' essere sostenuta elettricamente al di sopra della superficie dei quark.
\item E' formata da nuclei ed elettroni, non da materia autolegata di quark.
\item Il suo spessore dipende dalla gravita'.
\end{itemize}
\end{column}
\end{columns}
\medskip
La presenza di una crosta cambia la forma del ramo a bassa massa, pur lasciando simile la massa massima della sequenza mostrata.
\end{frame}

% 13
\begin{frame}{Crosta di una stella massiccia}
\small
\begin{columns}[T,onlytextwidth]
\begin{column}{0.49\textwidth}
\textbf{Core massiccio}
\begin{itemize}
\item La gravita' e' relativamente intensa.
\item La crosta subisce forte compressione.
\item Lo spessore radiale rimane piccolo.
\end{itemize}
\end{column}
\begin{column}{0.49\textwidth}
\textbf{Effetto sul diagramma $M$--$R$}
\begin{itemize}
\item Variazione modesta della massa.
\item Il raggio cresce di qualche centinaio di metri nel caso discusso nel libro.
\item Le masse massime dei due rami sono simili.
\end{itemize}
\end{column}
\end{columns}
\medskip
La crosta modifica soprattutto la posizione orizzontale della curva nelle vicinanze della massa massima.
\end{frame}

% 14
\begin{frame}{Stella leggera: come scala la gravita' superficiale?}
\small
Per una stella autolegata \emph{nuda} e quasi uniforme,
\[
M\simeq\frac{4\pi}{3}\rho_sR^3,
\qquad g\simeq\frac{GM}{R^2}.
\]
\begin{columns}[T,onlytextwidth]
\begin{column}{0.49\textwidth}
Sostituendo $M(R)$:
\[
\begin{aligned}
g&\simeq\frac{G}{R^2}\frac{4\pi}{3}\rho_sR^3\\
 &=\frac{4\pi G\rho_s}{3}R.
\end{aligned}
\]
Quindi $g\propto R$.
\end{column}
\begin{column}{0.49\textwidth}
Poiche' $R\propto M^{1/3}$,
\[
g\propto M^{1/3}.
\]
Diminuendo la massa, la gravita' superficiale del core diventa piu' debole.
\medskip
Questa stima vale nel regime quasi uniforme e non per una crosta molto estesa.
\end{column}
\end{columns}
\end{frame}

% 15
\begin{frame}{Equilibrio idrostatico e spessore della crosta}
\small
Nell'approssimazione newtoniana, per una crosta sottile:
\[
\frac{dP}{dr}\simeq-\rho(r)g,
\quad\Longrightarrow\quad
 dr\simeq-\frac{dP}{\rho(P)g}.
\]
Integrando dalla superficie ($P=0$) alla base ($P=P_b$):
\[
\Delta R\simeq\frac1g\int_0^{P_b}\frac{dP}{\rho(P)}.
\]
\begin{columns}[T,onlytextwidth]
\begin{column}{0.49\textwidth}
Per EOS della crosta e $P_b$ fissate:
\[
\Delta R\propto\frac1g.
\]
Una minore gravita' implica una crosta meno compressa e piu' spessa.
\end{column}
\begin{column}{0.49\textwidth}
\textbf{Limiti dell'approssimazione}
\begin{itemize}
\item $g$ quasi costante nello strato.
\item $\Delta R\ll R_{\rm core}$.
\item Poco accurata quando si forma un alone esteso.
\end{itemize}
\end{column}
\end{columns}
\end{frame}

% 16
\begin{frame}{Perche' il ramo con crosta si piega verso grandi raggi?}
\small
\begin{columns}[T,onlytextwidth]
\begin{column}{0.49\textwidth}
\textbf{Riducendo la massa}
\begin{itemize}
\item Il core autolegato si restringe.
\item La sua gravita' superficiale diminuisce.
\item La crosta si espande radialmente.
\end{itemize}
\end{column}
\begin{column}{0.49\textwidth}
\textbf{Conseguenza per $M(R)$}
Il raggio \emph{totale} puo' aumentare quando la massa diminuisce, anche se il raggio del core strange diminuisce.
\medskip
Il testo cita $M\sim0.1M_\odot$ e $R\sim7\,\mathrm{km}$ come regime in cui la crosta modifica fortemente il ramo.
\end{column}
\end{columns}
\medskip
L'andamento $M\propto R^3$ si riferisce al core nudo nel regime quasi uniforme; non si applica al raggio totale di una configurazione con crosta estesa.
\end{frame}

\subsection{Compattezza, causalita' e vincoli osservativi}

% 17
\begin{frame}{Compattezza: come leggere le rette nel piano $M$--$R$}
\small
Definiamo la compattezza adimensionale:
\[
\mathcal C=\frac{GM}{Rc^2}.
\]
A compattezza fissata:
\[
M=\frac{\mathcal C c^2}{G}R.
\]
\begin{columns}[T,onlytextwidth]
\begin{column}{0.49\textwidth}
Nel diagramma massa--raggio, le curve a $\mathcal C$ costante sono \textbf{rette passanti per l'origine}.
\medskip
A massa fissata, un raggio minore comporta compattezza maggiore.
\end{column}
\begin{column}{0.49\textwidth}
Le tre rette della figura rappresentano limiti sempre piu' restrittivi:
\[
\frac12,\qquad\frac49,\qquad0.354.
\]
Sono associati rispettivamente a orizzonte, Buchdahl e causalita' (nel contesto precisato dal testo).
\end{column}
\end{columns}
\end{frame}

% 18
\begin{frame}{Limite del buco nero: $\mathcal C=1/2$}
\small
\begin{columns}[T,onlytextwidth]
\begin{column}{0.49\textwidth}
Il raggio di Schwarzschild e'
\[
R_S=\frac{2GM}{c^2}.
\]
Pertanto
\[
\mathcal C_S=\frac{GM}{R_Sc^2}
=\frac12.
\]
\end{column}
\begin{column}{0.49\textwidth}
Una stella statica con superficie materiale al di fuori dell'orizzonte richiede
\[
R>\frac{2GM}{c^2},
\qquad \mathcal C<\frac12.
\]
La parte a sinistra della corrispondente retta non ospita stelle statiche ordinarie.
\end{column}
\end{columns}
\end{frame}

% 19
\begin{frame}{Limite di Buchdahl: $\mathcal C\leq4/9$}
\small
\begin{columns}[T,onlytextwidth]
\begin{column}{0.49\textwidth}
Sotto le ipotesi del teorema (fra cui simmetria sferica, staticita', pressione isotropa e densita' non crescente verso l'esterno):
\[
R\geq\frac94\frac{GM}{c^2}.
\]
Di conseguenza
\[
\mathcal C=\frac{GM}{Rc^2}\leq\frac49.
\]
\end{column}
\begin{column}{0.49\textwidth}
Il raggio minimo e' ora
\[
R_{\min}=2.25\frac{GM}{c^2},
\]
maggiore del raggio di Schwarzschild.
\medskip
Non e' un vincolo privo di ipotesi: stress anisotropi o altre condizioni possono modificarlo.
\end{column}
\end{columns}
\end{frame}

% 20
\begin{frame}{Limite di causalita': $\mathcal C\simeq0.354$}
\small
\begin{columns}[T,onlytextwidth]
\begin{column}{0.49\textwidth}
\textbf{Condizione locale}
\[
\frac{c_s^2}{c^2}=\frac{dP}{d\varepsilon}\leq1.
\]
La velocita' di propagazione del suono non puo' superare $c$.
\medskip
Per l'EOS del bag model:
\[
\frac{c_s^2}{c^2}=\frac13.
\]
\end{column}
\begin{column}{0.49\textwidth}
\textbf{Vincolo globale}
Per la classe di EOS e configurazioni stabili discussa nel testo, la massima compattezza causale e' circa
\[
\mathcal C_{\max}\simeq0.354.
\]
Da cui
\[
R\gtrsim\frac{1}{0.354}\frac{GM}{c^2}.
\]
Il numero $0.354$ non si ottiene imponendo localmente $dP/d\varepsilon=1$ in un singolo punto.
\end{column}
\end{columns}
\end{frame}

% 21
\begin{frame}{Raggi minimi a massa fissata: confronto}
\small
Dalla definizione $\mathcal C=GM/(Rc^2)$:
\[
R_{\min}(\mathcal C)=\frac{1}{\mathcal C}\frac{GM}{c^2}.
\]
\medskip
\begin{center}
\begin{tabular}{lll}
\hline
Limite & $\mathcal C$ & $R_{\min}$ \\
\hline
Schwarzschild & $1/2$ & $2GM/c^2$ \\
Buchdahl & $4/9$ & $2.25GM/c^2$ \\
Causalita' del testo & $0.354$ & $2.825GM/c^2$ \\
\hline
\end{tabular}
\end{center}
\medskip
Per $M=2M_\odot$, usando $GM_\odot/c^2\simeq1.477\,\mathrm{km}$:
\[
R_{\min,\rm caus}\simeq\frac{2\times1.477\,\mathrm{km}}{0.354}
\simeq8.3\,\mathrm{km}.
\]
\end{frame}

% 22
\begin{frame}{Vincolo osservativo: stelle con $M\simeq2M_\odot$}
\small
\begin{columns}[T,onlytextwidth]
\begin{column}{0.49\textwidth}
La Figura 8.4 indica la massa del pulsar PSR J0348+0432, vicina a
\[
M\simeq2M_\odot.
\]
Qualunque EOS proposta per l'oggetto deve permettere almeno
\[
M_{\max}\geq2M_\odot.
\]
\end{column}
\begin{column}{0.49\textwidth}
Questa e' una restrizione distinta dalla causalita':
\begin{itemize}
\item \emph{causalita'}: limita la risposta locale del fluido;
\item \emph{massa osservata}: richiede un ramo stabile sufficientemente massiccio.
\end{itemize}
Le due richieste limitano congiuntamente il piano $M$--$R$.
\end{column}
\end{columns}
\end{frame}

% 23
\begin{frame}{Da dove viene la regione tratteggiata esclusa?}
\small
\begin{columns}[T,onlytextwidth]
\begin{column}{0.49\textwidth}
Si considera un'EOS limite, con
\[
\frac{dP}{d\varepsilon}=1,
\]
ottimizzata in modo da sostenere una massa massima di circa $2M_\odot$.
\medskip
La sua curva $M(R)$ delimita una regione aggiuntiva del grafico, oltre alla sola retta di compattezza $0.354$.
\end{column}
\begin{column}{0.49\textwidth}
\textbf{Interpretazione}
\begin{itemize}
\item La linea di causalita' esclude le configurazioni troppo compatte.
\item La richiesta $M_{\max}\geq2M_\odot$ restringe ulteriormente le curve ammissibili.
\item La zona tratteggiata indica la regione esclusa nel confronto del testo.
\end{itemize}
\end{column}
\end{columns}
\medskip
Scoprire pulsar ancora piu' massicce rafforzerebbe questo tipo di vincolo.
\end{frame}

\subsection{Confronto con le stelle di neutroni}

% 24
\begin{frame}{Tabella 8.2: stelle autolegate e stelle di neutroni}
\footnotesize
\setlength{\tabcolsep}{4pt}
\begin{center}
\begin{tabular}{p{0.19\textwidth}p{0.36\textwidth}p{0.36\textwidth}}
\hline
 & \textbf{Stelle autolegate} & \textbf{Stelle di neutroni} \\
\hline
EOS & $P=0$ a $\varepsilon_s>0$ & $P=0$ a $\varepsilon\to0$ (modello ideale) \\
Legame & Interazioni microscopiche & Attrazione gravitazionale \\
Piccole masse & Curva che puo' iniziare dall'origine (senza crosta) & Raggi grandi per masse piccole \\
Massa minima & Arbitrariamente piccola nel modello macroscopico & $M_{\min}\sim0.1M_\odot$ nel modello citato \\
\hline
\end{tabular}
\end{center}
\medskip
L'andamento per stelle strane \emph{con crosta} puo' differire da quello del core autolegato nudo.
\end{frame}

% 25
\begin{frame}{Perche' le stelle di neutroni hanno una massa minima?}
\small
\begin{columns}[T,onlytextwidth]
\begin{column}{0.49\textwidth}
\textbf{Stella di neutroni ordinaria}
\begin{itemize}
\item La materia fredda neutra in equilibrio beta non ha lo stesso minimo autolegato del core strange.
\item Riducendo la massa, si indebolisce il confinamento gravitazionale.
\item Il ramo stabile termina a una massa minima.
\end{itemize}
\end{column}
\begin{column}{0.49\textwidth}
Nel modello citato nel libro:
\[
M_{\min}\simeq0.1M_\odot,
\qquad R\simeq200\,\mathrm{km}.
\]
Non e' un valore universale: dipende da EOS, condizioni termiche ed equilibrio considerato.
\end{column}
\end{columns}
\medskip
La stella autolegata \emph{nuda}, invece, non richiede una massa macroscopica minima perche' il materiale e' legato internamente.
\end{frame}

% 26
\begin{frame}{Che cosa rivela una misura di massa e raggio?}
\small
\begin{columns}[T,onlytextwidth]
\begin{column}{0.49\textwidth}
\textbf{Possibile indizio di autolegamento}
\begin{itemize}
\item Una configurazione con massa molto piccola e raggio molto piccolo.
\item Una sequenza che segue $M\propto R^3$ alle piccole masse, se il core e' nudo.
\end{itemize}
\end{column}
\begin{column}{0.49\textwidth}
\textbf{Ambiguita' osservativa}
\begin{itemize}
\item A masse stellari tipiche, diverse EOS possono dare $M$ e $R$ simili.
\item Una crosta modifica il ramo leggero.
\item Una misura isolata di $M$ e $R$ non determina univocamente la composizione.
\end{itemize}
\end{column}
\end{columns}
\end{frame}

% 27
\begin{frame}{Connessione fra i \S\S8.2--8.5}
\footnotesize
\begin{columns}[T,onlytextwidth]
\begin{column}{0.49\textwidth}
\textbf{\S8.2: quark liberi}
\[
P=\frac{\varepsilon}{3}.
\]
\textbf{\S8.3: bag e autolegamento}
\[
P=\frac{\varepsilon-4B}{3}.
\]
Un vuoto con energia positiva consente una superficie a densita' finita.
\end{column}
\begin{column}{0.49\textwidth}
\textbf{\S8.4: interazioni efficaci}
\[
P(\mu_q)=\frac{N_fa_{\rm eff}}{4\pi^2}\mu_q^4
-B_{\rm eff}.
\]
\textbf{\S8.5: struttura stellare}
\[
P(\varepsilon)\quad\longrightarrow\quad
\text{TOV}\quad\longrightarrow\quad M(R).
\]
Si aggiungono gli effetti della crosta e i vincoli relativistici.
\end{column}
\end{columns}
\end{frame}

% 28
\begin{frame}{\S8.5: risultati da conservare}
\small
\begin{columns}[T,onlytextwidth]
\begin{column}{0.49\textwidth}
\textbf{Stelle autolegate nude, piccole masse}
\[
P(R)=0,\quad \varepsilon_s=4B_{\rm eff},
\]
\[
M\simeq\frac{16\pi B_{\rm eff}}{3c^2}R^3.
\]
Il libro riporta $R^{-3}$ per refuso.
\medskip
\textbf{Dipendenza dai parametri dell'EOS}
\[
M_{\max},R_{\max}\propto B_{\rm eff}^{-1/2}.
\]
\end{column}
\begin{column}{0.49\textwidth}
\textbf{Crosta e limiti gravitazionali}
\begin{itemize}
\item La crosta amplia il raggio, soprattutto a piccola massa.
\item Orizzonte: $\mathcal C=1/2$.
\item Buchdahl: $\mathcal C\leq4/9$ sotto le sue ipotesi.
\item Causalita': $\mathcal C\lesssim0.354$ nel quadro del libro.
\end{itemize}
\medskip
La sola curva $M(R)$ non identifica necessariamente la composizione microscopica.
\end{column}
\end{columns}
\end{frame}

\section{Stelle ibride: EOS nucleonica e QCD}
\subsection{Regioni di densita' e vincoli teorici}

% 01
\begin{frame}{Capitolo 9: il problema delle stelle ibride}
\small
\begin{columns}[T,onlytextwidth]
\begin{column}{0.49\textwidth}
\textbf{Basse densita'}
\begin{itemize}
\item Descrizione con nucleoni, nuclei ed elettroni.
\item EOS nucleonica vincolata dall'interazione nucleare.
\item Vicino alla densita' di saturazione $n_0$.
\end{itemize}
\end{column}
\begin{column}{0.49\textwidth}
\textbf{Altissime densita'}
\begin{itemize}
\item Gradi di liberta': quark e gluoni.
\item QCD perturbativa, affidabile solo a potenziali chimici elevati.
\item EOS di materia di quark: $P_{\rm pQCD}(\mu_q)$.
\end{itemize}
\end{column}
\end{columns}
\medskip
La regione intermedia non e' fissata univocamente da questi due limiti.
\medskip
Scopo del capitolo: costruire EOS che li colleghino e studiarne le soluzioni TOV.
\end{frame}

% 02
\begin{frame}{Transizione adroni--quark: che cosa sappiamo?}
\small
\begin{columns}[T,onlytextwidth]
\begin{column}{0.49\textwidth}
Il passaggio fra gradi di liberta' adronici e quark puo' essere:
\begin{itemize}
\item un \emph{crossover} liscio;
\item una transizione continua con non analiticita' delle derivate;
\item una transizione del primo ordine con salto della densita' di energia.
\end{itemize}
\end{column}
\begin{column}{0.49\textwidth}
Le informazioni della QCD a $T$ elevata e $\mu_B$ piccolo non determinano la natura della transizione a
\[
T\simeq0,\qquad \mu_B\gg0.
\]
L'interpolazione fra le EOS non costituisce, da sola, prova di deconfinamento o di una transizione di fase.
\end{column}
\end{columns}
\end{frame}

% 03
\begin{frame}{Scala nucleare: densita' di saturazione}
\small
\begin{columns}[T,onlytextwidth]
\begin{column}{0.49\textwidth}
Per materia nucleare simmetrica a $T=0$:
\[
n_{\rm sat}\simeq 0.15\text{--}0.16\,\mathrm{fm}^{-3}.
\]
Massa media del nucleone:
\[
m_N\simeq 939\,\mathrm{MeV}.
\]
La scala di densita' di energia e'
\[
\varepsilon_0\simeq m_N n_{\rm sat}.
\]
\end{column}
\begin{column}{0.49\textwidth}
Usando $n_{\rm sat}=0.15\,\mathrm{fm}^{-3}$:
\[
\begin{aligned}
\varepsilon_0&\simeq 939\times0.15\\
&\simeq 141\,\mathrm{MeV\,fm}^{-3}.
\end{aligned}
\]
Il libro arrotonda a
\[
\varepsilon_0\simeq140\,\mathrm{MeV\,fm}^{-3}.
\]
Questa scala serve a esprimere le densita' dei capitoli 7--9.
\end{column}
\end{columns}
\end{frame}

% 04
\begin{frame}{Crosta e materia neutronica: gerarchia delle densita'}
\footnotesize
\begin{columns}[T,onlytextwidth]
\begin{column}{0.49\textwidth}
\textbf{Crosta esterna}
\[
\varepsilon\lesssim10^{-3}\varepsilon_0.
\]
Reticolo di nuclei ed elettroni degeneri; assenza di neutroni liberi.
\medskip
\textbf{Crosta interna}
\[
10^{-3}\varepsilon_0\lesssim\varepsilon
\lesssim\varepsilon_0/2.
\]
Reticolo di nuclei, elettroni e neutroni oltre il \emph{neutron drip}.
\end{column}
\begin{column}{0.49\textwidth}
\textbf{Materia nucleonica ricca di neutroni}
\[
\varepsilon\gtrsim\varepsilon_0/2.
\]
Si passa verso il fluido nucleonico del core.
\medskip
A densita' prossime e poco superiori a $\varepsilon_0$, il testo assume un'incertezza nucleonica che puo' gia' raggiungere circa il $20\%$.
\medskip
Le soglie indicate sono ordini di grandezza della descrizione utilizzata.
\end{column}
\end{columns}
\end{frame}

\subsection{Energia di simmetria: origine di $m_N-16+32$}

% 05
\begin{frame}{Composizione nucleonica e asimmetria isospinica}
\small
Definiamo
\[
n_B=n_n+n_p,\qquad Y_p=\frac{n_p}{n_B},
\]
\[
\delta\equiv\frac{n_n-n_p}{n_B}=1-2Y_p.
\]
\begin{columns}[T,onlytextwidth]
\begin{column}{0.49\textwidth}
\textbf{Materia simmetrica}
\[
n_n=n_p,\quad Y_p=\frac12,\quad\delta=0.
\]
Per simmetria di isospin, trascurando piccole rotture, $e(n_B,\delta)$ e' pari in $\delta$.
\end{column}
\begin{column}{0.49\textwidth}
\textbf{Materia neutronica pura}
\[
n_p=0,\quad Y_p=0,\quad\delta=1.
\]
La differenza energetica e' dominata dall'energia di simmetria, non dalla sola massa a riposo.
\end{column}
\end{columns}
\end{frame}

% 06
\begin{frame}{Derivazione dell'approssimazione parabolica}
\small
Sia $e(n_B,\delta)=\varepsilon(n_B,\delta)/n_B$ l'energia \emph{totale} per barione, inclusa la massa a riposo.
\medskip
Espandiamo a densita' fissata attorno a $\delta=0$:
\[
\begin{aligned}
e(n_B,\delta)={}&e(n_B,0)
+\left.\frac{\partial e}{\partial\delta}\right|_0\delta\\
&+\frac12\left.\frac{\partial^2e}{\partial\delta^2}\right|_0\delta^2
+O(\delta^4).
\end{aligned}
\]
La simmetria approssimata $n_n\leftrightarrow n_p$ implica $e(n_B,\delta)=e(n_B,-\delta)$: il termine lineare e' nullo.
\[
{e(n_B,\delta)\simeq e(n_B,0)+S(n_B)\delta^2},
\qquad
S(n_B)=\frac12\left.\frac{\partial^2e}{\partial\delta^2}\right|_0.
\]
\end{frame}

% 07
\begin{frame}{Da $m_N-16$ alla materia neutronica pura}
\small
\begin{columns}[T,onlytextwidth]
\begin{column}{0.49\textwidth}
\textbf{Input fenomenologici a saturazione}
\[
e(n_{\rm sat},0)\simeq m_N-16\,\mathrm{MeV},
\]
\[
S(n_{\rm sat})\equiv J\simeq32\,\mathrm{MeV}.
\]
$16$ MeV: legame per nucleone della materia simmetrica infinita.
\medskip
$32$ MeV: coefficiente di simmetria \emph{della materia nucleare} a saturazione.
\end{column}
\begin{column}{0.49\textwidth}
\textbf{Materia di neutroni puri}
\[
\delta=1\quad\Longrightarrow\quad \delta^2=1.
\]
\[
\begin{aligned}
e_{\rm PNM}(n_{\rm sat})
&\simeq e(n_{\rm sat},0)+J\\
&\simeq m_N-16+32\\
&=m_N+16\\
&\simeq955\,\mathrm{MeV}.
\end{aligned}
\]
\end{column}
\end{columns}
\medskip
\textit{Precisazione:} sono valori fenomenologici approssimati, non una previsione diretta della pQCD.
\end{frame}

% 08
\begin{frame}{Come nasce il termine di simmetria nelle masse nucleari?}
\small
Per un nucleo con $A$ nucleoni e $Z$ protoni, la formula semiempirica include
\[
E_b\supset a_vA-a_{\rm sym}^{(A)}\frac{(A-2Z)^2}{A}.
\]
Poiche' $\delta=(A-2Z)/A$, dividendo per $A$ e trascurando superficie, Coulomb e pairing:
\[
\frac{E_b}{A}\simeq a_v-a_{\rm sym}^{\rm bulk}\delta^2.
\]
L'energia totale per barione e'
\[
\begin{aligned}
e&\simeq m_N-\frac{E_b}{A}\\
&\simeq m_N-a_v+a_{\rm sym}^{\rm bulk}\delta^2.
\end{aligned}
\]
Nel limite della materia nucleare: $a_v\simeq16$ MeV e $a_{\rm sym}^{\rm bulk}=J\simeq32$ MeV.
\medskip
\textit{Nota:} il coefficiente di simmetria di un nucleo finito non coincide esattamente con $J$ perche' include effetti di superficie.
\end{frame}

% 09
\begin{frame}{Energia per barione e potenziale chimico non coincidono sempre}
\small
Per materia fredda con un solo numero barionico conservato:
\[
\mu_B=\frac{d\varepsilon}{dn_B},\qquad P=n_B\mu_B-\varepsilon,
\]
quindi
\[
{\mu_B=\frac{\varepsilon}{n_B}+\frac{P}{n_B}=e+\frac{P}{n_B}.}
\]
\begin{columns}[T,onlytextwidth]
\begin{column}{0.49\textwidth}
Materia \emph{simmetrica} alla saturazione:
\[
P(n_{\rm sat},\delta=0)=0
\]
e percio' $\mu_B=e\simeq m_N-16$ MeV.
\end{column}
\begin{column}{0.49\textwidth}
Materia \emph{neutronica} alla stessa densita':
\[
P_{\rm PNM}(n_{\rm sat})>0
\]
in generale, dunque $\mu_n>e_{\rm PNM}$.
\end{column}
\end{columns}
\medskip
Il valore $955$ MeV del testo fornisce la \emph{scala} $E/A$, non necessariamente il potenziale chimico esatto.
\end{frame}

% 10
\begin{frame}{Correzione esplicita: il ruolo della pendenza $L$}
\small
Scriviamo
\[
e(n_B,\delta)\simeq e_{\rm sat}(n_B)+S(n_B)\delta^2.
\]
Per neutroni puri a $n_{\rm sat}$, $\delta=1$, mentre $e_{\rm sat}'(n_{\rm sat})=0$.
\[
\begin{aligned}
P_{\rm PNM}(n_{\rm sat})
&=n_{\rm sat}^2\left.\frac{d e_{\rm PNM}}{dn_B}\right|_{n_{\rm sat}}\\
&\simeq n_{\rm sat}^2 S'(n_{\rm sat}).
\end{aligned}
\]
Definendo $L=3n_{\rm sat}S'(n_{\rm sat})$:
\[
P_{\rm PNM}(n_{\rm sat})\simeq\frac{n_{\rm sat}L}{3},
\qquad
{\mu_n\simeq m_N+16+\frac L3.}
\]
La stima del libro $\mu_B\sim955$ MeV trascura questo termine di pressione.
\end{frame}

\subsection{Vincolo perturbativo ad altissime densita'}

% 11
\begin{frame}{Limite pQCD: quanto e' distante dalla saturazione?}
\small
\begin{columns}[T,onlytextwidth]
\begin{column}{0.49\textwidth}
Il testo prende come scala indicativa:
\[
\mu_q\simeq900\,\mathrm{MeV},
\]
con
\[
\mu_B=3\mu_q\simeq2.7\,\mathrm{GeV}.
\]
All'estremo nucleonico:
\[
\mu_B\sim0.95\,\mathrm{GeV}.
\]
Differenza nei potenziali chimici: meno di un fattore $3$.
\end{column}
\begin{column}{0.49\textwidth}
Ma la densita' di energia cresce approssimativamente come
\[
\varepsilon_{\rm free}\propto\mu_q^4.
\]
Ne consegue un divario assai maggiore fra densita' nucleari e densita' a cui la pQCD e' controllata.
\medskip
La materia nel core delle stelle osservate ricade in gran parte nella regione intermedia.
\end{column}
\end{columns}
\end{frame}

% 12
\begin{frame}{Densita' del gas libero a $\mu_q=900$ MeV}
\small
Tre sapori massless, $T=0$:
\[
P_{\rm SB}=\frac{3\mu_q^4}{4\pi^2},
\qquad \varepsilon_{\rm SB}=3P_{\rm SB}=\frac{9\mu_q^4}{4\pi^2}.
\]
Ripristinando $\hbar c$:
\[
\varepsilon_{\rm SB}[\mathrm{MeV\,fm^{-3}}]
=\frac{9(900\,\mathrm{MeV})^4}{4\pi^2(\hbar c)^3},
\quad \hbar c\simeq197.327\,\mathrm{MeV\,fm}.
\]
Numericalmente:
\[
\varepsilon_{\rm SB}\simeq1.95\times10^4\,\mathrm{MeV\,fm^{-3}}
\simeq140\,\varepsilon_0.
\]
Con le correzioni interagenti, il libro stima $\varepsilon_{\rm pQCD}\sim70\text{--}100\,\varepsilon_0$.
\end{frame}

\subsection{Figura 9.1: il divario fra le EOS}

% 13
\begin{frame}{Figura 9.1: EOS nucleonica, regione incerta e pQCD}
\small
\begin{columns}[T,onlytextwidth]
\begin{column}{0.63\textwidth}
    \begin{figure}[!ht]\includegraphics[trim={0cm 0cm 0 0},clip,keepaspectratio,width=0.95\textwidth]{Schaffner/EOS-crust-pQCD}\label{fig:Schaffner_EOS-crust-pQCD}\end{figure}
\end{column}
\begin{column}{0.35\textwidth}
\textbf{Asse verticale}
\[
P\ [\mathrm{MeV\,fm^{-3}}].
\]
\textbf{Asse orizzontale}
\[
\widetilde\mu_q=\mu_q-\frac{\mu_{\rm Fe}}3.
\]
Entrambi gli assi: scala logaritmica.
\medskip
A sinistra: croste e materia nucleonica. A destra: pQCD.
\end{column}
\end{columns}
\end{frame}

% 14
\begin{frame}{Perche' traslare il potenziale chimico dei quark?}
\small
Il potenziale chimico di riferimento del ferro e' assunto come
\[
\mu_{\rm Fe}\simeq930\,\mathrm{MeV},\qquad \frac{\mu_{\rm Fe}}3\simeq310\,\mathrm{MeV}.
\]
Definendo $\widetilde\mu_q=\mu_q-\mu_{\rm Fe}/3$,
\[
\mu_q=310\,\mathrm{MeV}+\widetilde\mu_q,
\qquad \mu_B=930\,\mathrm{MeV}+3\widetilde\mu_q.
\]
Esempio:
\[
\widetilde\mu_q=10\,\mathrm{MeV}
\quad\Longrightarrow\quad\mu_B\simeq960\,\mathrm{MeV}.
\]
La traslazione rende leggibili su scala logaritmica piccoli incrementi di potenziale chimico rispetto alla materia nucleare stabile.
\end{frame}

% 15
\begin{frame}{Lettura dei confini nella Figura 9.1}
\small
Le linee verticali delimitano indicativamente:
\[
\begin{array}{ll}
\widetilde\mu_q\simeq3\,\mathrm{MeV} & \text{fine crosta esterna},\\[2mm]
\widetilde\mu_q\simeq7.4\,\mathrm{MeV} & \text{fine crosta interna},\\[2mm]
\widetilde\mu_q\simeq16\,\mathrm{MeV} & \text{estremo nucleonico}.
\end{array}
\]
Il margine inferiore pQCD si trova molto piu' a destra ($\widetilde\mu_q\sim560$ MeV nel grafico).
\medskip
La zona centrale, con il punto interrogativo, richiede una EOS interpolata; non equivale a una regione sperimentalmente accertata di quark liberi.
\medskip
La fascia vicina all'asse orizzontale segnala che le densita' centrali delle stelle di massa massima ricadono proprio nell'intervallo incerto.
\end{frame}

\section{Interpolazione con politropiche a tratti}
\subsection{Termodinamica a temperatura zero}

% 16
\begin{frame}{Politropiche a tratti: definizione}
\small
Per ogni intervallo di densita' si assume
\[
{P(n)=Kn^{\Gamma}},
\]
con $n=n_B$, $K$ costante e $\Gamma$ indice politropico.
\[
\Gamma=\frac{d\ln P}{d\ln n}.
\]
In una EOS a tratti:
\[
P(n)=\begin{cases}
K_1n^{\Gamma_1},&n_0\leq n<n_1,\\
K_2n^{\Gamma_2},&n_1\leq n<n_2,\\
\cdots&\cdots
\end{cases}
\]
Le costanti di ciascun segmento devono essere determinate con condizioni di raccordo termodinamico.
\end{frame}

% 17
\begin{frame}{Identita' termodinamiche essenziali}
\small
A temperatura zero, materia neutra in equilibrio con un'unica carica conservata indipendente:
\[
d\varepsilon=\mu\,dn,\qquad
P=\mu n-\varepsilon,\qquad dP=n\,d\mu.
\]
Da esse seguono
\[
{\mu=\frac{d\varepsilon}{dn}},
\qquad
{n=\frac{dP}{d\mu}}.
\]
\begin{columns}[T,onlytextwidth]
\begin{column}{0.49\textwidth}
La forma $P(n)$ \emph{non basta} a definire autonomamente l'energia: occorre integrare l'identita' termodinamica.
\end{column}
\begin{column}{0.49\textwidth}
La continuita' di $P(\mu)$ e della sua derivata garantisce un raccordo senza salto della densita' barionica.
\end{column}
\end{columns}
\end{frame}

% 18
\begin{frame}{Derivazione dell'identita' (9.2)}
\small
Partiamo da
\[
P=n\mu-\varepsilon,
\qquad\mu=\frac{d\varepsilon}{dn}.
\]
Sostituiamo il potenziale chimico:
\[
P=n\frac{d\varepsilon}{dn}-\varepsilon.
\]
Dividendo per $n^2$:
\[
\frac P{n^2}=\frac1n\frac{d\varepsilon}{dn}-\frac{\varepsilon}{n^2}.
\]
Il membro destro e' la derivata di un rapporto:
\[
{P=n^2\frac{d}{dn}\left(\frac{\varepsilon}{n}\right).}
\]
\end{frame}

% 19
\begin{frame}{Dalla politropica alla densita' di energia (9.3)}
\small
Assumiamo in un intervallo
\[
P(n)=Kn^\Gamma.
\]
Dall'identita' precedente:
\[
Kn^\Gamma=n^2\frac{d}{dn}\left(\frac{\varepsilon}{n}\right).
\]
Dividiamo per $n^2$:
\[
{\frac{d}{dn}\left(\frac{\varepsilon}{n}\right)=Kn^{\Gamma-2}.}
\]
Per $\Gamma\ne1$, integriamo:
\[
\frac{\varepsilon}{n}=K\int n^{\Gamma-2}dn+c
=\frac{K}{\Gamma-1}n^{\Gamma-1}+c.
\]
\end{frame}

% 20
\begin{frame}{Densita' di energia politropica (9.4)}
\small
Moltiplicando il risultato integrato per $n$:
\[
{\varepsilon(n)=\frac{K}{\Gamma-1}n^\Gamma+cn.}
\]
Poiche' $P=Kn^\Gamma$,
\[
{\varepsilon(n)=\frac{P(n)}{\Gamma-1}+cn.}
\]
\begin{columns}[T,onlytextwidth]
\begin{column}{0.49\textwidth}
Il primo termine e' determinato dalla legge politropica $P(n)$.
\end{column}
\begin{column}{0.49\textwidth}
Il termine $cn$ integra l'informazione sull'energia per barione al raccordo; in generale $c\ne m_N$.
\end{column}
\end{columns}
\end{frame}

% 21
\begin{frame}{Potenziale chimico della politropica (9.5)}
\small
Per definizione:
\[
\mu(n)=\frac{d\varepsilon(n)}{dn}.
\]
Derivando termine a termine:
\[
\begin{aligned}
\mu(n)&=\frac{d}{dn}\left[\frac{K}{\Gamma-1}n^\Gamma+cn\right]\\
&=\frac{K\Gamma}{\Gamma-1}n^{\Gamma-1}+c.
\end{aligned}
\]
Quindi
\[
{\mu(n)=c+\frac{K\Gamma}{\Gamma-1}n^{\Gamma-1}.}
\]
Verifica: $n\mu-\varepsilon=Kn^\Gamma=P$.
\end{frame}

\subsection{Raccordo al punto $(n_0,P_0,\mu_0)$}

% 22
\begin{frame}{Primo vincolo di raccordo: pressione}
\small
Sia $(n_0,P_0,\mu_0)$ il punto terminale dell'EOS nota.
\medskip
Per il segmento politropico imponiamo
\[
P(n_0)=P_0.
\]
Con $P(n)=Kn^\Gamma$:
\[
P_0=Kn_0^\Gamma
\quad\Longrightarrow\quad
{K=\frac{P_0}{n_0^\Gamma}.}
\]
\begin{columns}[T,onlytextwidth]
\begin{column}{0.49\textwidth}
$\Gamma$ determina la forma del segmento.
\end{column}
\begin{column}{0.49\textwidth}
$K$ non e' libero una volta scelta la pressione al raccordo.
\end{column}
\end{columns}
\medskip
Il simbolo $n_0$ qui identifica il \emph{punto di raccordo}, non necessariamente la densita' di saturazione nucleare.
\end{frame}

% 23
\begin{frame}{Secondo vincolo: continuita' del potenziale (9.6)}
\small
Partiamo da
\[
\mu(n)=c+\frac{K\Gamma}{\Gamma-1}n^{\Gamma-1}.
\]
Imponiamo $\mu(n_0)=\mu_0$:
\[
\mu_0=c+\frac{K\Gamma}{\Gamma-1}n_0^{\Gamma-1}.
\]
Si ricava
\[
{c=\mu_0-\frac{K\Gamma}{\Gamma-1}n_0^{\Gamma-1}.}
\]
Inserendo $K=P_0/n_0^\Gamma$:
\[
{c=\mu_0-\frac{\Gamma}{\Gamma-1}\frac{P_0}{n_0}.}
\]
Le due condizioni fissano $K$ e $c$; rimane da scegliere $\Gamma$.
\end{frame}

% 24
\begin{frame}{Controllo del raccordo: anche $\varepsilon$ e' continua}
\small
Alla densita' $n_0$,
\[
\varepsilon(n_0)=\frac{P_0}{\Gamma-1}+cn_0.
\]
Usiamo
\[
cn_0=\mu_0n_0-\frac{\Gamma}{\Gamma-1}P_0.
\]
Segue
\[
\begin{aligned}
\varepsilon(n_0)
&=\frac{P_0}{\Gamma-1}+\mu_0n_0
-\frac{\Gamma P_0}{\Gamma-1}\\
&=\mu_0n_0-P_0.
\end{aligned}
\]
\[
{\varepsilon_0+P_0=\mu_0n_0.}
\]
La densita' di energia coincide con quella dell'EOS nota se il punto soddisfa la stessa identita' termodinamica.
\end{frame}

\subsection{Da $P(n)$ a $P(\mu)$ e $\varepsilon(P)$}

% 25
\begin{frame}{Invertire il potenziale chimico: primo passaggio}
\small
Partiamo dalla (9.5):
\[
\mu-c=\frac{K\Gamma}{\Gamma-1}n^{\Gamma-1}.
\]
Isolando la potenza di $n$:
\[
n^{\Gamma-1}=\frac{\Gamma-1}{K\Gamma}(\mu-c).
\]
Elevando a $1/(\Gamma-1)$:
\[
{n(\mu)=\left[\frac{\Gamma-1}{K\Gamma}(\mu-c)\right]^{1/(\Gamma-1)}.}
\]
L'argomento della potenza deve rimanere positivo nel tratto fisico scelto.
\end{frame}

% 26
\begin{frame}{Densita' in funzione di $\mu$: raccordo (9.7)}
\small
Sostituiamo
\[
c=\mu_0-\frac{K\Gamma}{\Gamma-1}n_0^{\Gamma-1}
\]
nel termine $\mu-c$:
\[
\mu-c=(\mu-\mu_0)+\frac{K\Gamma}{\Gamma-1}n_0^{\Gamma-1}.
\]
Pertanto
\[
{
 n(\mu)=\left[n_0^{\Gamma-1}
 +\frac{\Gamma-1}{K\Gamma}(\mu-\mu_0)
 \right]^{1/(\Gamma-1)}.}
\]
Controllo immediato: $n(\mu_0)=n_0$.
\end{frame}

% 27
\begin{frame}{Pressione in funzione di $\mu$: Eq. (9.8)}
\small
Partendo da $P=Kn^\Gamma$ e dalla (9.7):
\[
{
P(\mu)=K\left[
 n_0^{\Gamma-1}+\frac{\Gamma-1}{K\Gamma}(\mu-\mu_0)
\right]^{\Gamma/(\Gamma-1)}.}
\]
Al raccordo $\mu=\mu_0$:
\[
P(\mu_0)=Kn_0^\Gamma=P_0.
\]
La funzione pressione e' quindi continua. Rimane da verificarne la derivata.
\end{frame}

% 28
\begin{frame}{Verifica: $dP/d\mu=n$ in ogni segmento}
\small
Definiamo
\[
X(\mu)=n_0^{\Gamma-1}+\frac{\Gamma-1}{K\Gamma}(\mu-\mu_0).
\]
Allora $P(\mu)=KX^{\Gamma/(\Gamma-1)}$ e
\[
\begin{aligned}
\frac{dP}{d\mu}
&=K\frac{\Gamma}{\Gamma-1}X^{1/(\Gamma-1)}
\frac{\Gamma-1}{K\Gamma}\\
&=X^{1/(\Gamma-1)}=n(\mu).
\end{aligned}
\]
Per $\mu=\mu_0$:
\[
{\left.\frac{dP}{d\mu}\right|_{\mu_0}=n_0.}
\]
Questa e' la condizione sulla pendenza usata per il raccordo liscio delle politropiche.
\end{frame}

% 29
\begin{frame}{Dalla pressione alla densita' per le TOV}
\small
Per integrare le TOV occorre $\varepsilon(P)$.
\medskip
Da $P=Kn^\Gamma$:
\[
n=\left(\frac PK\right)^{1/\Gamma}.
\]
Poiche' $P_0=Kn_0^\Gamma$, equivalentemente
\[
{n(P)=n_0\left(\frac P{P_0}\right)^{1/\Gamma}.}
\]
Usiamo la densita' di energia gia' integrata:
\[
\varepsilon(n)=\frac{P(n)}{\Gamma-1}+cn.
\]
Sostituiremo $n(P)$ e poi il valore di $c$ fissato al raccordo.
\end{frame}

% 30
\begin{frame}{EOS per le TOV: derivazione di (9.9)}
\small
Sostituendo $n(P)$:
\[
\varepsilon(P)=\frac{P}{\Gamma-1}+cn_0\left(\frac{P}{P_0}\right)^{1/\Gamma}.
\]
Dal vincolo sul potenziale:
\[
cn_0=\mu_0n_0-\frac{\Gamma}{\Gamma-1}P_0.
\]
Si ottiene
\[
{
\begin{aligned}
\varepsilon(P)={}&\frac{P}{\Gamma-1}\\
&+\left(\mu_0n_0-\frac{\Gamma}{\Gamma-1}P_0\right)
\left(\frac{P}{P_0}\right)^{1/\Gamma}.
\end{aligned}}
\]
Il termine in $P^{1/\Gamma}$ conserva l'informazione termodinamica del raccordo.
\end{frame}

% 31
\begin{frame}{Il caso speciale $\Gamma=1$}
\small
Le formule precedenti con $1/(\Gamma-1)$ richiedono $\Gamma\ne1$.
\medskip
Se $P=Kn$, l'equazione (9.2) diventa
\[
\frac{d}{dn}\left(\frac{\varepsilon}{n}\right)=\frac Kn.
\]
Integrando, con densita' di riferimento $n_{\rm ref}$,
\[
{\varepsilon(n)=Kn\ln\frac{n}{n_{\rm ref}}+cn.}
\]
Derivando:
\[
{\mu(n)=K\left[\ln\frac{n}{n_{\rm ref}}+1\right]+c.}
\]
Il segmento con $\Gamma=1$ si tratta dunque separatamente mediante un logaritmo.
\end{frame}

% 32
\begin{frame}{Algoritmo di raccordo per piu' politropiche}
\small
\begin{enumerate}
\item Partire dal punto nucleonico $(n_0,P_0,\mu_0)$.
\item Scegliere $\Gamma_1$ e determinare $K_1,c_1$ imponendo $P_0,\mu_0$.
\item Proseguire fino a un nuovo punto $(n_1,P_1,\mu_1)$.
\item Scegliere $\Gamma_2$ e imporre le stesse condizioni per il secondo segmento.
\item Raccordare all'EOS pQCD, controllando vincoli fisici lungo tutta la curva.
\end{enumerate}
\medskip
Al raccordo liscio:
\[
P_-(\mu_i)=P_+(\mu_i),\qquad
\left.\frac{dP_-}{d\mu}\right|_{\mu_i}
=\left.\frac{dP_+}{d\mu}\right|_{\mu_i}=n_i.
\]
\end{frame}

\subsection{Causalita', rigidita' e transizioni di fase}

% 33
\begin{frame}{Derivazione della velocita' del suono}
\small
A $T=0$, in un fluido barotropico:
\[
c_s^2=\frac{dP}{d\varepsilon}
=\frac{dP/dn}{d\varepsilon/dn}
\qquad(c=1).
\]
Per $P=Kn^\Gamma$:
\[
\frac{dP}{dn}=\Gamma Kn^{\Gamma-1}=\frac{\Gamma P}{n}.
\]
Dalla termodinamica:
\[
\frac{d\varepsilon}{dn}=\mu=\frac{\varepsilon+P}{n}.
\]
Quindi
\[
{c_s^2=\frac{\Gamma P}{\varepsilon+P}.}
\]
Causalita' e stabilita' locale richiedono $0\leq c_s^2\leq1$.
\end{frame}

% 34
\begin{frame}{Raccordo liscio e transizione del primo ordine}
\small
\begin{columns}[T,onlytextwidth]
\begin{column}{0.49\textwidth}
\textbf{Interpolazione liscia}
\[
P_H(\mu_i)=P_Q(\mu_i),
\]
\[
n_H(\mu_i)=n_Q(\mu_i).
\]
Per $\varepsilon=\mu n-P$ risulta anche $\varepsilon_H=\varepsilon_Q$ al punto.
\medskip
Nessun salto di densita' introdotto dal raccordo.
\end{column}
\begin{column}{0.49\textwidth}
\textbf{Primo ordine (Maxwell)}
\[
P_H(\mu_t)=P_Q(\mu_t),
\]
\[
n_Q\ne n_H.
\]
A pressione e potenziale comuni,
\[
{\Delta\varepsilon=\mu_t\Delta n_B.}
\]
La pressione continua, ma la derivata $dP/d\mu$ salta.
\end{column}
\end{columns}
\end{frame}

% 35
\begin{frame}{Quali politropiche sono fisicamente ammissibili?}
\small
\begin{columns}[T,onlytextwidth]
\begin{column}{0.49\textwidth}
\textbf{Vincoli locali}
\begin{itemize}
\item $n=dP/d\mu>0$;
\item stabilita' meccanica $dP/dn>0$;
\item causalita' $c_s^2=dP/d\varepsilon\leq1$.
\end{itemize}
\end{column}
\begin{column}{0.49\textwidth}
\textbf{Vincolo globale}
\[
M_{\max}\gtrsim2M_\odot
\]
per essere compatibili con l'esistenza di stelle massicce osservate.
\medskip
Una EOS troppo soffice non sostiene grandi masse; una troppo rigida puo' violare la causalita'.
\end{column}
\end{columns}
\medskip
La rigidita' stellare si giudica da $P(\varepsilon)$ e dalla velocita' del suono, non dal solo valore di $\Gamma$ isolatamente.
\end{frame}

\section{Massa--raggio delle stelle ibride}
\subsection{Figura 9.2 e vincoli astrofisici}

% 36
\begin{frame}{Figura 9.2: banda delle relazioni massa--raggio}
\small
\begin{columns}[T,onlytextwidth]
\begin{column}{0.64\textwidth}
    \begin{figure}[!ht]\includegraphics[trim={0cm 0cm 0 0},clip,keepaspectratio,width=0.95\textwidth]{Schaffner/MR-NS-hybrid}\label{fig:Schaffner_MR-NS-hybrid}\end{figure}
\end{column}
\begin{column}{0.34\textwidth}
Banda grigia: famiglie di EOS interpolate fra limite nucleonico e pQCD.
\medskip
Curve I, II e III: tre politropiche rappresentative.
\medskip
Linee: buchi neri, Buchdahl, limite causale e compattezze costanti.
\end{column}
\end{columns}
\end{frame}

% 37
\begin{frame}{Tre sequenze rappresentative nella Figura 9.2}
\footnotesize
\begin{columns}[T,onlytextwidth]
\begin{column}{0.33\textwidth}
\textbf{I: piu' compatta}
\medskip
Raggi relativamente piccoli.
\[
\varepsilon_{c,\max}\sim10\varepsilon_0.
\]
La configurazione di massa massima raggiunge densita' centrali elevate.
\end{column}
\begin{column}{0.33\textwidth}
\textbf{II: massa massima maggiore}
\[
M_{\max}\simeq2.4M_\odot,
\]
\[
\varepsilon_{c,\max}\sim5\varepsilon_0.
\]
Mostra come una EOS adeguatamente rigida possa sostenere masse elevate.
\end{column}
\begin{column}{0.33\textwidth}
\textbf{III: raggi maggiori}
\medskip
Sequenza piu' a destra nel diagramma.
\[
\varepsilon_{c,\max}\sim5\varepsilon_0.
\]
Raggi piu' ampi alle masse considerate.
\end{column}
\end{columns}
\end{frame}

% 38
\begin{frame}{Densita' centrali: cosa sondano le osservazioni?}
\small
Per la famiglia di EOS del testo:
\[
\begin{array}{c|c|c}
M & R & \varepsilon_c/\varepsilon_0\\ \hline
1.4M_\odot & 11\text{--}14.5\,\mathrm{km} & 1.6\text{--}4\\[1mm]
2.0M_\odot & 10\text{--}15\,\mathrm{km} & 2\text{--}8
\end{array}
\]
Confronto fra scale:
\[
\varepsilon_c\sim\text{pochi }\varepsilon_0,
\qquad \varepsilon_{\rm pQCD}\gtrsim70\,\varepsilon_0.
\]
Le osservazioni delle stelle compatte mettono alla prova la regione intermedia dell'EOS, non direttamente il regime di pQCD asintotica.
\medskip
{\scriptsize Intervalli del lavoro citato nel libro; non sono vincoli osservativi aggiornati.}
\end{frame}

% 39
\begin{frame}{Compattezza, causalita' e massa osservata}
\small
\begin{columns}[T,onlytextwidth]
\begin{column}{0.49\textwidth}
\textbf{Compattezza}
\[
\mathcal C=\frac{GM}{Rc^2}.
\]
Rette di riferimento nel grafico:
\[
\mathcal C=0.1,\ 0.2,\ 0.3.
\]
Limiti esaminati nel capitolo 8:
\[
\mathcal C_{\rm BH}=\tfrac12,\quad
\mathcal C_{\rm Buchdahl}=\tfrac49.
\]
\end{column}
\begin{column}{0.49\textwidth}
\textbf{Vincoli sull'EOS}
\[
c_s^2\leq1,\qquad M_{\max}\geq2M_\odot.
\]
La causalita' impedisce una EOS arbitrariamente rigida.
\medskip
La massa massima osservata impedisce una EOS arbitrariamente soffice.
\medskip
La loro combinazione restringe la banda ammissibile nel diagramma $M$--$R$.
\end{column}
\end{columns}
\end{frame}

\subsection{Riepilogo}

% 40
\begin{frame}{Risultato del \S9.1: interpolare senza inventare la microfisica}
\small
\begin{enumerate}
\item EOS nucleonica vincolata vicino a $\varepsilon_0$.
\item EOS pQCD controllata solo intorno a $\varepsilon\gtrsim70\text{--}100\varepsilon_0$.
\item Politropiche a tratti raccordate con continuita' di $P$, $\mu$ e $n$.
\item Verifica di causalita', stabilita' e massa massima.
\item Integrazione delle TOV per ottenere una banda di curve $M(R)$.
\end{enumerate}
\medskip
Il raccordo liscio e' una \emph{parametrizzazione} dell'EOS intermedia, non una dimostrazione che la transizione adroni--quark sia un crossover.
\end{frame}

% 41
\begin{frame}{Formule da ricordare: politropiche termodinamiche}
\footnotesize
\begin{columns}[T,onlytextwidth]
\begin{column}{0.49\textwidth}
\[
P(n)=Kn^\Gamma,
\]
\[
\varepsilon(n)=\frac{Kn^\Gamma}{\Gamma-1}+cn,
\]
\[
\mu(n)=\frac{K\Gamma}{\Gamma-1}n^{\Gamma-1}+c,
\]
\[
c=\mu_0-\frac{\Gamma P_0}{(\Gamma-1)n_0}.
\]
\end{column}
\begin{column}{0.49\textwidth}
\[
n=\frac{dP}{d\mu},
\]
\[
\varepsilon(P)=\frac{P}{\Gamma-1}
+\left(\mu_0n_0-\frac{\Gamma P_0}{\Gamma-1}\right)
\left(\frac{P}{P_0}\right)^{1/\Gamma},
\]
\[
c_s^2=\frac{\Gamma P}{\varepsilon+P}\leq1.
\]
\end{column}
\end{columns}
\medskip
\textbf{Energia di simmetria:} $e_{\rm PNM}(n_{\rm sat})\simeq m_N-16+32=m_N+16\simeq955$ MeV.
\end{frame}

